%% ============================================================================
%%  CAPITOLO 3 — Limiti di successioni
%%  Fonte principale per argomenti e dimostrazioni: Marcellini–Sbordone, cap. 3
%%  (§16 premessa, §17–22 definizioni, operazioni, confronto, §23 limiti
%%  notevoli, §24 successioni monotòne, §25 il numero e; appendice §26–28:
%%  infiniti, Bolzano–Weierstrass, successioni di Cauchy).
%%  Arricchimenti da Giusti, Analisi Matematica 1, cap. 2 (§7, §10–13 e
%%  notizie storiche): seconda dimostrazione per (1+1/n)^n, massimo e minimo
%%  limite, completamento dei razionali, Mengoli e Bolzano.
%% ============================================================================
\aperturacapitolo{Limiti di successioni}{cap:successioni}

\epigraph{Quando i valori successivamente attribuiti a una variabile si avvicinano
indefinitamente a un valore fisso, così da finire per differirne di tanto poco quanto si
vuole, quest'ultimo è chiamato il limite di tutti gli altri.}{Augustin-Louis Cauchy,
\emph{Cours d'analyse} (1821)}

\lettrine[lines=3, findent=2pt, nindent=0pt]{C}{on questo capitolo} entriamo nel cuore
dell'Analisi Matematica. I due capitoli precedenti hanno preparato il terreno: abbiamo
fissato gli assiomi dei numeri reali e abbiamo imparato a usare l'assioma di completezza
attraverso l'estremo superiore. Ora introduciamo il concetto che distingue l'Analisi
dall'Algebra: il \emph{limite}.

Un limite risponde a una domanda che l'aritmetica, da sola, non sa porre: che cosa
succede a una quantità quando un indice cresce «senza fine»? L'area dei poligoni
regolari inscritti in un cerchio, al crescere del numero dei lati; il capitale prodotto
da un interesse capitalizzato sempre più spesso; l'errore di un metodo iterativo dopo
sempre più iterazioni. In nessuno di questi casi esiste un «ultimo» termine da calcolare,
eppure la successione dei valori si comporta in un modo preciso, e quel comportamento è
proprio ciò che ci interessa. La definizione di limite che daremo, dovuta a Cauchy e
resa definitivamente rigorosa da Weierstrass, è il frutto di due millenni di riflessione,
da Archimede a Newton: è breve, ma ogni sua parola è necessaria, e conviene leggerla con
calma.

\medskip
\noindent\textbf{In questo capitolo.}
\begin{itemize}
  \item Il problema dell'area del cerchio e l'idea di limite (\S\ref{sec:premessa-limiti}).
  \item La definizione di limite, finito e infinito; unicità (\S\ref{sec:def-limite}).
  \item Successioni limitate; le operazioni con i limiti (\S\ref{sec:operazioni-limiti}).
  \item Limiti infiniti e forme indeterminate (\S\ref{sec:forme-indeterminate}).
  \item Permanenza del segno, teorema dei carabinieri e altri confronti
  (\S\ref{sec:confronto}).
  \item I limiti notevoli: $a^{n}$, $\sqrt[n]a$, $\sqrt[n]n$, $\frac{\sen a_{n}}{a_{n}}$; l'area del
  cerchio (\S\ref{sec:limiti-notevoli}).
  \item Il teorema sulle successioni monotòne e il numero $e$ (\S\ref{sec:monotone},
  \S\ref{sec:numero-e}).
  \item Appendice: infiniti di ordine crescente, successioni estratte, teorema di
  Bolzano--Weierstrass, criterio di Cauchy (\S\ref{sec:infiniti}, \S\ref{sec:bw-cauchy}).
  \item Approfondimento: massimo e minimo limite (\S\ref{subsec:massimo-minimo-limite}).
  \item Punti fondamentali del capitolo (\S\ref{sec:sintesi-cap03}).
\end{itemize}

%% ----------------------------------------------------------------------------
%%  §3.1  Premessa: l'area del cerchio e l'idea di limite
%% ----------------------------------------------------------------------------
\section{Premessa: l'area del cerchio}
\label{sec:premessa-limiti}

Uno dei problemi che affronteremo nel capitolo sugli integrali sarà il calcolo delle
aree di figure piane. Tale calcolo fa uso in modo essenziale del concetto di limite. Per
mostrarlo, consideriamo un metodo per il calcolo dell'area di un cerchio simile a quello
che Archimede usò nel III secolo a.C.

Si tratta di «approssimare» un cerchio di raggio $1$ (potremmo scegliere anche un
diverso valore del raggio) con figure che «differiscono di poco» dal cerchio e di cui sia
possibile calcolare l'area. Nel capitolo sugli integrali vedremo che è comodo
considerare, come figure approssimanti, unioni di rettangoli. Qui, per ogni $n\ge3$,
consideriamo il poligono regolare $P_{n}$ di $n$ lati inscritto nel cerchio di raggio
$1$: per $n=3$ un triangolo equilatero, per $n=4$ un quadrato, per $n=6$ un esagono, e
così via (\cref{fig:poligoni-inscritti}).

\begin{figure}[htb]
\centering
\begin{tikzpicture}[scale=1.0, every node/.style={font=\footnotesize}]
  \foreach \n/\xx/\col in {3/0/blu, 4/2.45/rame, 6/4.9/verde} {
    \begin{scope}[shift={(\xx,0)}]
      \draw[grigio!70] (0,0) circle (1);
      \draw[\col, thick, fill=\col!12] (90:1) \foreach \k in {1,...,\n} { -- ({90+360*\k/\n}:1)} -- cycle;
      \fill (0,0) circle (0.8pt);
      \node at (0,-1.3) {$P_{\n}$};
    \end{scope}}
  % dettaglio di un triangolo
  \begin{scope}[shift={(6.6,0)}, scale=1.75]
    \pgfmathsetmacro{\al}{30}
    \draw[grigio!70] ({-\al-10}:1) arc ({-\al-10}:{\al+10}:1);
    \draw[blu, thick] (0,0) -- (\al:1) -- (-\al:1) -- cycle;
    \draw[blu, dashed] (0,0) -- ({cos(\al)},0);
    \draw[blu] ({cos(\al)},0) -- (\al:1);
    \draw ({cos(\al)-0.06},0) -- ({cos(\al)-0.06},0.06) -- ({cos(\al)},0.06);
    \draw[rame] (0.3,0) arc (0:\al:0.3); \node[rame] at (15:0.45) {$\alpha$};
    \node[below] at ({cos(\al)/2},-0.02) {$\cos\alpha$};
    \node[right] at ({cos(\al)},{sin(\al)/2}) {$\sen\alpha$};
    \node[above left] at ({cos(\al)/2},{sin(\al)/2}) {$1$};
    \node at (0.55,-0.85) {$\alpha=\pi/n$};
  \end{scope}
\end{tikzpicture}
\caption{I poligoni regolari $P_{3}$, $P_{4}$, $P_{6}$ inscritti nel cerchio di raggio
$1$. A destra, uno dei $2n$ triangoli rettangoli in cui si scompone $P_{n}$: l'angolo al
centro vale $\alpha=\pi/n$, i cateti valgono $\sen\alpha$ e $\cos\alpha$.}
\label{fig:poligoni-inscritti}
\end{figure}

Agli $n$ lati del poligono $P_{n}$ corrispondono $n$ angoli al centro uguali fra loro,
che sono l'$n$-esima parte dell'angolo giro e che quindi misurano $2\pi/n$ radianti.
Tracciando le altezze dal centro, $P_{n}$ si scompone in $2n$ triangoli rettangoli
uguali, con l'angolo al centro $\alpha=\pi/n$ e ipotenusa $1$; i cateti misurano quindi
$\sen\alpha$ e $\cos\alpha$, e l'area di ciascun triangolo vale
$\frac12\sen\alpha\cos\alpha$. Indicando con $a_{n}$ l'area di $P_{n}$ otteniamo
\begin{equation}\label{eq:area-poligono-1}
  a_{n}=\operatorname{area}(P_{n})=2n\cdot\frac12\sen\alpha\cos\alpha,
\end{equation}
e, usando la formula di duplicazione $\sen2\alpha=2\sen\alpha\cos\alpha$
\eqref{eq:duplicazione} e ricordando che $\alpha=\pi/n$,
\begin{equation}\label{eq:area-poligono}
  a_{n}=\frac n2\,\sen\frac{2\pi}{n}.
\end{equation}
Riportiamo $a_{n}$ per alcuni valori di $n$ nella \cref{tab:aree-poligoni}; il lettore
controlli i primi valori, fino a $n=12$, usando la \cref{tab:valori-notevoli} (per $n=3$
si ricordi che $\sen\frac{2\pi}3=\sen\bigl(\pi-\frac\pi3\bigr)=\sen\frac\pi3$).

\begin{table}[htb]
\centering
\renewcommand{\arraystretch}{1.5}
\footnotesize\setlength{\tabcolsep}{4pt}
\begin{tabular}{l*{10}{c}}
\toprule
$n$ & $3$ & $4$ & $6$ & $8$ & $12$ & $20$ & $50$ & $100$ & $200$ & $1000$\\
\midrule
$a_{n}$ & $\tfrac{3\sqrt3}{4}$ & $2$ & $\tfrac{3\sqrt3}{2}$ & $2\sqrt2$ & $3$ & $3{,}0902$ &
$3{,}1333$ & $3{,}1395$ & $3{,}1411$ & $3{,}14157$\\
\bottomrule
\end{tabular}
\caption{Le aree $a_{n}$ dei poligoni regolari di $n$ lati inscritti nel cerchio di
raggio $1$, calcolate con la \eqref{eq:area-poligono}.}
\label{tab:aree-poligoni}
\end{table}

Come ci aspettiamo, per $n$ grande i valori di $a_{n}$ si avvicinano a
$\pi=3{,}14159265\ldots$ Alla fine del \cref{sec:limiti-notevoli} dimostreremo, tramite
la \eqref{eq:area-poligono}, che l'area del cerchio di raggio $1$ vale $\pi$. La
\eqref{eq:area-poligono} non è però utile per il calcolo \emph{numerico} di $\pi$, perché
$\pi$ entra nella definizione stessa di $a_{n}$ (Archimede, come vedremo nella nota
storica, aggirò il problema in modo geniale).

Qui desideriamo sottolineare lo schema del procedimento: abbiamo introdotto una
\emph{successione} di numeri reali
\begin{equation}\label{eq:succ-aree}
  a_{3},\ a_{4},\ a_{5},\ a_{6},\ \dots,\ a_{n},\ \dots
\end{equation}
Di tale successione ci interessa il comportamento per $n$ grande. Infatti il numero
$a_{n}$ rappresenta l'area del poligono regolare di $n$ lati inscritto nel cerchio, ed
è tanto più vicino all'obiettivo --- l'area del cerchio --- quanto più $n$ è grande. Se
avesse un senso, diremmo che ci interessa «l'ultimo» termine della successione; però ciò
non ha senso, perché non esiste l'ultimo termine della successione.

Ciò che realmente ci interessa è il \emph{limite} della successione $a_{n}$, cioè un
numero $a\in\R$ che sia «vicino» ai termini della successione che hanno l'indice $n$
«grande». Ciò si esprime più precisamente così: $a$, limite della successione $a_{n}$,
è un numero reale tale che, comunque si scelga un intervallo di numeri intorno ad $a$,
diciamo $(a-\epsilon,a+\epsilon)$ con $\epsilon>0$, esiste un indice $\nu$ tale che, per
$n>\nu$, $a_{n}$ sta in questo intervallo, cioè $a-\epsilon<a_{n}<a+\epsilon$.

\begin{figure}[htb]
\centering
\begin{tikzpicture}[x=12cm, every node/.style={font=\footnotesize}]
  % a_n = 1 - 0.9 (-1)^n / n ;  |a_n - 1| = 0.9/n < eps  sse  n > 0.9/eps
  \foreach \row/\eps/\nuu/\lab in {0/0.30/3/{0{,}30}, 1/0.12/7/{0{,}12}, 2/0.04/22/{0{,}04}} {
    \begin{scope}[yshift={-\row*1.25cm}]
      \fill[verde!18] ({1-\eps},-0.16) rectangle ({1+\eps},0.16);
      \draw[-Latex, grigio] (0.70,0) -- (1.40,0);
      \foreach \n in {3,...,60} {
        \pgfmathsetmacro{\x}{1-0.9/\n*cos(180*\n)}
        \pgfmathparse{\n>\nuu ? 1 : 0}
        \ifnum\pgfmathresult=1 \fill[verde!70!black] (\x,0) circle (1.1pt);
        \else \fill[rame] (\x,0) circle (1.4pt); \fi
      }
      \draw[verde!60!black, thick] ({1-\eps},-0.16) -- ({1-\eps},0.16);
      \draw[verde!60!black, thick] ({1+\eps},-0.16) -- ({1+\eps},0.16);
      \node[verde!60!black, above] at ({1-\eps},0.16) {$a-\epsilon$};
      \node[verde!60!black, above] at ({1+\eps},0.16) {$a+\epsilon$};
      \draw[blu, thick] (1,-0.1) -- (1,0.1); \node[blu, below=2pt] at (1,-0.1) {$a$};
      \node[anchor=west] at (1.41,0) {$\epsilon=\lab$, \ $\nu=\nuu$};
    \end{scope}}
\end{tikzpicture}
\caption{L'idea di limite, per la successione $a_{n}=1-\frac{9}{10}\frac{(-1)^{n}}{n}$ e $a=1$
(qui $\abs{a_{n}-a}<\epsilon$ se e solo se $n>\frac{9}{10\epsilon}$). Per tre valori sempre
più piccoli di $\epsilon$, i termini
$a_{n}$ con $n>\nu$ (in verde) cadono tutti nell'intervallo $(a-\epsilon,a+\epsilon)$;
quelli con $n\le\nu$ (in rame) possono cadere fuori. Diminuendo $\epsilon$, l'indice
$\nu$ in genere aumenta; l'importante è che, \emph{comunque} si fissi $\epsilon>0$, un
tale $\nu$ esista.}
\label{fig:idea-limite}
\end{figure}

Abbiamo rappresentato nella \cref{fig:idea-limite} tre casi per diversi valori di
$\epsilon$; se si cambia $\epsilon$ con un numero minore, l'indice $\nu$ in genere
aumenta. L'importante è che, comunque si fissi $\epsilon>0$, in corrispondenza esista un
indice $\nu$ con le proprietà suddette.

\begin{definizione}[Intorno]\index{intorno}
Dati $a\in\R$ ed $\epsilon>0$, l'intervallo
\begin{equation}\label{eq:intorno}
  (a-\epsilon,a+\epsilon)=\{x\in\R : a-\epsilon<x<a+\epsilon\}=\{x\in\R : \abs{x-a}<\epsilon\}
\end{equation}
si chiama \emph{intorno} del numero $a$ di raggio $\epsilon$. L'ultima uguaglianza è la
\eqref{eq:mod-lt}: l'intorno di raggio $\epsilon$ è l'insieme dei punti che distano da
$a$ meno di $\epsilon$.
\end{definizione}

\begin{notastorica}[Archimede e il calcolo di $\pi$]
Nella \emph{Misura del cerchio} Archimede (287--212 a.C.) dimostrò che l'area del
cerchio è uguale a quella di un triangolo rettangolo con un cateto uguale al raggio e
l'altro uguale alla circonferenza --- in linguaggio moderno, $\pi r^{2}=\frac12\,r\cdot2\pi
r$ --- con un raffinato doppio ragionamento per assurdo (il \emph{metodo di esaustione})
che è, in sostanza, un passaggio al limite. Nella stessa opera calcolò il perimetro dei
poligoni regolari inscritti e circoscritti partendo dall'esagono e raddoppiando quattro
volte il numero dei lati, fino a $96$; ottenne così
\[
  3+\tfrac{10}{71}<\pi<3+\tfrac17,\qquad\text{cioè}\qquad 3{,}1408<\pi<3{,}1429 .
\]
Il punto geniale è che il procedimento \emph{non} richiede di conoscere $\pi$: le
grandezze dei poligoni con $2n$ lati si ottengono da quelle dei poligoni con $n$ lati
mediante sole operazioni aritmetiche ed estrazioni di radice quadrata (si veda il codice
qui sotto). Archimede eseguì a mano tutte le radici con approssimazioni razionali,
sempre per difetto o per eccesso nel verso giusto, per non compromettere le
disuguaglianze.
\end{notastorica}

\begin{python}[Il metodo di Archimede, senza conoscere $\pi$]
# s = semiperimetro del poligono inscritto, t = del circoscritto
# (raggio 1). Da n a 2n lati (media armonica e media geometrica):
#     t_2n = 2 s t / (s + t),      s_2n = sqrt(s * t_2n)
# Si parte dall'esagono: s_6 = 3, t_6 = 2*sqrt(3). Niente pi greco!
import math

n, s, t = 6, 3.0, 2 * math.sqrt(3)
print(f"{n:4d} lati: {s:.6f} < pi < {t:.6f}")
for _ in range(4):               # quattro raddoppi, fino a 96 lati
    t = 2 * s * t / (s + t)
    s = math.sqrt(s * t)
    n *= 2
    print(f"{n:4d} lati: {s:.6f} < pi < {t:.6f}")
print(f"Archimede: {223/71:.6f} < pi < {22/7:.6f}")
\end{python}
\begin{risultato}
   6 lati: 3.000000 < pi < 3.464102
  12 lati: 3.105829 < pi < 3.215390
  24 lati: 3.132629 < pi < 3.159660
  48 lati: 3.139350 < pi < 3.146086
  96 lati: 3.141032 < pi < 3.142715
Archimede: 3.140845 < pi < 3.142857
\end{risultato}

Le due successioni del programma, quella dei semiperimetri inscritti (crescente) e
quella dei circoscritti (decrescente), si stringono l'una contro l'altra intorno a $\pi$:
è una situazione che ritroveremo, in forma generale, nel teorema sulle successioni
monotòne (\cref{sec:monotone}) e nella definizione del numero $e$ (\cref{sec:numero-e}).

%% ----------------------------------------------------------------------------
%%  §3.2  Definizioni e prime proprietà
%% ----------------------------------------------------------------------------
\section{Definizione di limite e prime proprietà}
\label{sec:def-limite}

\begin{definizione}[Successione]\index{successione}
Una \emph{successione} è una legge che a ogni numero naturale $n$ fa corrispondere uno e
un solo numero reale $a_{n}$. Ricordando la definizione di funzione (\cref{sec:funzioni})
si può dire che \emph{una successione è una funzione da $\N$ in $\R$}. Indichiamo una
successione con il simbolo $(a_{n})$, o più semplicemente con $a_{n}$, o per esteso con
\begin{equation}\label{eq:succ-esteso}
  a_{1},\ a_{2},\ a_{3},\ a_{4},\ \dots,\ a_{n},\ \dots
\end{equation}
Il numero $a_{n}$ si chiama \emph{termine $n$-esimo} (o \emph{termine generale}) della
successione, e $n$ il suo \emph{indice}.
\end{definizione}

Parleremo allo stesso modo di successione anche se non sono definiti i termini $a_{n}$
per i primi indici $n$, cioè per un numero finito di indici, come nel caso considerato
nel paragrafo precedente, dove era $n\ge3$. Esempi di successioni sono:
\begin{align}
  a_{n}&=\frac1n: & &1,\ \tfrac12,\ \tfrac13,\ \tfrac14,\ \dots,\ \tfrac1n,\ \dots \label{eq:succ-1n}\\
  a_{n}&=\frac{n-1}{n}: & &0,\ \tfrac12,\ \tfrac23,\ \tfrac34,\ \dots,\ \tfrac{n-1}{n},\ \dots \label{eq:succ-n1n}\\
  a_{n}&=\frac{(-1)^{n}}{n}: & &-1,\ \tfrac12,\ -\tfrac13,\ \tfrac14,\ \dots,\ \tfrac{(-1)^{n}}{n},\ \dots \label{eq:succ-alt-n}\\
  a_{n}&=(-1)^{n}: & &-1,\ 1,\ -1,\ 1,\ \dots,\ (-1)^{n},\ \dots \label{eq:succ-alt}\\
  a_{n}&=n^{2}: & &1,\ 4,\ 9,\ 16,\ \dots,\ n^{2},\ \dots \label{eq:succ-n2}
\end{align}
I grafici delle successioni \eqref{eq:succ-alt-n} e \eqref{eq:succ-alt} sono riportati
nella \cref{fig:grafici-successioni}.

\subsection{Limite finito}

\begin{definizione}[Limite finito di una successione]\label{def:limite}\index{limite!di una successione}\index{successione!convergente}
Un numero reale $a$ è il \emph{limite} della successione $a_{n}$ (si dice anche che
$a_{n}$ \emph{tende} o \emph{converge} ad $a$), e si scrive
\begin{equation}\label{eq:lim-notazione}
  \lim_{n\to+\infty}a_{n}=a \qquad(\text{oppure}\quad a_{n}\to a),
\end{equation}
se, qualunque sia $\epsilon>0$, esiste un numero $\nu$ tale che
$a-\epsilon<a_{n}<a+\epsilon$ per ogni $n>\nu$.
\end{definizione}

\begin{figure}[htb]
\centering
\begin{tikzpicture}
\begin{axis}[width=5.8cm, height=4.4cm, xmin=0, xmax=16.5, ymin=-1.3, ymax=1.3,
  xtick={5,10,15}, ytick={-1,1}, title={\footnotesize $a_{n}=\frac{(-1)^{n}}{n}$}, title style={at={(0.62,0.92)}}]
  \addplot[only marks, mark=*, mark size=1.4pt, blu, samples at={1,...,16}] {(-1)^x/x};
\end{axis}
\end{tikzpicture}\hfill
\begin{tikzpicture}
\begin{axis}[width=5.8cm, height=4.4cm, xmin=0, xmax=16.5, ymin=-1.3, ymax=1.3,
  xtick={5,10,15}, ytick={-1,1}, title={\footnotesize $a_{n}=(-1)^{n}$}, title style={at={(0.62,0.92)}}]
  \addplot[only marks, mark=*, mark size=1.4pt, rame, samples at={1,...,16}] {(-1)^x};
\end{axis}
\end{tikzpicture}
\caption{Il grafico di una successione è un insieme di punti isolati $(n,a_{n})$. A
sinistra la successione \eqref{eq:succ-alt-n}, i cui termini oscillano attorno a $0$
avvicinandosi sempre più; a destra la successione \eqref{eq:succ-alt}, che oscilla
indefinitamente tra $-1$ e $1$.}
\label{fig:grafici-successioni}
\end{figure}

La relazione $a-\epsilon<a_{n}<a+\epsilon$ si può anche scrivere
$-\epsilon<a_{n}-a<\epsilon$, e ciò equivale, per la \eqref{eq:mod-lt}, a
$\abs{a_{n}-a}<\epsilon$. Quindi possiamo ripetere la definizione in simboli:
\begin{equation}\label{eq:def-limite}
  \lim_{n\to+\infty}a_{n}=a \quad\sse\quad
  \perogni\epsilon>0\ \ \esiste\nu : \ \abs{a_{n}-a}<\epsilon\quad\perogni n>\nu .
\end{equation}
Si dice anche che una proprietà vale \emph{definitivamente}\index{definitivamente} se vale
per ogni $n$ maggiore di un opportuno indice $\nu$. Con questo linguaggio:
\emph{$a_{n}\to a$ se e solo se, per ogni $\epsilon>0$, definitivamente
$\abs{a_{n}-a}<\epsilon$}. Geometricamente: comunque si fissi un intorno di $a$, i termini
della successione vi cadono definitivamente, cioè tutti tranne al più un numero finito.

\begin{attenzione}[L'ordine dei quantificatori]
Nella \eqref{eq:def-limite} l'ordine «$\perogni\epsilon\ \esiste\nu$» è essenziale:
l'indice $\nu$ si sceglie \emph{dopo} $\epsilon$, e dipende da esso (per sottolinearlo si
scrive talvolta $\nu=\nu(\epsilon)$). Se si scambiassero i quantificatori,
«$\esiste\nu\ \perogni\epsilon$: $\abs{a_{n}-a}<\epsilon$ per $n>\nu$», si otterrebbe
un'affermazione molto più forte: per $n>\nu$ il numero $\abs{a_{n}-a}$, non negativo,
sarebbe minore di \emph{ogni} $\epsilon>0$, e quindi nullo; la successione sarebbe
definitivamente \emph{uguale} ad $a$. Scambiare due quantificatori di tipo diverso
cambia quasi sempre il significato di una frase: «per ogni serratura esiste una chiave
che la apre» non è la stessa cosa che «esiste una chiave che apre ogni serratura».
\end{attenzione}

\begin{idea}[Il limite come tolleranza: il residuo di un solutore iterativo]
Chi ha usato un codice di fluidodinamica computazionale conosce bene questo schema. Un
solutore iterativo produce una successione di soluzioni approssimate $u_{1},u_{2},\dots$,
e l'utente fissa una \emph{tolleranza} sul residuo, ad esempio $10^{-6}$: il calcolo si
considera «converso» quando, da un certo numero di iterazioni in poi, l'errore resta sotto
la tolleranza. La definizione di limite dice la stessa cosa, ma per \emph{tutte} le
tolleranze contemporaneamente: la tolleranza è $\epsilon$, il numero di iterazioni
necessario è $\nu(\epsilon)$, e un metodo converge se, comunque si scelga la tolleranza,
esiste un numero di iterazioni oltre il quale l'errore non la supera più. Si noti il
«non la supera \emph{più}»: un residuo che scende sotto la soglia e poi risale non
soddisfa la definizione.
\end{idea}

Osserviamo, perché ci sarà utile nel seguito, che la \eqref{eq:def-limite} è equivalente
a
\begin{equation}\label{eq:def-limite-c}
  \lim_{n\to+\infty}a_{n}=a \quad\sse\quad \esiste c>0 :\ \perogni\epsilon>0\ \ \esiste\nu :\
  \abs{a_{n}-a}<c\,\epsilon\quad\perogni n>\nu ;
\end{equation}
infatti la \eqref{eq:def-limite} è il caso $c=1$ della \eqref{eq:def-limite-c}, e
viceversa dalla \eqref{eq:def-limite-c} si riottiene la \eqref{eq:def-limite} applicandola
con $\epsilon/c$ al posto di $\epsilon$. Analogamente, nella \eqref{eq:def-limite} ci si può
limitare a considerare i valori $\epsilon<\epsilon_{0}$, con $\epsilon_{0}>0$ fissato: se
la condizione vale per gli $\epsilon$ piccoli, vale a maggior ragione per quelli grandi
(l'indice $\nu$ trovato per $\epsilon_{0}/2$ va bene per ogni $\epsilon\ge\epsilon_{0}$).
Infine il numero $\nu$ si può sempre supporre naturale: se $\nu$ va bene, va bene anche
ogni numero più grande, in particolare un naturale maggiore di $\nu$, che esiste per la
proprietà di Archimede (\cref{teo:archimede}).

\subsection{Primi esempi}

\begin{esempio}\label{es:lim-1n}
Usando la definizione, verifichiamo che
\begin{equation}\label{eq:lim-1n}
  \lim_{n\to+\infty}\frac1n=0 .
\end{equation}
Risulta $\abs{a_{n}-a}=\abs{1/n}=1/n$. Dato che $1/n<\epsilon$ equivale a $n>1/\epsilon$,
basta scegliere $\nu=1/\epsilon$. Cioè abbiamo verificato che per ogni $\epsilon>0$
esiste $\nu=1/\epsilon$ per cui $\abs{a_{n}-a}=1/n<\epsilon$ per ogni $n>\nu$.
\end{esempio}

\begin{esempio}\label{es:lim-n1n}
Verifichiamo ora che
\begin{equation}\label{eq:lim-n1n}
  \lim_{n\to+\infty}\frac{n-1}{n}=1 .
\end{equation}
Risulta
\[
  \abs{a_{n}-a}=\Bigl|\frac{n-1}{n}-1\Bigr|=\Bigl|-\frac1n\Bigr|=\frac1n ;
\]
in questo esempio si vede l'importanza di considerare il valore assoluto di $a_{n}-a$: si
trova $\abs{a_{n}-a}=1/n$ e poi si procede come nel caso precedente.
\end{esempio}

\begin{esempio}[Un limite sbagliato]\label{es:lim-sbagliato}
Naturalmente non abbiamo scelto a caso i valori $a=0$ e $a=1$ nei due esempi precedenti.
Proviamo a vedere che cosa succede se tentassimo invece di dimostrare che
$\lim_{n\to+\infty}\frac1n=1$. Avremmo
\[
  \abs{a_{n}-a}=\Bigl|\frac1n-1\Bigr|=1-\frac1n<\epsilon
  \quad\sse\quad \frac1n>1-\epsilon .
\]
Questa relazione non crea problemi se $\epsilon$ è grande: ad esempio, se $\epsilon=1$ è
verificata da ogni $n\in\N$. Ma se $\epsilon$ è più piccolo, ad esempio $\epsilon=1/2$,
allora $1/n>1/2$ è verificata solo se $n<2$. Cioè solo $a_{1}$ verifica la relazione
data, e non gli $a_{n}$ con $n>\nu$, comunque si scelga $\nu$. Ciò prova che $a=1$
\emph{non} è il limite della successione $1/n$: la definizione deve valere per
\emph{ogni} $\epsilon>0$, e basta un solo $\epsilon$ per cui fallisce a farla cadere.
\end{esempio}

\begin{python}[Quanto deve essere grande $\nu$?]
# a_n = (2n+1)/(n+3) -> 2. Si ha |a_n - 2| = 5/(n+3), e la condizione
# |a_n - 2| < eps equivale a n > 5/eps - 3. Confrontiamo la formula
# con la ricerca "brutale" del primo indice da cui essa vale sempre.
from fractions import Fraction as F
import math

def a(n):
    return F(2 * n + 1, n + 3)

for eps in [F(1, 10), F(1, 100), F(1, 1000), F(1, 10**6)]:
    nu = max(0, math.floor(5 / eps - 3))   # disuguaglianza risolta
    primo = 1
    while abs(a(primo) - 2) >= eps:        # |a_n - 2| decresce
        primo += 1
    print(f"eps = {float(eps):.0e}:  nu = {nu:>7d},"
          f"  primo indice buono = {primo}")
\end{python}
\begin{risultato}
eps = 1e-01:  nu =      47,  primo indice buono = 48
eps = 1e-02:  nu =     497,  primo indice buono = 498
eps = 1e-03:  nu =    4997,  primo indice buono = 4998
eps = 1e-06:  nu = 4999997,  primo indice buono = 4999998
\end{risultato}

Il programma mostra concretamente la dipendenza di $\nu$ da $\epsilon$: per guadagnare
una cifra decimale di precisione occorre moltiplicare per $10$ l'indice. Si noti che la
ricerca «brutale» del primo indice buono funziona solo perché $\abs{a_{n}-2}$ decresce al
crescere di $n$; in generale un calcolatore non può verificare la definizione di limite,
che riguarda infiniti termini: per questo servono le dimostrazioni.

\subsection{Unicità del limite}

Con un argomento simile a quello dell'\cref{es:lim-sbagliato} verifichiamo che il limite,
se esiste, è uno solo: è ciò che giustifica l'articolo determinativo e la notazione
$\lim a_{n}=a$.

\begin{teorema}[Unicità del limite]\label{teo:unicita-limite}\index{limite!unicità}
Una successione convergente non può avere due limiti distinti.
\end{teorema}

\begin{schema}
\Ipotesi{$a_{n}\to a$ \ e \ $a_{n}\to b$, con $a,b\in\R$.}
\Tesi{$a=b$.}
\end{schema}

\begin{strategia}
Ragioniamo per assurdo: se fosse $a\neq b$, i due numeri avrebbero una distanza
$\abs{a-b}>0$, e potremmo scegliere due intorni di $a$ e di $b$ \emph{disgiunti}: basta
prendere entrambi di raggio $\epsilon=\abs{a-b}/2$, metà della distanza
(\cref{fig:unicita-limite}). Ma la successione dovrebbe cadere definitivamente in
entrambi, e un numero non può stare in due intervalli disgiunti. Per rendere rigoroso
il «non può stare in entrambi» usiamo la disuguaglianza triangolare: se $a_{n}$ distasse
meno di $\epsilon$ sia da $a$ sia da $b$, la distanza tra $a$ e $b$ sarebbe minore di
$2\epsilon=\abs{a-b}$. Un altro ingrediente, che useremo continuamente: se due proprietà
valgono definitivamente, valgono definitivamente \emph{insieme} (basta prendere il più
grande dei due indici).
\end{strategia}

\begin{dimostrazione}
Supponiamo per assurdo che esistano due limiti distinti, cioè supponiamo che
$a_{n}\to a$, $a_{n}\to b$, con $a\neq b$. Poniamo $\epsilon=\abs{a-b}/2$ $(>0)$. Si ha
\begin{equation}\label{eq:unicita-passo}
  \esiste\nu_{1}:\ \abs{a_{n}-a}<\epsilon\quad\perogni n>\nu_{1}; \qquad
  \esiste\nu_{2}:\ \abs{a_{n}-b}<\epsilon\quad\perogni n>\nu_{2}.
\end{equation}
Ponendo $\nu=\max\{\nu_{1},\nu_{2}\}$, le relazioni sopra scritte valgono
contemporaneamente per $n>\nu$, e si ha (utilizzando la disuguaglianza triangolare,
\cref{teo:triangolare})
\begin{equation}\label{eq:unicita-triangolare}
\begin{split}
  \abs{a-b}&=\abs{(a-a_{n})+(a_{n}-b)}\le\abs{a-a_{n}}+\abs{a_{n}-b}\\
  &=\abs{a_{n}-a}+\abs{a_{n}-b}<\epsilon+\epsilon=\abs{a-b}.
\end{split}
\end{equation}
Abbiamo così trovato che $\abs{a-b}<\abs{a-b}$, che è assurdo.
\end{dimostrazione}

\begin{figure}[htb]
\centering
\begin{tikzpicture}[x=1cm, every node/.style={font=\footnotesize}]
  \draw[-Latex, grigio] (-0.5,0) -- (9.5,0);
  \fill[blu!15] (1,-0.15) rectangle (4,0.15);
  \fill[rame!20] (4,-0.15) rectangle (7,0.15);
  \draw[blu, thick] (2.5,-0.12) -- (2.5,0.12) node[above=4pt] {$a$};
  \draw[rame, thick] (5.5,-0.12) -- (5.5,0.12) node[above=4pt] {$b$};
  \draw[blu, decorate, decoration={brace, amplitude=3pt, mirror}] (1,-0.3) -- (4,-0.3) node[midway, below=4pt] {$(a-\epsilon,a+\epsilon)$};
  \draw[rame, decorate, decoration={brace, amplitude=3pt, mirror}] (4.02,-0.3) -- (7,-0.3) node[midway, below=4pt] {$(b-\epsilon,b+\epsilon)$};
  \draw[grigio, decorate, decoration={brace, amplitude=3pt}] (2.5,0.85) -- (5.5,0.85) node[midway, above=3pt] {$\abs{a-b}=2\epsilon$};
\end{tikzpicture}
\caption{L'idea della dimostrazione del \cref{teo:unicita-limite}: con
$\epsilon=\abs{a-b}/2$ i due intorni sono disgiunti, e i termini della successione non
possono cadere definitivamente in entrambi.}
\label{fig:unicita-limite}
\end{figure}

\subsection{Successioni oscillanti}

Esaminiamo ora le successioni \eqref{eq:succ-alt-n} e \eqref{eq:succ-alt}. Si ha che
\begin{equation}\label{eq:lim-alt-n}
  \lim_{n\to+\infty}\frac{(-1)^{n}}{n}=0 ;
\end{equation}
infatti $\abs{a_{n}-a}=\abs{(-1)^{n}/n}=1/n$, e poi si procede come fatto per il limite
\eqref{eq:lim-1n}. Invece il limite della successione $(-1)^{n}$ non esiste.

\begin{proposizione}\label{prop:alterna-non-converge}
La successione $a_{n}=(-1)^{n}$ non ha limite finito.
\end{proposizione}

\begin{schema}
\Ipotesi{$a_{n}=(-1)^{n}$; \ $a\in\R$ arbitrario.}
\Tesi{$\esiste\epsilon>0$ tale che $\perogni\nu$ $\esiste n>\nu$ con
$\abs{a_{n}-a}\ge\epsilon$, cioè $a$ non è il limite di $a_{n}$.}
\end{schema}

\begin{strategia}
Dobbiamo negare la definizione: per ogni candidato $a$ va trovato un $\epsilon>0$ tale
che, oltre qualunque indice $\nu$, qualche termine disti da $a$ almeno $\epsilon$. Con
$\epsilon=1$ funziona: i valori $1$ e $-1$ distano $2$, e per la disuguaglianza
triangolare $a$ non può distare meno di $1$ da entrambi; poiché sia i termini di posto
pari sia quelli di posto dispari compaiono oltre ogni $\nu$, quello «lontano» da $a$
compare sempre.
\end{strategia}

\begin{dimostrazione}
Sia $a\in\R$ qualsiasi. Per la disuguaglianza triangolare
\[
  2=\abs{1-(-1)}=\abs{(1-a)+(a+1)}\le\abs{1-a}+\abs{-1-a},
\]
quindi almeno uno dei due numeri $\abs{1-a}$, $\abs{-1-a}$ è maggiore o uguale a $1$.
Supponiamo, ad esempio, $\abs{-1-a}\ge1$ (l'altro caso è analogo, con «pari» al posto di
«dispari»). Allora, scelto $\epsilon=1$, per ogni $\nu$ esiste un indice dispari $n>\nu$
(ad esempio $n=2m+1$, dove $m$ è un naturale maggiore di $\nu$, che esiste per la
proprietà di Archimede), e per tale $n$ si ha $\abs{a_{n}-a}=\abs{-1-a}\ge1=\epsilon$.
Dunque per $\epsilon=1$ nessun indice $\nu$ soddisfa la definizione: $a$ non è il limite.
Poiché $a$ era arbitrario, la successione non ha limite finito.
\end{dimostrazione}

\subsection{Limiti infiniti}

Infine la successione \eqref{eq:succ-n2} tende a $+\infty$ secondo la definizione
seguente.

\begin{definizione}[Limite infinito]\label{def:limite-infinito}\index{limite!infinito}\index{successione!divergente}
Una successione $a_{n}$ ha limite uguale a $+\infty$ (si dice anche che $a_{n}$
\emph{tende} o \emph{diverge} a $+\infty$), e si scrive
\begin{equation}\label{eq:lim-piu-infinito}
  \lim_{n\to+\infty}a_{n}=+\infty \qquad(\text{oppure}\quad a_{n}\to+\infty),
\end{equation}
se, qualunque sia $M>0$, esiste un numero $\nu$ tale che $a_{n}>M$ per ogni $n>\nu$. Si
dà una definizione analoga nel caso di limite uguale a $-\infty$. In simboli:
\begin{align}
  \lim_{n\to+\infty}a_{n}=+\infty \quad&\sse\quad \perogni M>0\ \ \esiste\nu :\ a_{n}>M \quad\perogni n>\nu; \label{eq:def-piu-inf}\\
  \lim_{n\to+\infty}a_{n}=-\infty \quad&\sse\quad \perogni M>0\ \ \esiste\nu :\ a_{n}<-M \quad\perogni n>\nu. \label{eq:def-meno-inf}
\end{align}
\end{definizione}

\begin{esempio}
$n^{2}\to+\infty$: fissato $M>0$, si ha $n^{2}\ge n>M$ per ogni $n>\nu=M$. Analogamente
$-n^{2}\to-\infty$. Anche la successione $a_{n}=n$ diverge a $+\infty$: è la proprietà di
Archimede, riletta nel linguaggio dei limiti.
\end{esempio}

\begin{marginfigure}
\centering
\begin{tikzpicture}
\begin{axis}[width=5.4cm, height=4.6cm, xmin=0, xmax=8.5, ymin=0, ymax=40,
  xtick={2,6,8}, extra x ticks={4.47}, extra x tick labels={$\nu$},
  ytick={20}, yticklabels={$M$}]
  \fill[verde!15] (axis cs:4.47,20) rectangle (axis cs:8.5,40);
  \addplot[tratteggio, domain=0:8.5] {20};
  \addplot[only marks, mark=*, mark size=1.4pt, blu, samples at={1,...,6}] {x^2};
  \addplot[tratteggio] coordinates {(4.47,0) (4.47,20)};
\end{axis}
\end{tikzpicture}
\caption{$a_{n}=n^{2}\to+\infty$: fissato un livello $M$, i termini con $n>\nu=\sqrt M$
stanno tutti al di sopra di $M$.}
\end{marginfigure}

\begin{definizione}[Terminologia]\index{successione!regolare}\index{successione!infinitesima}
Una successione si dice \emph{convergente} se ammette limite finito, mentre si dice
\emph{divergente} se ammette limite uguale a $+\infty$ oppure a $-\infty$. Le successioni
convergenti o divergenti si dicono \emph{regolari}, mentre le successioni che non
ammettono limite (né finito né infinito) si dicono \emph{non regolari} o
\emph{oscillanti}. Infine, una successione che converge a zero si dice
\emph{infinitesima}, mentre una successione divergente si dice anche \emph{infinita}.
\end{definizione}

La successione $(-1)^{n}$ è non regolare: non ha limite finito
(\cref{prop:alterna-non-converge}) e non diverge, perché è compresa tra $-1$ e $1$ (con
$M=1$ nessun termine supera $M$). Anche per i limiti infiniti vale l'unicità: una
successione non può tendere sia a $+\infty$ sia a $-\infty$ (con $M=1$, definitivamente
dovrebbe essere $a_{n}>1$ e $a_{n}<-1$), né contemporaneamente a un limite finito $a$ e a
$\pm\infty$ (con $\epsilon=1$ e $M=\abs a+1$, definitivamente sarebbe
$\abs{a_{n}}<\abs a+1$ e $\abs{a_{n}}>\abs a+1$).

\begin{notastorica}[Dalle «ultime ragioni» all'$\epsilon$]
Newton, nei \emph{Principia} (1687), parlava di «ultime ragioni di quantità
evanescenti», cioè del rapporto tra due quantità «nell'istante in cui svaniscono»: un
concetto efficace ma oscuro, che il vescovo George Berkeley ridicolizzò nel 1734
chiamando gli incrementi evanescenti «fantasmi di quantità defunte». Eppure, quasi
trent'anni prima dei \emph{Principia}, una definizione sorprendentemente vicina alla
nostra era già stata scritta. Nei \emph{Geometriae speciosae elementa} (1659) il
bolognese Pietro Mengoli chiamava «quasi infinita» una quantità variabile che può
superare ogni numero assegnato, «quasi nulla» una che può diventare minore di ogni
numero positivo, e «quasi $a$» una che può essere resa minore di ogni numero maggiore di
$a$ e maggiore di ogni numero minore di $a$: è, in sostanza, la condizione
$a-\epsilon<a_{n}<a+\epsilon$. Ne ricavò anche le regole per i limiti di somme, prodotti e
quozienti (\cref{sec:operazioni-limiti}); ma le sue idee restarono ignorate, anche perché
l'interesse per le successioni sarebbe nato molto più tardi (sulle sue ricerche sulle
serie si veda la nota storica del \S\ref{sec:serie-numeriche}). D'Alembert,
nell'\emph{Encyclopédie} (1765), definì il limite a parole come una grandezza a cui un'altra «si avvicina più di
qualsiasi quantità data», e Cauchy, nel \emph{Cours d'analyse} (1821) da cui è tratta
l'epigrafe del capitolo, ne fece il fondamento di tutta l'Analisi. Ma la definizione nella
forma che abbiamo dato, con i quantificatori «per ogni $\epsilon$ esiste $\nu$» e le
disuguaglianze esplicite, si deve a Karl Weierstrass (1815--1897), che la insegnò nelle
sue lezioni berlinesi a partire dagli anni Sessanta dell'Ottocento; i suoi allievi la
diffusero. La lettera $\epsilon$ era già usata da Cauchy (\emph{Résumé des leçons}, 1823); per
tradizione la si legge come iniziale di \emph{erreur}, «errore».
\end{notastorica}

%% ----------------------------------------------------------------------------
%%  §3.3  Successioni limitate. Operazioni con i limiti
%% ----------------------------------------------------------------------------
\section{Successioni limitate e operazioni con i limiti}
\label{sec:operazioni-limiti}

\subsection{Successioni limitate}

Abbiamo detto nel paragrafo precedente che una successione si dice \emph{regolare} se
ammette limite (finito o infinito). Introduciamo ora una nozione di natura diversa, che
non riguarda il comportamento «per $n$ grande» ma l'insieme di tutti i valori.

\begin{definizione}[Successione limitata]\index{successione!limitata}
Una successione $a_{n}$ si dice \emph{limitata} se esiste un numero reale $M$ tale che
\begin{equation}\label{eq:succ-limitata}
  \abs{a_{n}}\le M \qquad\perogni n\in\N,
\end{equation}
o, equivalentemente, in base alla \cref{prop:modulo-r},
\begin{equation}\label{eq:succ-limitata-2}
  -M\le a_{n}\le M \qquad\perogni n\in\N.
\end{equation}
Si dice \emph{limitata superiormente} (risp.\ \emph{inferiormente}) se esiste $L$ tale che
$a_{n}\le L$ (risp.\ $a_{n}\ge L$) per ogni $n$.
\end{definizione}

Si noti il legame con la definizione \eqref{eq:limitato} di insieme limitato e con la
\cref{prop:limitato-modulo}: una successione è limitata se e solo se è limitato l'insieme
$\{a_{n} : n\in\N\}$ dei suoi valori. Il lettore provi a verificare per esercizio che le
successioni \eqref{eq:succ-1n}, \eqref{eq:succ-n1n}, \eqref{eq:succ-alt-n},
\eqref{eq:succ-alt} sono limitate, mentre la successione \eqref{eq:succ-n2} non lo è; ad
esempio, la successione \eqref{eq:succ-alt-n} è limitata dalla costante $M=1$, infatti
\[
  \abs{a_{n}}=\Bigl|\frac{(-1)^{n}}{n}\Bigr|=\frac1n\le1 \qquad\perogni n\in\N .
\]

\emph{Esistono successioni limitate non regolari}, cioè successioni limitate che non
ammettono limite; un esempio è fornito dalla successione $a_{n}=(-1)^{n}$, che è
limitata (perché $\abs{a_{n}}=1$ per ogni $n\in\N$) ma che, come dimostrato nella
\cref{prop:alterna-non-converge}, non ammette limite. Viceversa, ogni successione che
ammette limite finito è limitata, come mostrato dal seguente teorema.

\begin{teorema}\label{teo:convergente-limitata}\index{successione!convergente!è limitata}
Ogni successione convergente è limitata.
\end{teorema}

\begin{schema}
\Ipotesi{$a_{n}\to a\in\R$.}
\Tesi{$\esiste M\ge0 : \abs{a_{n}}\le M$ per ogni $n\in\N$.}
\end{schema}

\begin{strategia}
La definizione di limite controlla i termini solo \emph{definitivamente}, cioè per
$n>\nu$: con $\epsilon=1$ essi stanno tutti nell'intorno $(a-1,a+1)$, e quindi hanno
modulo minore di $1+\abs a$. Restano fuori i primi termini $a_{1},\dots,a_{\nu}$; ma sono
in numero \emph{finito}, e un insieme finito di numeri ha sempre un massimo. Prendendo
come $M$ il più grande tra questi valori e $1+\abs a$ si controllano tutti i termini.
È uno schema che ricorrerà spesso: \emph{un controllo definitivo più un numero finito di
eccezioni dà un controllo globale}.
\end{strategia}

\begin{dimostrazione}
Supponiamo che $a_{n}$ converga ad $a$ e scegliamo $\epsilon=1$. In base alla definizione
di limite \eqref{eq:def-limite} esiste un indice $\nu$, che possiamo supporre naturale,
per cui $\abs{a_{n}-a}<1$ per ogni $n>\nu$. Quindi, utilizzando la disuguaglianza
triangolare \eqref{eq:triangolare},
\begin{equation}\label{eq:conv-lim-passo}
  \abs{a_{n}}=\abs{(a_{n}-a)+a}\le\abs{a_{n}-a}+\abs a<1+\abs a \qquad\perogni n>\nu .
\end{equation}
Ma allora, per ogni $n\in\N$, si ha
\begin{equation}\label{eq:conv-lim-M}
  \abs{a_{n}}\le M=\max\{\abs{a_{1}},\abs{a_{2}},\dots,\abs{a_{\nu}},1+\abs a\},
\end{equation}
dove il massimo esiste perché si tratta di un insieme finito di numeri.
\end{dimostrazione}

\begin{osservazione}
Il viceversa del \cref{teo:convergente-limitata} è falso, come mostra $(-1)^{n}$; e
anche una successione divergente non è limitata (se $a_{n}\to+\infty$, nessun $M$ può
maggiorare tutti i termini). Vedremo però nel \cref{sec:bw-cauchy}, con il teorema di
Bolzano--Weierstrass, che da ogni successione limitata si può sempre «estrarre» una
successione convergente.
\end{osservazione}

\subsection{Le operazioni con i limiti}

Valgono le seguenti regole di calcolo, che permettono di calcolare limiti di successioni
complicate a partire da limiti noti di successioni semplici, senza ricorrere ogni volta
alla definizione.

\begin{teorema}[Operazioni con i limiti]\label{teo:operazioni-limiti}\index{limite!operazioni}
Se $\displaystyle\lim_{n\to+\infty}a_{n}=a$ \ e \ $\displaystyle\lim_{n\to+\infty}b_{n}=b$,
con $a,b\in\R$, si ha
\begin{align}
  &\lim_{n\to+\infty}(a_{n}\pm b_{n})=a\pm b; \label{eq:lim-somma}\\
  &\lim_{n\to+\infty}a_{n}\cdot b_{n}=a\,b; \label{eq:lim-prodotto}\\
  &\lim_{n\to+\infty}\frac{a_{n}}{b_{n}}=\frac ab \qquad(\text{se } b_{n}\neq0 \text{ per ogni } n, \text{ e } b\neq0). \label{eq:lim-quoziente}
\end{align}
\end{teorema}

\begin{schema}
\Ipotesi{$a_{n}\to a\in\R$, \ $b_{n}\to b\in\R$; per il quoziente, inoltre,
$b_{n}\neq0$ per ogni $n$ e $b\neq0$.}
\Tesi{$a_{n}\pm b_{n}\to a\pm b$, \ $a_{n}b_{n}\to ab$, \ $a_{n}/b_{n}\to a/b$.}
\end{schema}

\begin{strategia}
In tutti e tre i casi dobbiamo stimare la distanza dal limite, $\abs{(\text{risultato})-
(\text{limite atteso})}$, in termini delle due distanze $\abs{a_{n}-a}$ e $\abs{b_{n}-b}$,
che sappiamo rendere piccole quanto vogliamo. L'arte sta nello scrivere la differenza
in modo che compaiano proprio queste due quantità:
\begin{itemize}
  \item per la somma basta riordinare: $(a_{n}+b_{n})-(a+b)=(a_{n}-a)+(b_{n}-b)$, e la
  disuguaglianza triangolare dà la stima con $2\epsilon$;
  \item per il prodotto si aggiunge e si toglie un termine «misto»:
  $a_{n}b_{n}-ab=a_{n}(b_{n}-b)+b(a_{n}-a)$, e per controllare il fattore $a_{n}$ serve
  sapere che la successione è limitata (\cref{teo:convergente-limitata});
  \item per il quoziente ci si riduce al prodotto $a_{n}\cdot\frac1{b_{n}}$, e la
  difficoltà è il reciproco: occorre impedire che $b_{n}$ si avvicini a zero, cosa che
  otterremo mostrando che definitivamente $\abs{b_{n}}>\abs b/2$.
\end{itemize}
In ogni caso, per avere entrambe le stime contemporaneamente, si sceglie come indice
$\nu$ il più grande dei due indici forniti dalle ipotesi, e alla fine si usa la forma
\eqref{eq:def-limite-c} della definizione, che ammette un fattore costante davanti a
$\epsilon$.
\end{strategia}

\begin{dimostrazione}[Dimostrazione della \eqref{eq:lim-somma}]
Dimostriamo la \eqref{eq:lim-somma} con il segno $+$. Per ipotesi, per ogni $\epsilon>0$,
\begin{equation}\label{eq:ipotesi-ab}
  \esiste\nu_{1}:\ \abs{a_{n}-a}<\epsilon\quad\perogni n>\nu_{1}; \qquad
  \esiste\nu_{2}:\ \abs{b_{n}-b}<\epsilon\quad\perogni n>\nu_{2}.
\end{equation}
Ponendo $\nu=\max\{\nu_{1},\nu_{2}\}$, per ogni $n>\nu$ si ha
\begin{equation}\label{eq:somma-passo}
  \abs{(a_{n}+b_{n})-(a+b)}=\abs{(a_{n}-a)+(b_{n}-b)}\le\abs{a_{n}-a}+\abs{b_{n}-b}<2\epsilon .
\end{equation}
La prova è completa per la \eqref{eq:def-limite-c} con $c=2$. Si noti che abbiamo usato la
disuguaglianza triangolare e la definizione di limite nella forma \eqref{eq:def-limite-c}.
Con il segno $-$ si procede allo stesso modo, scrivendo
$(a_{n}-b_{n})-(a-b)=(a_{n}-a)-(b_{n}-b)$ e usando $\abs{x-y}\le\abs x+\abs y$.
\end{dimostrazione}

Dimostriamo la \eqref{eq:lim-prodotto} relativa al limite di un prodotto con due metodi,
il primo dei quali non fa uso del concetto di successione limitata.

\begin{dimostrazione}[Dimostrazione della \eqref{eq:lim-prodotto}, primo metodo]
Utilizziamo l'identità
\begin{equation}\label{eq:identita-prodotto}
\begin{split}
  a_{n}b_{n}-ab&=(a_{n}-a)(b_{n}-b)+a_{n}b+ab_{n}-2ab\\
  &=(a_{n}-a)(b_{n}-b)+(a_{n}-a)\,b+a\,(b_{n}-b).
\end{split}
\end{equation}
Nell'ipotesi \eqref{eq:ipotesi-ab}, e supponendo $\epsilon\le1$ (il che è lecito, come
osservato dopo la \eqref{eq:def-limite-c}), per ogni $n>\nu=\max\{\nu_{1},\nu_{2}\}$
risulta quindi
\begin{equation}\label{eq:prodotto-passo-1}
\begin{split}
  \abs{a_{n}b_{n}-ab}&\le\abs{a_{n}-a}\,\abs{b_{n}-b}+\abs{a_{n}-a}\,\abs b+\abs a\,\abs{b_{n}-b}\\
  &<\epsilon^{2}+\epsilon\abs b+\abs a\epsilon\le(1+\abs a+\abs b)\,\epsilon,
\end{split}
\end{equation}
e la tesi segue dalla \eqref{eq:def-limite-c} con $c=1+\abs a+\abs b$.
\end{dimostrazione}

\begin{dimostrazione}[Dimostrazione della \eqref{eq:lim-prodotto}, secondo metodo]
Per il \cref{teo:convergente-limitata} la successione $a_{n}$ è limitata, cioè esiste un
numero reale $M>0$ tale che
\begin{equation}\label{eq:an-limitata}
  \abs{a_{n}}\le M \qquad\perogni n\in\N .
\end{equation}
Dall'ipotesi \eqref{eq:ipotesi-ab}, per ogni $n>\nu=\max\{\nu_{1},\nu_{2}\}$ si ottiene
\begin{equation}\label{eq:prodotto-passo-2}
\begin{split}
  \abs{a_{n}b_{n}-ab}&=\abs{a_{n}b_{n}-a_{n}b+a_{n}b-ab}=\abs{a_{n}(b_{n}-b)+b(a_{n}-a)}\\
  &\le\abs{a_{n}}\,\abs{b_{n}-b}+\abs b\,\abs{a_{n}-a}<M\epsilon+\abs b\epsilon=(M+\abs b)\,\epsilon,
\end{split}
\end{equation}
e si conclude con la \eqref{eq:def-limite-c}, con $c=M+\abs b$.
\end{dimostrazione}

Il testo di riferimento \cite{marcellini-sbordone} si limita ad affermare che la prova per
il quoziente è «simile» a quella del prodotto. La difficoltà nuova è il denominatore, che
dobbiamo tenere lontano da zero; la isoliamo in un lemma, preceduto da una semplice
conseguenza della disuguaglianza triangolare.

\begin{lemma}[Limite del valore assoluto]\label{lem:lim-modulo}
Se $a_{n}\to a$, allora $\abs{a_{n}}\to\abs a$.
\end{lemma}

\begin{schema}
\Ipotesi{$a_{n}\to a\in\R$.}
\Tesi{$\abs{a_{n}}\to\abs a$.}
\end{schema}

\begin{strategia}
La distanza tra i moduli non supera quella tra i numeri (disuguaglianza triangolare
«inversa»): l'indice $\nu$ che serve per $a_{n}$ serve anche per $\abs{a_{n}}$.
\end{strategia}

\begin{dimostrazione}
Per il \cref{cor:triangolare-inversa},
$\bigl|\abs{a_{n}}-\abs a\bigr|\le\abs{a_{n}-a}$, e il secondo membro è definitivamente
minore di $\epsilon$, per ogni $\epsilon>0$.
\end{dimostrazione}

\begin{lemma}[Il denominatore resta lontano da zero]\label{lem:denominatore}
Se $b_{n}\to b$ con $b\neq0$, allora esiste $\nu$ tale che $\abs{b_{n}}>\dfrac{\abs b}{2}$
per ogni $n>\nu$. Se inoltre $b_{n}\neq0$ per ogni $n$, allora
$\dfrac1{b_{n}}\to\dfrac1b$.
\end{lemma}

\begin{schema}
\Ipotesi{$b_{n}\to b\in\R$, \ $b\neq0$; per la seconda parte, inoltre, $b_{n}\neq0$ per
ogni $n$.}
\Tesi{$\abs{b_{n}}>\abs b/2$ definitivamente; \ $1/b_{n}\to1/b$.}
\end{schema}

\begin{strategia}
Per la prima affermazione basta applicare la definizione di limite alla successione
$\abs{b_{n}}\to\abs b$ (\cref{lem:lim-modulo}) con $\epsilon=\abs b/2$: definitivamente
$\abs{b_{n}}$ sta nell'intorno $(\abs b/2,\,3\abs b/2)$, quindi supera $\abs b/2$. Per la
seconda, scriviamo la differenza dei reciproci con un denominatore comune:
$\frac1{b_{n}}-\frac1b=\frac{b-b_{n}}{b_{n}b}$. Il numeratore è piccolo, e il
denominatore, grazie alla prima parte, non è troppo piccolo: $\abs{b_{n}b}>\abs b^{2}/2$.
\end{strategia}

\begin{dimostrazione}
Per il \cref{lem:lim-modulo}, $\abs{b_{n}}\to\abs b>0$. Scelto $\epsilon=\abs b/2$, esiste
$\nu_{0}$ tale che $\bigl|\abs{b_{n}}-\abs b\bigr|<\abs b/2$ per $n>\nu_{0}$; in
particolare $\abs{b_{n}}>\abs b-\abs b/2=\abs b/2$ per $n>\nu_{0}$.

Sia ora $\epsilon>0$ e sia $\nu_{2}$ tale che $\abs{b_{n}-b}<\epsilon$ per $n>\nu_{2}$.
Per $n>\max\{\nu_{0},\nu_{2}\}$ si ha
\begin{equation}\label{eq:reciproco-passo}
  \Bigl|\frac1{b_{n}}-\frac1b\Bigr|=\frac{\abs{b-b_{n}}}{\abs{b_{n}}\,\abs b}
  <\frac{\epsilon}{\frac{\abs b}{2}\,\abs b}=\frac{2}{\abs b^{2}}\,\epsilon,
\end{equation}
e la tesi segue dalla \eqref{eq:def-limite-c}, con $c=2/\abs b^{2}$.
\end{dimostrazione}

\begin{dimostrazione}[Dimostrazione della \eqref{eq:lim-quoziente}]
Per il \cref{lem:denominatore} si ha $1/b_{n}\to1/b$, e applicando la
\eqref{eq:lim-prodotto} alle successioni $a_{n}$ e $1/b_{n}$ si ottiene
\[
  \frac{a_{n}}{b_{n}}=a_{n}\cdot\frac1{b_{n}}\to a\cdot\frac1b=\frac ab. \qedhere
\]
\end{dimostrazione}

\begin{osservazione}
Sono casi particolari del teorema: se $c\in\R$ è una costante, la successione costante
$c_{n}=c$ converge a $c$ (per ogni $\epsilon$ va bene qualsiasi $\nu$), e quindi
$c\,a_{n}\to c\,a$ e $a_{n}+c\to a+c$. L'ipotesi $b_{n}\neq0$ per ogni $n$ serve solo a
garantire che il quoziente sia definito per tutti gli indici; per il
\cref{lem:denominatore} essa è comunque soddisfatta definitivamente, e si potrebbe
considerare la successione $a_{n}/b_{n}$ definita solo da un certo indice in poi.
\end{osservazione}

\begin{idea}[La propagazione degli errori]
La dimostrazione del limite del prodotto è la stessa che si usa in laboratorio per la
\emph{propagazione degli errori}. Se si misura una forza $F$ con un errore $\delta F$ e
una velocità $V$ con un errore $\delta V$, la potenza $P=FV$ calcolata con i valori
misurati differisce da quella vera di
\[
  \abs{(F+\delta F)(V+\delta V)-FV}\le\abs F\,\abs{\delta V}+\abs V\,\abs{\delta F}+\abs{\delta F}\,\abs{\delta V},
\]
che è la \eqref{eq:prodotto-passo-1} con altri nomi. Dividendo per $\abs{FV}$ si ritrova la
regola pratica «gli errori relativi si sommano» (a meno del termine di secondo ordine
$\abs{\delta F\,\delta V}$). Il teorema sui limiti dice che, se gli errori di misura
possono essere resi piccoli a piacere, anche l'errore sulla potenza può esserlo.
\end{idea}

\subsection{Esempi}

Esaminiamo alcuni esempi di applicazione delle operazioni con i limiti.

\begin{esempio}\label{es:n-alla-meno-b}
Avendo già verificato che la successione $1/n$ converge a zero, dal limite del prodotto
\eqref{eq:lim-prodotto} con $a_{n}=b_{n}=1/n$ si deduce che
\begin{equation}\label{eq:lim-1n2}
  \lim_{n\to+\infty}\frac1{n^{2}}=0
\end{equation}
e, iterando il procedimento (cioè per induzione su $b$: $n^{-(b+1)}=n^{-b}\cdot n^{-1}$),
\begin{equation}\label{eq:lim-n-meno-b}
  \lim_{n\to+\infty}n^{-b}=0 \qquad\perogni b\in\N .
\end{equation}
\end{esempio}

\begin{esempio}[Quoziente di polinomi]\label{es:quoziente-polinomi}
Mediante le operazioni con i limiti si calcola, ad esempio, il seguente limite, dividendo
numeratore e denominatore per $n^{2}$:
\begin{equation}\label{eq:lim-razionale}
  \lim_{n\to+\infty}\frac{n^{2}+5n-4}{3n^{2}+1}
  =\lim_{n\to+\infty}\frac{1+5/n-4/n^{2}}{3+1/n^{2}}=\frac13 .
\end{equation}
Il passaggio decisivo è il primo: così com'è scritta, la frazione è il quoziente di due
successioni divergenti, a cui il \cref{teo:operazioni-limiti} non si applica; dopo la
divisione per $n^{2}$ numeratore e denominatore convergono (a $1$ e a $3\neq0$), e il
teorema dà il risultato. Si noti che non è un caso che il limite sia il rapporto dei
coefficienti dei termini di grado massimo: è un fatto generale per i quozienti di
polinomi dello stesso grado.
\end{esempio}

\begin{python}[Controllare un limite con \texttt{sympy} e numericamente]
import sympy as sp

n = sp.symbols("n", positive=True, integer=True)
a = (n**2 + 5*n - 4) / (3*n**2 + 1)
print(sp.limit(a, n, sp.oo))           # calcolo simbolico del limite

# Controllo numerico: la distanza dal limite 1/3 si riduce come 1/n
terzo = sp.Rational(1, 3)
for k in [10, 100, 1000, 10**4]:
    v = a.subs(n, k)
    print(f"n = {k:>5d}:  a_n = {float(v):.8f},"
          f"  |a_n - 1/3| = {float(abs(v - terzo)):.2e}")
\end{python}
\begin{risultato}
1/3
n =    10:  a_n = 0.48504983,  |a_n - 1/3| = 1.52e-01
n =   100:  a_n = 0.34985500,  |a_n - 1/3| = 1.65e-02
n =  1000:  a_n = 0.33499856,  |a_n - 1/3| = 1.67e-03
n = 10000:  a_n = 0.33349999,  |a_n - 1/3| = 1.67e-04
\end{risultato}

%% ----------------------------------------------------------------------------
%%  §3.4  Limiti infiniti e forme indeterminate
%% ----------------------------------------------------------------------------
\section{Limiti infiniti e forme indeterminate}
\label{sec:forme-indeterminate}

Il \cref{teo:operazioni-limiti} riguarda successioni con limiti \emph{finiti}. Quando
uno o entrambi i limiti sono infiniti, alcune regole continuano a valere in una forma
naturale, che conviene ricordare con una specie di «aritmetica dei simboli $\pm\infty$»;
altre, invece, non danno alcuna informazione.

\begin{teorema}[Operazioni con limiti infiniti]\label{teo:operazioni-infinite}
Siano $a\in\R$ e $a_{n}$, $b_{n}$ due successioni. Valgono le implicazioni
\begin{align}
  &a_{n}\to a, &&b_{n}\to\pm\infty &&\implica\ a_{n}+b_{n}\to\pm\infty; \label{eq:inf1}\\
  &a_{n}\to\pm\infty, &&b_{n}\to\pm\infty &&\implica\ a_{n}+b_{n}\to\pm\infty; \label{eq:inf2}\\
  &a_{n}\to a\neq0, &&b_{n}\to\pm\infty &&\implica\ \abs{a_{n}b_{n}}\to+\infty; \label{eq:inf3}\\
  &a_{n}\to\pm\infty, &&b_{n}\to\pm\infty &&\implica\ \abs{a_{n}b_{n}}\to+\infty; \label{eq:inf4}\\
  &a_{n}\to a, &&b_{n}\to\pm\infty &&\implica\ a_{n}/b_{n}\to0; \label{eq:inf5}\\
  &a_{n}\to a, &&b_{n}\to\pm\infty &&\implica\ \abs{b_{n}/a_{n}}\to+\infty \quad(a_{n}\neq0); \label{eq:inf6}\\
  &a_{n}\to a\neq0, &&b_{n}\to0 &&\implica\ \abs{a_{n}/b_{n}}\to+\infty \quad(b_{n}\neq0). \label{eq:inf7}
\end{align}
Nella \eqref{eq:inf2} i due segni sono concordi; nelle \eqref{eq:inf5} e
\eqref{eq:inf6} si sottintende che $b_{n}\neq0$ e, rispettivamente, $a_{n}\neq0$, perché i
quozienti abbiano senso. Nelle \eqref{eq:inf3}, \eqref{eq:inf4}, \eqref{eq:inf6},
\eqref{eq:inf7} il segno del limite (quando esiste) si determina con la regola dei segni:
ad esempio, se $a>0$ e $b_{n}\to-\infty$, allora $a_{n}b_{n}\to-\infty$.
\end{teorema}

\begin{schema}
\Ipotesi{in ciascuna riga, i limiti di $a_{n}$ e $b_{n}$ scritti a sinistra (con
$a\in\R$, e $a\neq0$ dove indicato); i denominatori diversi da zero.}
\Tesi{il comportamento della somma, del prodotto o del quoziente scritto a destra,
\eqref{eq:inf1}--\eqref{eq:inf7}.}
\end{schema}

\begin{strategia}
Le dimostrazioni seguono tutte lo stesso schema: per dimostrare che una successione
diverge occorre, fissato un livello $M$ arbitrariamente grande, trovare un indice oltre
il quale la successione lo supera. La successione divergente $b_{n}$ «trascina» il
risultato: basta chiedere a $b_{n}$ di superare un livello un po' più grande di $M$, che
compensi il contributo dell'altra successione, la quale --- essendo convergente --- è
definitivamente compresa tra due numeri fissi. Per i quozienti con limite zero si
procede in modo simmetrico: si rende il denominatore così grande che la frazione diventa
minore di qualunque $\epsilon$ assegnato.
\end{strategia}

\begin{dimostrazione}
Proviamo le \eqref{eq:inf1}, \eqref{eq:inf3}, \eqref{eq:inf5} e \eqref{eq:inf7} nel caso
del segno $+$; le altre si dimostrano con argomenti dello stesso tipo e sono lasciate per
esercizio.

\eqref{eq:inf1}: poiché $a_{n}\to a$, scelto $\epsilon=1$ si ha definitivamente
$a_{n}>a-1$. Fissato $M>0$, poiché $b_{n}\to+\infty$ si ha definitivamente
$b_{n}>M+\abs{a-1}$ (lo si richiede con il livello $M+\abs{a-1}>0$). Per gli $n$ per cui
valgono entrambe le disuguaglianze, cioè definitivamente,
\[
  a_{n}+b_{n}>(a-1)+M+\abs{a-1}\ge M .
\]

\eqref{eq:inf3}: per il \cref{lem:lim-modulo}, $\abs{a_{n}}\to\abs a>0$, e per il
\cref{lem:denominatore} (applicato ad $a_{n}$) definitivamente $\abs{a_{n}}>\abs a/2$.
Fissato $M>0$, definitivamente $\abs{b_{n}}\ge b_{n}>2M/\abs a$; quindi definitivamente
$\abs{a_{n}b_{n}}>\frac{\abs a}{2}\cdot\frac{2M}{\abs a}=M$.

\eqref{eq:inf5}: per il \cref{teo:convergente-limitata} esiste $K>0$ con $\abs{a_{n}}\le K$
per ogni $n$. Fissato $\epsilon>0$, definitivamente $b_{n}>K/\epsilon>0$, e quindi
$\abs{a_{n}/b_{n}}<K\cdot\frac{\epsilon}{K}=\epsilon$.

\eqref{eq:inf7}: definitivamente $\abs{a_{n}}>\abs a/2$ (come sopra). Fissato $M>0$, poiché
$b_{n}\to0$, definitivamente $\abs{b_{n}}<\frac{\abs a}{2M}$; allora definitivamente
$\abs{a_{n}/b_{n}}>\frac{\abs a}{2}\cdot\frac{2M}{\abs a}=M$.
\end{dimostrazione}

\begin{attenzione}[$1/0$ non è $+\infty$]
La \eqref{eq:inf7} afferma che il \emph{valore assoluto} del quoziente diverge, non il
quoziente stesso. Ad esempio $b_{n}=\frac{(-1)^{n}}{n}\to0$, ma $\frac1{b_{n}}=(-1)^{n}n$
assume alternativamente valori positivi e negativi sempre più grandi in modulo: non
tende né a $+\infty$ né a $-\infty$. Vale però il seguente fatto: \emph{se $b_{n}\to0$, allora
$1/b_{n}\to+\infty$ se e solo se $b_{n}>0$ definitivamente} (e $1/b_{n}\to-\infty$ se e solo
se $b_{n}<0$ definitivamente). Infatti, se $1/b_{n}\to+\infty$, definitivamente
$1/b_{n}>0$, cioè $b_{n}>0$; viceversa, se $b_{n}>0$ definitivamente, allora
definitivamente $1/b_{n}=\abs{1/b_{n}}$, che diverge a $+\infty$ per la \eqref{eq:inf7}.
\end{attenzione}

\subsection{Le forme indeterminate}

Risultano esclusi dal \cref{teo:operazioni-infinite} alcuni casi, che schematizziamo
nelle forme seguenti, dette \emph{forme indeterminate}:\index{forma indeterminata}
\begin{equation}\label{eq:forme-indeterminate}
  \infty-\infty, \qquad 0\cdot\infty, \qquad \frac\infty\infty, \qquad \frac00 .
\end{equation}
Altre forme indeterminate, di tipo esponenziale, saranno introdotte nel
\cref{sec:numero-e}. Dire che un limite è una forma indeterminata \emph{non} significa dire
che il limite non esiste: significa semplicemente che la conoscenza dei limiti delle
singole successioni non basta a determinarlo, e che occorre preliminarmente eseguire
trasformazioni o semplificazioni per eliminare, se possibile, l'indeterminazione.

\begin{esempio}[Perché «indeterminata»]\label{es:zero-per-infinito}
Tutte le successioni seguenti hanno la forma $0\cdot\infty$, cioè sono prodotti di una
successione infinitesima per una divergente; eppure hanno comportamenti completamente
diversi:
\[
\renewcommand{\arraystretch}{1.6}
\begin{array}{ll}
  \dfrac1n\cdot n^{2}=n\to+\infty, \qquad & -\dfrac1n\cdot n^{2}=-n\to-\infty,\\[2ex]
  \dfrac1n\cdot 3n=3\to3, & \dfrac1{n^{2}}\cdot n=\dfrac1n\to0,\\[2ex]
  \multicolumn{2}{l}{\dfrac{(-1)^{n}}{n}\cdot n=(-1)^{n}\ \text{ non regolare}.}
\end{array}
\]
Il risultato dipende dalla «velocità» con cui i due fattori tendono ai rispettivi limiti,
non dai limiti soltanto. Lo stesso accade per le altre forme: $n-n\to0$,
$(n+1)-n\to1$, $n^{2}-n\to+\infty$ sono tutte forme $\infty-\infty$.
\end{esempio}

\begin{esempio}\label{es:forme-indeterminate}
Le successioni seguenti esprimono forme indeterminate:
\begin{equation}\label{eq:es-forme}
  (n+1)^{2}-(n-1)^{2}; \qquad \frac1{n^{2}}\cdot(n+1); \qquad \frac{n-1}{n}; \qquad
  \frac{1/n^{2}}{1/n}.
\end{equation}
La prima è una forma $+\infty-\infty$; però svolgendo i quadrati si trova che la
successione vale $4n$, e quindi tende a $+\infty$. La seconda è una forma $0\cdot\infty$,
ma si può anche scrivere $\frac1n+\frac1{n^{2}}$, che tende a zero. La terza è una forma
$\infty/\infty$, ma già sappiamo che tende a $1$ \eqref{eq:lim-n1n}. La quarta è del tipo
$0/0$, ma semplificando si trova $1/n\to0$.
\end{esempio}

\begin{esempio}[Razionalizzare]\label{es:razionalizzare}
La successione $a_{n}=\sqrt{n+1}-\sqrt n$ è una forma $\infty-\infty$ (che
$\sqrt n\to+\infty$ si vede subito: $\sqrt n>M$ per $n>M^{2}$). Moltiplicando e dividendo
per $\sqrt{n+1}+\sqrt n$ si ottiene
\[
  \sqrt{n+1}-\sqrt n=\frac{(n+1)-n}{\sqrt{n+1}+\sqrt n}=\frac1{\sqrt{n+1}+\sqrt n}\to0,
\]
per la \eqref{eq:inf5} con $a_{n}=1$ e $b_{n}=\sqrt{n+1}+\sqrt n\to+\infty$ (per la
\eqref{eq:inf2}). È un esempio tipico di come una trasformazione algebrica possa
eliminare l'indeterminazione.
\end{esempio}

\begin{python}[Forme indeterminate: stessi limiti, comportamenti diversi]
# Prodotti "zero per infinito": un fattore tende a 0, l'altro
# diverge; il comportamento dipende dalle velocità relative.
casi = {
    "(1/n) * n^2":      lambda n: (1 / n) * n**2,
    "(1/n) * 3n":       lambda n: (1 / n) * 3 * n,
    "(1/n^2) * n":      lambda n: (1 / n**2) * n,
    "((-1)^n/n) * n":   lambda n: ((-1)**n / n) * n,
}
indici = (10, 11, 1000, 1001)
for nome, a in casi.items():
    valori = "  ".join(f"{a(n):>10.4g}" for n in indici)
    print(f"{nome:>16s}:  {valori}")
\end{python}
\begin{risultato}
     (1/n) * n^2:          10          11        1000        1001
      (1/n) * 3n:           3           3           3           3
     (1/n^2) * n:         0.1     0.09091       0.001    0.000999
  ((-1)^n/n) * n:           1          -1           1          -1
\end{risultato}

%% ----------------------------------------------------------------------------
%%  §3.5  Teoremi di confronto e altre proprietà dei limiti
%% ----------------------------------------------------------------------------
\section{Teoremi di confronto}
\label{sec:confronto}

Studiamo in questo paragrafo alcune relazioni tra limiti e ordinamento: che cosa si può
dedurre sul limite conoscendo il segno dei termini, e viceversa; e come si può
determinare il limite di una successione «intrappolandola» tra due successioni di cui
conosciamo il comportamento.

\subsection{Permanenza del segno}

\begin{teorema}[Permanenza del segno]\label{teo:permanenza-segno}\index{teorema!della permanenza del segno}
Se $\displaystyle\lim_{n\to+\infty}a_{n}=a>0$, esiste un numero $\nu$ tale che $a_{n}>0$
per ogni $n>\nu$.
\end{teorema}

\begin{schema}
\Ipotesi{$a_{n}\to a$, \ $a>0$.}
\Tesi{$a_{n}>0$ definitivamente.}
\end{schema}

\begin{strategia}
Se il limite è positivo, i termini della successione finiscono per stare in un intorno
di $a$; basta scegliere l'intorno abbastanza piccolo da non contenere lo zero. Il raggio
$\epsilon=a/2$ va bene: l'intorno $(a/2,\,3a/2)$ sta tutto a destra dello zero. Anzi,
otteniamo un'informazione più precisa: definitivamente $a_{n}>a/2$, cioè i termini non
solo sono positivi, ma si mantengono lontani da zero.
\end{strategia}

\begin{dimostrazione}
Dato che $a>0$, possiamo scegliere $\epsilon=a/2$. Esiste quindi un numero $\nu$ per cui
$\abs{a_{n}-a}<a/2$ per ogni $n>\nu$. Ciò equivale a $-a/2<a_{n}-a<a/2$. In particolare
abbiamo
\begin{equation}\label{eq:permanenza-passo}
  a_{n}>a-\frac a2=\frac a2>0 \qquad\perogni n>\nu. \qedhere
\end{equation}
\end{dimostrazione}

Sottolineiamo che, se una successione $a_{n}$ converge a un numero reale positivo $a$,
\emph{non} si può affermare in generale che \emph{tutti} i termini della successione
siano positivi. Ad esempio, in analogia con la \eqref{eq:succ-n1n}, la successione
$a_{n}=(n-7)/n$ converge al numero $1$, però i primi termini della successione (da $a_{1}$
fino ad $a_{6}$) sono negativi e $a_{7}=0$; il numero $\nu$ nella tesi del teorema della
permanenza del segno, in questo caso, può essere scelto uguale a $7$, e $7$ è il più
piccolo valore possibile.

\begin{osservazione}
Con la stessa dimostrazione (scegliendo $\epsilon=a-L$) si prova la versione più
generale: \emph{se $a_{n}\to a$ e $a>L$, allora definitivamente $a_{n}>L$}; e
analogamente, se $a<L$, definitivamente $a_{n}<L$. Il caso $L=0$ è il teorema.
\end{osservazione}

\begin{corollario}\label{cor:limite-non-negativo}
Se $\displaystyle\lim_{n\to+\infty}a_{n}=a$, e se $a_{n}\ge0$ per ogni $n$, allora anche
$a\ge0$.
\end{corollario}

\begin{schema}
\Ipotesi{$a_{n}\to a\in\R$; \ $a_{n}\ge0$ per ogni $n$.}
\Tesi{$a\ge0$.}
\end{schema}

\begin{strategia}
Ragioniamo per assurdo: se il limite fosse negativo, la permanenza del segno (applicata
a $-a_{n}$) renderebbe i termini definitivamente negativi, contro l'ipotesi.
\end{strategia}

\begin{dimostrazione}
Se per assurdo fosse $a<0$, il teorema della permanenza del segno, applicato alla
successione $-a_{n}$ (che converge a $-a>0$ per il \cref{teo:operazioni-limiti}),
comporterebbe che $-a_{n}>0$, cioè $a_{n}<0$, per $n$ grande, contro l'ipotesi.
\end{dimostrazione}

\begin{corollario}[Confronto dei limiti]\label{cor:confronto-limiti}
Se $\displaystyle\lim_{n\to+\infty}a_{n}=a$, $\displaystyle\lim_{n\to+\infty}b_{n}=b$, e
se $a_{n}\ge b_{n}$ per ogni $n$, allora $a\ge b$.
\end{corollario}

\begin{schema}
\Ipotesi{$a_{n}\to a\in\R$, \ $b_{n}\to b\in\R$; \ $a_{n}\ge b_{n}$ per ogni $n$.}
\Tesi{$a\ge b$.}
\end{schema}

\begin{strategia}
Confrontare due successioni equivale a confrontare la loro differenza con zero: basta
applicare il corollario precedente ad $a_{n}-b_{n}$, il cui limite è noto per il
\cref{teo:operazioni-limiti}.
\end{strategia}

\begin{dimostrazione}
Per ottenere la dimostrazione di quest'ultimo corollario basta applicare il corollario
precedente alla successione $a_{n}-b_{n}\ge0$, che converge ad $a-b$.
\end{dimostrazione}

Possiamo schematizzare i risultati ottenuti nel modo seguente (si noti la differenza tra
i segni $>$ e $\ge$):
\begin{align}
  a_{n}\to a,\ \ a>0\ \ &\implica\ \ a_{n}>0\ \text{ definitivamente}; \label{eq:schema-perm1}\\
  a_{n}\to a,\ \ a_{n}\ge0\ \ \perogni n\in\N\ \ &\implica\ \ a\ge0; \label{eq:schema-perm2}\\
  a_{n}\to a,\ \ b_{n}\to b,\ \ a_{n}\ge b_{n}\ \ \perogni n\in\N\ \ &\implica\ \ a\ge b. \label{eq:schema-perm3}
\end{align}

\begin{attenzione}[Le disuguaglianze strette non passano al limite]
Nelle \eqref{eq:schema-perm2} e \eqref{eq:schema-perm3} non si può sostituire $\ge$ con
$>$ nella tesi, nemmeno se l'ipotesi è stretta: $a_{n}=1/n>0$ per ogni $n$, ma il limite
è $0$, che non è positivo. Passando al limite, «$>$» si indebolisce in «$\ge$». Al
contrario, nella \eqref{eq:schema-perm1} l'ipotesi deve essere stretta: se fosse $a=0$, i
termini potrebbero avere qualsiasi segno, come accade per $(-1)^{n}/n$.
\end{attenzione}

\subsection{Il teorema dei carabinieri}

\begin{teorema}[dei carabinieri]\label{teo:carabinieri}\index{teorema!dei carabinieri}
Siano $a_{n}$, $b_{n}$, $c_{n}$ tre successioni tali che
\begin{equation}\label{eq:carabinieri-ipotesi}
  a_{n}\le c_{n}\le b_{n} \qquad\perogni n\in\N .
\end{equation}
Se $\displaystyle\lim_{n\to+\infty}a_{n}=\lim_{n\to+\infty}b_{n}=a$, allora anche la
successione $c_{n}$ è convergente e $\displaystyle\lim_{n\to+\infty}c_{n}=a$.
\end{teorema}

\begin{schema}
\Ipotesi{$a_{n}\le c_{n}\le b_{n}$ per ogni $n$ (basta definitivamente); \
$a_{n}\to a$, \ $b_{n}\to a$ (lo \emph{stesso} limite).}
\Tesi{$c_{n}\to a$ \ (in particolare, $c_{n}$ ammette limite).}
\end{schema}

\begin{strategia}
Si noti che la tesi contiene un'affermazione di \emph{esistenza}: non si suppone che
$c_{n}$ abbia limite, lo si dimostra. Non possiamo quindi usare il
\cref{cor:confronto-limiti}, che presuppone l'esistenza dei limiti; dobbiamo tornare alla
definizione. L'idea è visiva (\cref{fig:carabinieri}): fissato un intorno
$(a-\epsilon,a+\epsilon)$, definitivamente vi cadono sia $a_{n}$ sia $b_{n}$; e poiché
$c_{n}$ è «stretto» tra i due, ci cade anche lui. In formule, basta leggere le due
disuguaglianze $\abs{a_{n}-a}<\epsilon$ e $\abs{b_{n}-a}<\epsilon$ nella forma «doppia
disuguaglianza» e concatenarle con la \eqref{eq:carabinieri-ipotesi}.
\end{strategia}

\begin{dimostrazione}
Per ipotesi, per ogni $\epsilon>0$,
\begin{equation}\label{eq:carabinieri-passo1}
  \esiste\nu_{1}:\ \abs{a_{n}-a}<\epsilon\quad\perogni n>\nu_{1}; \qquad
  \esiste\nu_{2}:\ \abs{b_{n}-a}<\epsilon\quad\perogni n>\nu_{2}.
\end{equation}
Ricordiamo che le disuguaglianze con il valore assoluto si possono anche scrivere
\begin{equation}\label{eq:carabinieri-passo2}
  a-\epsilon<a_{n}<a+\epsilon; \qquad a-\epsilon<b_{n}<a+\epsilon .
\end{equation}
Quindi, se $n>\nu=\max\{\nu_{1},\nu_{2}\}$, risulta
\begin{equation}\label{eq:carabinieri-passo3}
  a-\epsilon<a_{n}\le c_{n}\le b_{n}<a+\epsilon .
\end{equation}
Perciò $\abs{c_{n}-a}<\epsilon$ per ogni $n>\nu$, come volevasi dimostrare.
\end{dimostrazione}

\begin{figure}[htb]
\centering
\begin{tikzpicture}
\begin{axis}[width=0.92\textwidth, height=5.2cm, xmin=0, xmax=31, ymin=-0.1, ymax=2.1,
  xtick={5,10,15,20,25,30}, ytick={1}, yticklabels={$a$}, xlabel={$n$}]
  \fill[verde!15] (axis cs:0,0.85) rectangle (axis cs:31,1.15);
  \addplot[tratteggio, domain=0:31] {1};
  \addplot[only marks, mark=triangle*, mark size=1.8pt, rame, samples at={1,...,30}] {1-1/x};
  \addplot[only marks, mark=triangle*, mark size=1.8pt, rame, mark options={rotate=180}, samples at={1,...,30}] {1+1/x};
  \addplot[only marks, mark=*, mark size=1.3pt, blu, samples at={1,...,30}] {1+0.8*sin(deg(1.3*x))/x};
  \node[font=\scriptsize, rame, anchor=west] at (axis cs:1.6,0.22) {$a_{n}$};
  \node[font=\scriptsize, rame, anchor=west] at (axis cs:1.6,1.93) {$b_{n}$};
  \node[font=\scriptsize, blu, anchor=west] at (axis cs:1.6,1.72) {$c_{n}$};
  \node[font=\scriptsize, verde!50!black, anchor=east] at (axis cs:30.5,1.25) {$(a-\epsilon,a+\epsilon)$};
\end{axis}
\end{tikzpicture}
\caption{Il teorema dei carabinieri: la successione $c_{n}$ (in blu) è compresa tra
$a_{n}$ e $b_{n}$ (i «carabinieri», in rame), che convergono allo stesso limite $a$.
Quando i due carabinieri sono entrati nella striscia verde $(a-\epsilon,a+\epsilon)$,
anche $c_{n}$, che sta in mezzo, è costretta a entrarvi.}
\label{fig:carabinieri}
\end{figure}

\begin{osservazione}
Nell'enunciato basta che la \eqref{eq:carabinieri-ipotesi} valga definitivamente, cioè
per $n$ maggiore di un certo $\nu_{3}$: nella dimostrazione si prende allora
$\nu=\max\{\nu_{1},\nu_{2},\nu_{3}\}$. In generale, il comportamento di una successione
per $n$ grande non dipende dai primi termini: modificando un numero finito di termini il
limite (se esiste) non cambia. Il nome del teorema, d'uso comune in Italia, evoca un
arrestato condotto tra due carabinieri: dove vanno loro, deve andare anche lui. In altri
paesi si chiama «teorema del sandwich» o «del confronto».
\end{osservazione}

\begin{idea}[Le inviluppanti di un'oscillazione smorzata]
Un esempio fisico: lo spostamento di una struttura che vibra con ampiezza decrescente,
campionato a istanti successivi, è una successione $c_{n}$ che oscilla attorno alla
posizione di equilibrio; ma se l'ampiezza $A_{n}$ tende a zero, le due curve
\emph{inviluppanti} $-A_{n}$ e $A_{n}$ fanno da carabinieri, e lo spostamento tende a
zero qualunque sia la forma dettagliata delle oscillazioni. È esattamente la situazione
della \cref{fig:carabinieri}, e del \cref{teo:limitata-infinitesima} qui sotto.
\end{idea}

Valgono analoghi teoremi di confronto anche per i limiti infiniti; qui basta un solo
«carabiniere».

\begin{teorema}[Confronto per limiti infiniti]\label{teo:confronto-infinito}
Se $a_{n}\le b_{n}$ per ogni $n\in\N$, allora
\begin{align}
  a_{n}\to+\infty &\ \implica\ b_{n}\to+\infty; \label{eq:confronto-piu-inf}\\
  b_{n}\to-\infty &\ \implica\ a_{n}\to-\infty. \label{eq:confronto-meno-inf}
\end{align}
\end{teorema}

\begin{schema}
\Ipotesi{$a_{n}\le b_{n}$ per ogni $n$ (basta definitivamente); \ $a_{n}\to+\infty$
(risp.\ $b_{n}\to-\infty$).}
\Tesi{$b_{n}\to+\infty$ (risp.\ $a_{n}\to-\infty$).}
\end{schema}

\begin{strategia}
Basta la definizione: se la successione più piccola supera definitivamente un livello
$M$, la più grande lo supera a maggior ragione, con lo stesso indice $\nu$. Non occorre
un secondo carabiniere, perché verso $+\infty$ non c'è nulla da «stringere».
\end{strategia}

\begin{dimostrazione}
Dimostriamo la \eqref{eq:confronto-piu-inf} (la prova della \eqref{eq:confronto-meno-inf}
è analoga). Per ipotesi $a_{n}\to+\infty$, che per la definizione
\eqref{eq:def-piu-inf} significa
\begin{equation}\label{eq:confronto-inf-passo}
  \perogni M>0\ \ \esiste\nu :\ a_{n}>M \quad\perogni n>\nu .
\end{equation}
Dato che $b_{n}\ge a_{n}$ per ogni $n\in\N$, si ottiene la tesi:
$b_{n}\ge a_{n}>M$ per ogni $n>\nu$.
\end{dimostrazione}

\subsection{Successioni infinitesime}

Riportiamo infine due ulteriori proprietà dei limiti di successioni; in particolare la
seconda, enunciata sotto forma di teorema, è importante per le applicazioni.

\begin{proposizione}\label{prop:infinitesima-modulo}
La successione $a_{n}$ converge a zero se e soltanto se $\abs{a_{n}}$ converge a zero.
\end{proposizione}

\begin{schema}
\Ipotesi{$a_{n}$ successione di numeri reali.}
\Tesi{$a_{n}\to0$ \ $\sse$ \ $\abs{a_{n}}\to0$.}
\end{schema}

\begin{strategia}
Quando il limite è zero, la definizione coinvolge soltanto la distanza da zero, cioè il
valore assoluto: scritte per esteso, le due condizioni sono la stessa frase, perché
$\bigl|\abs{a_{n}}\bigr|=\abs{a_{n}}$.
\end{strategia}

\begin{dimostrazione}
Posto $b_{n}=\abs{a_{n}}$, in base alla definizione \eqref{eq:def-limite} $b_{n}$
converge a zero se e solo se
\begin{equation}\label{eq:modulo-zero-passo}
  \perogni\epsilon>0\ \ \esiste\nu :\ \abs{b_{n}}<\epsilon \quad\perogni n>\nu .
\end{equation}
Dato che $\abs{b_{n}}=\bigl|\abs{a_{n}}\bigr|=\abs{a_{n}}$ per ogni $n\in\N$, la
\eqref{eq:modulo-zero-passo} è equivalente alla convergenza a zero della successione
$a_{n}$.
\end{dimostrazione}

Si noti che, nella proposizione precedente, è importante considerare non solo
successioni convergenti, ma più in particolare successioni \emph{convergenti a zero}. Ad
esempio, se $a_{n}=(-1)^{n}$, allora $a_{n}$ non è convergente, mentre
$b_{n}=\abs{a_{n}}$ è la successione costante $b_{n}=1$, che ovviamente converge a $1$. In
generale vale solo l'implicazione del \cref{lem:lim-modulo}: se $a_{n}\to a$ allora
$\abs{a_{n}}\to\abs a$.

\begin{teorema}[Prodotto di una limitata per un'infinitesima]\label{teo:limitata-infinitesima}\index{successione!infinitesima}
Se $a_{n}$ è una successione limitata e $b_{n}$ è una successione che converge a zero,
allora la successione prodotto $a_{n}\cdot b_{n}$ converge a zero.
\end{teorema}

\begin{schema}
\Ipotesi{$\abs{a_{n}}\le M$ per ogni $n$; \ $b_{n}\to0$.}
\Tesi{$a_{n}b_{n}\to0$.}
\end{schema}

\begin{strategia}
Si noti che $a_{n}$ \emph{non} è supposta convergente (potrebbe essere $(-1)^{n}$): il
\cref{teo:operazioni-limiti} non si applica. Ma $\abs{a_{n}b_{n}}\le M\abs{b_{n}}$, e la
successione a destra tende a zero: possiamo quindi usare i carabinieri $\pm M\abs{b_{n}}$
(primo metodo), oppure verificare direttamente la definizione nella forma
\eqref{eq:def-limite-c} (secondo metodo).
\end{strategia}

\begin{dimostrazione}[Primo metodo]
Per l'ipotesi di limitatezza si ha
\begin{equation}\label{eq:lim-inf-passo1}
  \abs{a_{n}b_{n}}=\abs{a_{n}}\cdot\abs{b_{n}}\le M\abs{b_{n}} \qquad\perogni n\in\N,
\end{equation}
che, per la proprietà \eqref{eq:mod-le} del valore assoluto, equivale a
\begin{equation}\label{eq:lim-inf-passo2}
  -M\abs{b_{n}}\le a_{n}b_{n}\le M\abs{b_{n}} \qquad\perogni n\in\N .
\end{equation}
Dato che per ipotesi $b_{n}\to0$, per la proposizione precedente anche la successione
$\abs{b_{n}}$ converge a zero, e quindi $\pm M\abs{b_{n}}\to0$. Per il teorema dei
carabinieri, dalla \eqref{eq:lim-inf-passo2} si deduce infine che $a_{n}b_{n}\to0$.
\end{dimostrazione}

\begin{dimostrazione}[Secondo metodo]
Per la definizione di limite si ha
$\perogni\epsilon>0\ \esiste\nu:\ \abs{b_{n}}<\epsilon$ per ogni $n>\nu$. Dall'ipotesi di
limitatezza si ottiene poi
\begin{equation}\label{eq:lim-inf-passo3}
  \abs{a_{n}b_{n}}=\abs{a_{n}}\cdot\abs{b_{n}}\le M\abs{b_{n}}<M\epsilon \qquad\perogni n>\nu,
\end{equation}
che equivale (si veda la \eqref{eq:def-limite-c}) al fatto che la successione prodotto
$a_{n}b_{n}$ converge a zero.
\end{dimostrazione}

\begin{esempio}
A titolo di esempio verifichiamo che
\begin{equation}\label{eq:es-limitata-inf}
  \lim_{n\to+\infty}\frac{(-1)^{n}(n+5)}{3n^{2}+1}=0;
\end{equation}
infatti si tratta del limite del prodotto della successione limitata $a_{n}=(-1)^{n}$
per la successione infinitesima
\begin{equation}\label{eq:es-limitata-inf-b}
  b_{n}=\frac{n+5}{3n^{2}+1}=\frac{1/n+5/n^{2}}{3+1/n^{2}}\to\frac03=0 .
\end{equation}
\end{esempio}

\begin{esempio}\label{es:senn-n}
Come ulteriore esempio verifichiamo che
\begin{equation}\label{eq:lim-senn-n}
  \lim_{n\to+\infty}\frac{\sen n}{n}=0 .
\end{equation}
Il limite è zero perché è il limite del prodotto della successione limitata
$a_{n}=\sen n$ ($\abs{a_{n}}=\abs{\sen n}\le1$ per ogni $n\in\N$, per la
\eqref{eq:sen-cos-limitate}) per la successione infinitesima $b_{n}=1/n$. Il lettore non
confonda questo limite con quello che dimostreremo nel \cref{teo:sen-x-su-x}, uguale a
$1$: là la successione $a_{n}$ converge a zero, mentre la successione in esame nella
\eqref{eq:lim-senn-n} è $\sen a_{n}/a_{n}$ con $a_{n}=n\to+\infty$.
\end{esempio}

%% ----------------------------------------------------------------------------
%%  §3.6  Alcuni limiti notevoli
%% ----------------------------------------------------------------------------
\section{Alcuni limiti notevoli}
\label{sec:limiti-notevoli}

In questo paragrafo esaminiamo alcuni esempi di limiti particolarmente importanti, che
useremo continuamente nel seguito come «mattoni» per calcolarne altri.

\subsection{La successione geometrica $a^{n}$}

\begin{teorema}[Limite della successione geometrica]\label{teo:lim-geometrica}\index{successione!geometrica}
Sia $a\in\R$ fissato. Allora
\begin{equation}\label{eq:lim-an}
  \lim_{n\to+\infty}a^{n}=
  \begin{cases}
    +\infty & \text{se } a>1,\\
    1 & \text{se } a=1,\\
    0 & \text{se } -1<a<1,\\
    \text{non esiste} & \text{se } a\le-1 .
  \end{cases}
\end{equation}
\end{teorema}

\begin{schema}
\Ipotesi{$a\in\R$ fissato.}
\Tesi{$a^{n}\to+\infty$ se $a>1$; \ $a^{n}\to1$ se $a=1$; \ $a^{n}\to0$ se
$\abs a<1$; \ $a^{n}$ non ha limite se $a\le-1$.}
\end{schema}

\begin{strategia}
Il caso chiave è $a>1$: bisogna mostrare che le potenze di un numero maggiore di $1$
crescono oltre ogni limite. Scrivendo $a=1+(a-1)$ con $a-1>0$, la disuguaglianza di
Bernoulli (\cref{teo:bernoulli}) fornisce una stima dal basso \emph{lineare},
$a^{n}\ge1+n(a-1)$, e il secondo membro diverge; il confronto (\cref{teo:confronto-infinito})
conclude. Tutti gli altri casi si riconducono a questo: per $0<\abs a<1$ si passa al
reciproco, $1/\abs a>1$; per $a\le-1$ si osserva che i termini cambiano segno a ogni
passo mentre i moduli non tendono a zero.
\end{strategia}

\begin{dimostrazione}
\emph{Caso $a>1$.} È possibile utilizzare la disuguaglianza di Bernoulli
\eqref{eq:bernoulli} con $x=a-1>0$:
\begin{equation}\label{eq:an-bernoulli}
  a^{n}=\bigl(1+(a-1)\bigr)^{n}\ge1+n(a-1).
\end{equation}
Dato che $a>1$, il secondo membro tende a $+\infty$ per $n\to+\infty$ (per la
\eqref{eq:inf3} e la \eqref{eq:inf1}); per il \cref{teo:confronto-infinito} anche
$a^{n}\to+\infty$.

\emph{Casi $a=1$ e $a=0$.} Sono ovvi: le successioni sono costanti, uguali a $1$ e a $0$.

\emph{Caso $0<\abs a<1$.} Risulta $1/\abs a>1$, e quindi dal caso già trattato
$(1/\abs a)^{n}\to+\infty$; per la \eqref{eq:inf5} otteniamo
\begin{equation}\label{eq:an-zero}
  \lim_{n\to+\infty}\abs{a^{n}}=\lim_{n\to+\infty}\abs a^{n}=\lim_{n\to+\infty}\frac{1}{(1/\abs a)^{n}}=0,
\end{equation}
e quindi, per la \cref{prop:infinitesima-modulo}, $a^{n}\to0$.

\emph{Caso $a=-1$.} Si riottiene la successione $(-1)^{n}$, che non ha limite finito
(\cref{prop:alterna-non-converge}) e non diverge, essendo limitata.

\emph{Caso $a<-1$.} Si ha $\abs{a^{n}}=\abs a^{n}\to+\infty$, perché $\abs a>1$; dunque la
successione non è limitata, e per il \cref{teo:convergente-limitata} non può convergere.
Non può nemmeno divergere a $+\infty$, perché in tal caso sarebbe definitivamente
positiva, mentre i termini di indice dispari $a^{2k+1}=a\cdot(a^{2})^{k}$ sono negativi;
né a $-\infty$, perché i termini di indice pari $a^{2k}=(a^{2})^{k}$ sono positivi. Si noti
che invece esiste, ed è uguale a $+\infty$, il limite di $\abs{a^{n}}$ se $a<-1$.
\end{dimostrazione}

\begin{marginfigure}
\centering
\begin{tikzpicture}
\begin{axis}[width=5.4cm, height=4.8cm, xmin=0, xmax=12.5, ymin=-2.2, ymax=2.6,
  xtick={5,10}, ytick={-1,1}]
  \addplot[only marks, mark=*, mark size=1.1pt, blu, samples at={0,...,12}] {1.08^x};
  \addplot[only marks, mark=square*, mark size=1.1pt, verde!60!black, samples at={0,...,12}] {0.75^x};
  \addplot[only marks, mark=triangle*, mark size=1.4pt, rame, samples at={0,...,12}] {(-0.8)^x};
  \node[font=\tiny, blu] at (axis cs:5,1.95) {$1{,}08^{n}$};
  \node[font=\tiny, verde!60!black] at (axis cs:9.5,0.35) {$0{,}75^{n}$};
  \node[font=\tiny, rame] at (axis cs:3.5,-1.4) {$(-0{,}8)^{n}$};
\end{axis}
\end{tikzpicture}
\caption{Successioni geometriche: per $a>1$ divergono (anche se lentamente, come
$1{,}08^{n}$), per $\abs a<1$ tendono a zero, eventualmente oscillando se $a<0$.}
\end{marginfigure}

\subsection{Le radici $\sqrt[n]{a}$ e $\sqrt[n]{n}$}

\begin{proposizione}\label{prop:radice-n-a}
Se $a$ è un numero reale positivo, risulta
\begin{equation}\label{eq:lim-radice-a}
  \lim_{n\to+\infty}\sqrt[n]{a}=\lim_{n\to+\infty}a^{1/n}=1 .
\end{equation}
\end{proposizione}

\begin{schema}
\Ipotesi{$a>0$.}
\Tesi{$\sqrt[n]a\to1$.}
\end{schema}

\begin{strategia}
Nel caso $a\ge1$ la radice è $\ge1$, e scriviamo $\sqrt[n]a=1+b_{n}$ con $b_{n}\ge0$: la
tesi diventa $b_{n}\to0$. Elevando alla $n$ e usando ancora Bernoulli,
$a=(1+b_{n})^{n}\ge1+nb_{n}$, cioè $b_{n}\le\frac{a-1}{n}$: la successione $b_{n}$ è
intrappolata tra $0$ e una successione infinitesima. È lo stesso trucco del teorema
precedente, usato «al contrario»: là Bernoulli serviva a mostrare che $a^{n}$ cresce
almeno linearmente, qui a mostrare che la radice $n$-esima si avvicina a $1$ almeno come
$1/n$. Il caso $0<a<1$ si riduce al precedente passando al reciproco.
\end{strategia}

\begin{dimostrazione}
Nel caso $a\ge1$ si ha $\sqrt[n]a\ge1$ (se fosse $\sqrt[n]a<1$, elevando alla $n$ sarebbe
$a<1$). Poniamo $b_{n}=\sqrt[n]a-1$; risulta $b_{n}\ge0$ e inoltre, sempre per la
disuguaglianza di Bernoulli,
\begin{equation}\label{eq:radice-a-passo}
  a=(1+b_{n})^{n}\ge1+n\,b_{n}.
\end{equation}
Perciò
\begin{equation}\label{eq:radice-a-stima}
  0\le b_{n}\le\frac{a-1}{n}.
\end{equation}
Dal teorema dei carabinieri segue che $b_{n}\to0$, cioè che $\sqrt[n]a\to1$. Se $0<a<1$,
allora $1/a>1$ e quindi, per quanto già dimostrato e per la \eqref{eq:lim-quoziente},
\begin{equation}\label{eq:radice-a-reciproco}
  \sqrt[n]a=\frac{1}{\sqrt[n]{1/a}}\to\frac11=1. \qedhere
\end{equation}
\end{dimostrazione}

Il limite \eqref{eq:lim-radice-a} si ricorda facilmente per mezzo dei passaggi
$\sqrt[n]{a}=a^{1/n}\to a^{0}=1$; questi passaggi, però, non sono una dimostrazione,
perché presuppongono la «continuità» della funzione esponenziale, che dimostreremo solo
più avanti.

Più sorprendente è il fatto che anche la radice $n$-esima di $n$ tenda a $1$
(\cref{fig:radici-n}): qui la base cresce, ma l'indice della radice cresce più in fretta. Nel \cref{sec:binomio} avevamo
anticipato che per dimostrarlo ci sarebbe servita la stima (iii) del
\cref{cor:binomio}; manteniamo la promessa.

\begin{proposizione}\label{prop:radice-n-n}
Risulta
\begin{equation}\label{eq:lim-radice-n}
  \lim_{n\to+\infty}\sqrt[n]{n}=1 .
\end{equation}
\end{proposizione}

\begin{schema}
\Ipotesi{nessuna ipotesi aggiuntiva: $a_{n}=\sqrt[n]n$, \ $n\in\N$.}
\Tesi{$\sqrt[n]n\to1$.}
\end{schema}

\begin{strategia}
Ripetiamo l'idea della proposizione precedente: $\sqrt[n]n=1+b_{n}$ con $b_{n}\ge0$, e
$n=(1+b_{n})^{n}$. Ma ora la stima di Bernoulli, $n\ge1+nb_{n}$, dà soltanto
$b_{n}\le\frac{n-1}{n}$, che \emph{non} tende a zero: Bernoulli è troppo debole. Occorre
una stima dal basso di $(1+b_{n})^{n}$ che cresca più velocemente di $n$, e la formula del
binomio la fornisce: tenendo il termine di secondo grado,
$(1+b_{n})^{n}\ge\frac{n(n-1)}{2}b_{n}^{2}$, che contiene $n^{2}$. Da qui
$b_{n}^{2}\le\frac2{n-1}$.
\end{strategia}

\begin{dimostrazione}
Per $n\ge1$ si ha $\sqrt[n]n\ge1$; poniamo $b_{n}=\sqrt[n]n-1\ge0$. Per $n\ge2$, per la
\cref{cor:binomio}(iii) con $x=b_{n}\ge0$,
\begin{equation}\label{eq:radice-n-binomio}
  n=(1+b_{n})^{n}\ge1+nb_{n}+\frac{n(n-1)}{2}b_{n}^{2}\ge\frac{n(n-1)}{2}b_{n}^{2},
\end{equation}
da cui $b_{n}^{2}\le\frac{2}{n-1}$, cioè
\begin{equation}\label{eq:radice-n-stima}
  0\le b_{n}\le\sqrt{\frac{2}{n-1}} \qquad (n\ge2).
\end{equation}
Il secondo membro tende a zero: fissato $\epsilon>0$, si ha
$\sqrt{2/(n-1)}<\epsilon$ se e solo se $n>1+2/\epsilon^{2}$. Per il teorema dei
carabinieri $b_{n}\to0$, cioè $\sqrt[n]n\to1$.
\end{dimostrazione}

\begin{corollario}\label{cor:radice-n-nb}
Per ogni $b\in\R$ risulta
\begin{equation}\label{eq:lim-radice-nb}
  \lim_{n\to+\infty}\sqrt[n]{n^{b}}=\lim_{n\to+\infty}n^{b/n}=1 .
\end{equation}
\end{corollario}

\begin{schema}
\Ipotesi{$b\in\R$.}
\Tesi{$n^{b/n}\to1$.}
\end{schema}

\begin{strategia}
Se $b=k$ è intero, $n^{k/n}=\bigl(\sqrt[n]n\bigr)^{k}$ è una potenza fissa di una
successione che tende a $1$, e basta il teorema sul prodotto. Per $b$ reale qualsiasi
intrappoliamo l'esponente tra due interi consecutivi, $\lfloor b\rfloor\le b<\lfloor
b\rfloor+1$, e usiamo la monotonia di $x\mapsto n^{x}$ e il teorema dei carabinieri.
\end{strategia}

\begin{dimostrazione}
Consideriamo dapprima il caso in cui l'esponente $b$ sia un intero $k\in\Z$. In tal caso
$n^{k/n}=\bigl(n^{1/n}\bigr)^{k}$ è il prodotto di $\abs k$ fattori uguali a $n^{1/n}$ (o
il reciproco di tale prodotto, se $k<0$; per $k=0$ vale costantemente $1$), e per la
\cref{prop:radice-n-n} e il \cref{teo:operazioni-limiti} tende a $1^{k}=1$.

Per trattare il caso generale $b\in\R$ ricorriamo alla parte intera $\lfloor b\rfloor$
(\cref{prop:parte-intera}): $\lfloor b\rfloor\le b<\lfloor b\rfloor+1$. Poiché per $n\ge2$
la funzione $x\mapsto n^{x}$ è crescente (\cref{teo:esponenziale}), e per $n=1$ è
costante, si ha
\begin{equation}\label{eq:radice-nb-stima}
  n^{\lfloor b\rfloor/n}\le n^{b/n}\le n^{(\lfloor b\rfloor+1)/n} \qquad\perogni n\in\N .
\end{equation}
Il primo e il terzo membro tendono a $1$ per il caso intero appena visto; per il teorema
dei carabinieri si ottiene la tesi.
\end{dimostrazione}

\begin{osservazione}[Un'altra strada]
Il testo di riferimento \cite{marcellini-sbordone} dimostra la \eqref{eq:lim-radice-nb}
partendo dal caso $b=\frac12$, per il quale basta Bernoulli: posto $b_{n}=n^{1/(2n)}-1$,
si ha $\sqrt n=(1+b_{n})^{n}\ge1+nb_{n}$, da cui
$0\le b_{n}\le\frac{\sqrt n-1}{n}<\frac1{\sqrt n}\to0$. Si ottiene poi il caso intero
come $n^{k/n}=(n^{1/(2n)})^{2k}$. Il lettore si chieda perché si è scelto proprio
$b=\frac12$, perché la dimostrazione non funziona con $b=1$, e per quali altri valori di
$b$ funziona (si veda l'\cref{es:quesiti-b}).
\end{osservazione}

\begin{figure}[htb]
\centering
\begin{tikzpicture}
\begin{axis}[width=0.92\textwidth, height=5cm, xmin=0, xmax=61, ymin=0.9, ymax=1.75,
  xtick={10,20,30,40,50,60}, ytick={1,1.2,1.4,1.6},
  yticklabels={$1$,$1{,}2$,$1{,}4$,$1{,}6$}, xlabel={$n$},
  legend style={font=\scriptsize, draw=none, fill=none, at={(0.98,0.97)}, anchor=north east}]
  \addplot[only marks, mark=*, mark size=1.1pt, blu, samples at={1,...,60}] {x^(1/x)};
  \addplot[only marks, mark=square*, mark size=1pt, rame, samples at={3,...,60}] {5^(1/x)};
  \addplot[tratteggio, domain=0:61] {1};
  \legend{$\sqrt[n]{n}$, $\sqrt[n]{5}$}
\end{axis}
\end{tikzpicture}
\caption{Le successioni $\sqrt[n]5$ (decrescente) e $\sqrt[n]n$: entrambe tendono a $1$.
La seconda prima cresce (fino a $n=3$, dove vale $\sqrt[3]3\approx1{,}442$) e poi decresce
lentamente verso $1$. Per $\sqrt[n]5$ sono riportati i termini da $n=3$ in poi
($\sqrt[1]5=5$ e $\sqrt5\approx2{,}236$ uscirebbero dal grafico).}
\label{fig:radici-n}
\end{figure}

\subsection{Limiti trigonometrici}

Esaminiamo ora tre limiti relativi alle funzioni trigonometriche. Ci serve una
disuguaglianza elementare.

\begin{lemma}\label{lem:sen-minore-x}
Per ogni $x\in\R$ si ha $\abs{\sen x}\le\abs x$, con l'uguaglianza solo per $x=0$.
\end{lemma}

\begin{schema}
\Ipotesi{$x\in\R$.}
\Tesi{$\abs{\sen x}\le\abs x$, \ e $\abs{\sen x}=\abs x$ solo per $x=0$.}
\end{schema}

\begin{strategia}
Per gli angoli piccoli e positivi la disuguaglianza è la \eqref{eq:sen-x-tg}, che viene
da un confronto geometrico nel cerchio unitario; gli angoli negativi si trattano con la disparità del
seno; per gli angoli grandi, $\abs x\ge\pi/2$, basta che il seno non superi $1$.
\end{strategia}

\begin{dimostrazione}
Se $0<x<\pi/2$, per la \eqref{eq:sen-x-tg} si ha $0<\sen x<x$. Se $-\pi/2<x<0$, poiché il
seno è dispari, $\abs{\sen x}=\abs{\sen(-x)}=\sen(-x)<-x=\abs x$. Se $\abs x\ge\pi/2$,
allora $\abs{\sen x}\le1<\pi/2\le\abs x$ (si ricordi che $\pi>3$). Infine per $x=0$ vale
l'uguaglianza.
\end{dimostrazione}

\begin{teorema}[Seno e coseno di una successione infinitesima]\label{teo:lim-sen-cos}
Se $a_{n}\to0$, allora
\begin{equation}\label{eq:lim-sen-cos}
  \sen a_{n}\to0 \qquad\text{e}\qquad \cos a_{n}\to1 .
\end{equation}
\end{teorema}

\begin{schema}
\Ipotesi{$a_{n}\to0$.}
\Tesi{$\sen a_{n}\to0$, \ $\cos a_{n}\to1$.}
\end{schema}

\begin{strategia}
Per il seno, il \cref{lem:sen-minore-x} dà $\abs{\sen a_{n}}\le\abs{a_{n}}$: il seno è
«schiacciato» verso zero insieme ad $a_{n}$ (carabinieri). Per il coseno, la formula
$\cos x=1-2\sen^{2}\frac x2$ riconduce tutto al seno, già trattato, e alle operazioni con
i limiti.
\end{strategia}

\begin{dimostrazione}
Per il \cref{lem:sen-minore-x}, $0\le\abs{\sen a_{n}}\le\abs{a_{n}}$. Per la
\cref{prop:infinitesima-modulo} $\abs{a_{n}}\to0$; per il teorema dei carabinieri
$\abs{\sen a_{n}}\to0$, e ancora per la \cref{prop:infinitesima-modulo} $\sen a_{n}\to0$.

Per il coseno usiamo la formula di duplicazione \eqref{eq:duplicazione} e la relazione
fondamentale \eqref{eq:relazione-fondamentale}: per ogni $x$,
$\cos x=\cos^{2}\frac x2-\sen^{2}\frac x2=1-2\sen^{2}\frac x2$. Poiché $a_{n}/2\to0$, per
la prima parte $\sen(a_{n}/2)\to0$, e quindi, per il \cref{teo:operazioni-limiti},
\begin{equation}\label{eq:cos-passo}
  \cos a_{n}=1-2\sen^{2}\frac{a_{n}}{2}\to1-2\cdot0=1. \qedhere
\end{equation}
\end{dimostrazione}

Ad esempio $\sen(1/n)\to0$ e $\cos(1/n)\to1$, perché $1/n\to0$.

\begin{osservazione}
Il testo \cite{marcellini-sbordone} ottiene il limite del coseno dalla relazione
$\cos a_{n}=\sqrt{1-\sen^{2}a_{n}}$, valida definitivamente perché $\cos a_{n}\ge0$ per
$\abs{a_{n}}\le\pi/2$; ciò richiede di sapere che $b_{n}\to1$ implica $\sqrt{b_{n}}\to1$.
La formula di duplicazione evita questo passaggio.
\end{osservazione}

Il terzo limite di funzione trigonometrica che prendiamo in considerazione è il più
importante, ed è una forma indeterminata $0/0$.

\begin{figure}[htb]
\centering
\begin{tikzpicture}
\begin{axis}[width=0.92\textwidth, height=5.2cm, xmin=-1.75, xmax=1.75, ymin=-0.05, ymax=1.18,
  xtick={-1.5708,-0.7854,0.7854,1.5708}, xticklabels={$-\frac\pi2$,$-\frac\pi4$,$\frac\pi4$,$\frac\pi2$},
  ytick={0.5}, yticklabels={$0{,}5$}, legend style={font=\scriptsize, draw=none, fill=none, at={(0.5,0.32)}, anchor=north}, legend columns=3]
  \addplot[curva2, domain=-1.5708:1.5708, samples=120] {cos(deg(x))};
  \addplot[curva, domain=-1.5708:-0.001, samples=120, forget plot] {sin(deg(x))/x};
  \addplot[curva, domain=0.001:1.5708, samples=120] {sin(deg(x))/x};
  \addplot[curva3, domain=-1.5708:1.5708] {1};
  \addplot[only marks, mark=o, mark size=1.8pt, blu, fill=white] coordinates {(0,1)};
  \legend{$\cos x$, $\frac{\sen x}{x}$, $1$}
\end{axis}
\end{tikzpicture}
\caption{La disuguaglianza \eqref{eq:cos-sen-x-x}: per $0<\abs x<\pi/2$ il grafico di
$\frac{\sen x}{x}$ è compreso tra quello di $\cos x$ e la retta $y=1$. Avvicinandosi a
$x=0$, dove la funzione non è definita (cerchietto vuoto), i due «carabinieri» si
stringono entrambi verso il valore $1$.}
\label{fig:sen-x-su-x}
\end{figure}

\begin{teorema}[Il limite fondamentale $\sen x/x$]\label{teo:sen-x-su-x}\index{limite!notevole $\sen x/x$}
Se $a_{n}\to0$ e $a_{n}\neq0$ per ogni $n$, allora
\begin{equation}\label{eq:lim-sen-x-x}
  \lim_{n\to+\infty}\frac{\sen a_{n}}{a_{n}}=1 .
\end{equation}
\end{teorema}

\begin{schema}
\Ipotesi{$a_{n}\to0$, \ $a_{n}\neq0$ per ogni $n$.}
\Tesi{$\dfrac{\sen a_{n}}{a_{n}}\to1$.}
\end{schema}

\begin{strategia}
Il quoziente è una forma $0/0$, quindi le regole sulle operazioni non servono. Usiamo
il teorema dei carabinieri: cerchiamo due successioni, entrambe convergenti a $1$, tra
cui $\frac{\sen a_{n}}{a_{n}}$ sia compresa. Le fornisce la geometria, con la catena di
disuguaglianze $\sen x<x<\tg x$ per $0<x<\pi/2$ \eqref{eq:sen-x-tg}: dividendo per
$\sen x$ e passando ai reciproci si ottiene
$\cos x<\frac{\sen x}{x}<1$ (\cref{fig:sen-x-su-x}). Il carabiniere inferiore $\cos a_{n}$ tende a $1$ per il
teorema precedente; quello superiore è costante. Occorre solo un po' di attenzione per i
valori negativi di $x$, che si trattano con la parità delle funzioni coinvolte.
\end{strategia}

\begin{dimostrazione}
Cominciamo col dimostrare che
\begin{equation}\label{eq:cos-sen-x-x}
  0<\abs x<\frac\pi2 \quad\implica\quad \cos x<\frac{\sen x}{x}<1 .
\end{equation}
Infatti, se $x$ è positivo, dalla \eqref{eq:sen-x-tg} si ottiene
\begin{equation}\label{eq:sen-x-tg-bis}
  \sen x<x<\tg x=\frac{\sen x}{\cos x},
\end{equation}
da cui, dividendo la prima disuguaglianza per $x>0$ e moltiplicando la seconda per
$\frac{\cos x}{x}>0$ (si ricordi che $\cos x>0$ per $0<x<\frac\pi2$), e tenendo presente
che moltiplicare per un numero positivo conserva le disuguaglianze strette,
\begin{equation}\label{eq:cos-sen-passo}
  \cos x<\frac{\sen x}{x}<1 .
\end{equation}
Se invece $x$ è negativo, applicando a $-x$ la \eqref{eq:cos-sen-passo} e ricordando che il
coseno è pari e il seno dispari, si ottiene
\begin{equation}\label{eq:cos-sen-negativo}
  \cos x=\cos(-x)<\frac{\sen(-x)}{-x}=\frac{\sen x}{x}<1,
\end{equation}
dunque la \eqref{eq:cos-sen-passo} vale anche se $x$ è negativo, e la \eqref{eq:cos-sen-x-x}
è dimostrata.

Dato che $a_{n}\to0$, per la definizione di limite esiste un indice $\nu$ tale che
$\abs{a_{n}}<\pi/2$ per ogni $n>\nu$. Dalla \eqref{eq:cos-sen-x-x} si ottiene
\begin{equation}\label{eq:sen-x-x-carabinieri}
  \cos a_{n}<\frac{\sen a_{n}}{a_{n}}<1 \qquad\perogni n>\nu .
\end{equation}
Per $n\to+\infty$, $\cos a_{n}\to1$ (\cref{teo:lim-sen-cos}). La tesi
\eqref{eq:lim-sen-x-x} discende quindi dal teorema dei carabinieri.
\end{dimostrazione}

\begin{idea}[L'approssimazione dei piccoli angoli]
Il \cref{teo:sen-x-su-x} giustifica l'approssimazione $\sen\alpha\approx\alpha$ (e, di
conseguenza, $\tg\alpha\approx\alpha$, perché $\tg\alpha/\alpha=\frac{\sen\alpha}{\alpha}\cdot
\frac1{\cos\alpha}\to1$) per angoli piccoli \emph{espressi in radianti}, usata
ovunque in meccanica del volo. Nelle equazioni del volo librato o in salita, con
angolo di rampa $\gamma$ piccolo, si scrive $W\sen\gamma\approx W\gamma$ e
$W\cos\gamma\approx W$; la velocità ascensionale $V\sen\gamma$ diventa $V\gamma$; nella
teoria dei profili sottili la linearizzazione delle condizioni al contorno per piccole
incidenze porta alla celebre $C_{L}=2\pi\alpha$. L'errore relativo è modesto: per
$\alpha=5^{\circ}=0{,}0873$~rad si ha $\frac{\sen\alpha}{\alpha}=0{,}99873$, cioè un errore
dello $0{,}13\%$; a $15^{\circ}$ l'errore sale all'$1{,}1\%$ (il programma che segue
calcola alcuni valori). La stima \eqref{eq:cos-sen-x-x} fornisce anche un controllo
dell'errore: $1-\frac{\sen\alpha}{\alpha}<1-\cos\alpha=2\sen^{2}\frac\alpha2\le\frac{\alpha^{2}}2$.
\end{idea}

\begin{python}[Quanto è buona l'approssimazione $\sen\alpha\approx\alpha$?]
import math

print(" gradi   radianti   sen(a)/a   errore %   stima a^2/2 %")
for gradi in (1, 2, 5, 10, 15, 20, 30):
    a = math.radians(gradi)
    r = math.sin(a) / a
    print(f"{gradi:5d}   {a:8.5f}   {r:8.6f}   {100*(1-r):7.4f}"
          f"    {100*a*a/2:8.4f}")
\end{python}
\begin{risultato}
 gradi   radianti   sen(a)/a   errore %   stima a^2/2 %
    1    0.01745   0.999949    0.0051      0.0152
    2    0.03491   0.999797    0.0203      0.0609
    5    0.08727   0.998731    0.1269      0.3808
   10    0.17453   0.994931    0.5069      1.5231
   15    0.26180   0.988616    1.1384      3.4269
   20    0.34907   0.979816    2.0184      6.0923
   30    0.52360   0.954930    4.5070     13.7078
\end{risultato}

\subsection{L'area del cerchio}

Applichiamo il risultato appena dimostrato per concludere l'argomento del
\cref{sec:premessa-limiti}.

\begin{teorema}[Area del cerchio]\label{teo:area-cerchio}
L'area del cerchio di raggio $1$ è uguale a $\pi$.
\end{teorema}

\begin{schema}
\Ipotesi{$\pi$ è la lunghezza della semicirconferenza di raggio $1$; l'area del cerchio
è, per definizione, il limite delle aree $a_{n}$ dei poligoni regolari inscritti.}
\Tesi{$\displaystyle\lim_{n\to+\infty}a_{n}=\pi$.}
\end{schema}

\begin{strategia}
La formula \eqref{eq:area-poligono}, $a_{n}=\frac n2\sen\frac{2\pi}n$, è una forma
$\infty\cdot0$. Basta riscriverla mettendo in evidenza il quoziente
$\frac{\sen x_{n}}{x_{n}}$ con $x_{n}=\frac{2\pi}n\to0$: il fattore che resta è proprio
$\pi$, e il limite fondamentale appena dimostrato conclude.
\end{strategia}

\begin{dimostrazione}
In base alla \eqref{eq:area-poligono} si ottiene
\begin{equation}\label{eq:area-cerchio-1}
  \text{area del cerchio di raggio }1=\lim_{n\to+\infty}\frac n2\sen\frac{2\pi}{n}.
\end{equation}
Dato che, per $n\to+\infty$, la successione $2\pi/n$ tende a zero ed è sempre diversa da
zero, siamo nella situazione del limite notevole \eqref{eq:lim-sen-x-x}, e quindi
\begin{equation}\label{eq:area-cerchio-2}
  \text{area del cerchio di raggio }1=\lim_{n\to+\infty}\pi\,\frac{\sen(2\pi/n)}{2\pi/n}=\pi\cdot1=\pi. \qedhere
\end{equation}
\end{dimostrazione}

\begin{osservazione}[Coerenza della definizione di $\pi$]
Lo stesso ragionamento applicato ai perimetri dà un controllo di coerenza: il perimetro
di $P_{n}$ è $p_{n}=2n\sen\frac\pi n=2\pi\,\frac{\sen(\pi/n)}{\pi/n}\to2\pi$, cioè la
lunghezza della circonferenza, come deve essere se $\pi$ è la lunghezza della
semicirconferenza. Si osservi però che il ragionamento poggia su due definizioni
(lunghezza della circonferenza e area del cerchio) e sulla disuguaglianza
\eqref{eq:sen-x-tg}, la cui giustificazione nel \cref{sec:trigonometriche} era di natura
geometrica (confronto tra corda, arco e aree). Una teoria pienamente rigorosa della
lunghezza delle curve e dell'area delle figure piane sarà sviluppata con il calcolo
integrale; per ora, il \cref{teo:area-cerchio} va letto come la prova che le due
nozioni geometriche di $\pi$ --- come rapporto tra circonferenza e diametro e come area
del cerchio unitario --- coincidono.
\end{osservazione}

%% ----------------------------------------------------------------------------
%%  §3.7  Successioni monotòne
%% ----------------------------------------------------------------------------
\section{Successioni monotòne}
\label{sec:monotone}

Abbiamo già dato nel \cref{sec:invertibili} la definizione di funzioni (strettamente)
crescenti o decrescenti. Allo stesso modo, per le successioni diciamo che:

\begin{definizione}[Successioni monotòne]\index{successione!monotòna}
Una successione $a_{n}$ si dice
\begin{align}
  &\text{\emph{strettamente crescente}} \text{ se } & a_{n}&<a_{n+1} & &\perogni n\in\N; \label{eq:succ-str-cresc}\\
  &\text{\emph{crescente}} \text{ se } & a_{n}&\le a_{n+1} & &\perogni n\in\N; \label{eq:succ-cresc}\\
  &\text{\emph{strettamente decrescente}} \text{ se } & a_{n}&>a_{n+1} & &\perogni n\in\N; \label{eq:succ-str-decr}\\
  &\text{\emph{decrescente}} \text{ se } & a_{n}&\ge a_{n+1} & &\perogni n\in\N. \label{eq:succ-decr}
\end{align}
Una successione $a_{n}$ è \emph{monotòna} se verifica una delle quattro condizioni sopra
scritte; è \emph{strettamente monotòna} se verifica la \eqref{eq:succ-str-cresc} oppure la
\eqref{eq:succ-str-decr}.
\end{definizione}

Basta confrontare ogni termine con il successivo: per induzione, se $a_{n}\le a_{n+1}$
per ogni $n$, allora $a_{n}\le a_{m}$ per ogni $n\le m$, e la successione è crescente
anche nel senso della definizione data per le funzioni. Se $a_{n}=a$ per ogni $n\in\N$,
con $a$ numero reale fissato, si dice che $a_{n}$ è una successione \emph{costante}. Le
successioni costanti sono allo stesso tempo crescenti e decrescenti.

\begin{esempio}
Le successioni \eqref{eq:succ-1n} e \eqref{eq:succ-n1n} sono strettamente monotòne: la
\eqref{eq:succ-1n} $a_{n}=1/n$ è strettamente decrescente, mentre la \eqref{eq:succ-n1n}
$a_{n}=(n-1)/n$ è strettamente crescente; infatti, se $a_{n}=1/n$, allora
\begin{equation}\label{eq:es-monotona-1}
  a_{n}>a_{n+1} \quad\sse\quad \frac1n>\frac1{n+1} \quad\sse\quad n+1>n,
\end{equation}
e l'ultima disuguaglianza è verificata per ogni $n\in\N$. Invece, se $a_{n}=(n-1)/n$,
\begin{equation}\label{eq:es-monotona-2}
  a_{n}<a_{n+1} \quad\sse\quad \frac{n-1}{n}<\frac{n}{n+1} \quad\sse\quad n^{2}-1<n^{2},
\end{equation}
e anche questa volta l'ultima disuguaglianza è verificata per ogni $n\in\N$.
\end{esempio}

Il risultato seguente è di fondamentale importanza: nello schema della
\cref{fig:schema-teoremi} è il primo dei grandi teoremi di esistenza che seguono
l'esistenza dell'estremo superiore, e da essa discende direttamente.

\begin{teorema}[sulle successioni monotòne]\label{teo:monotone}\index{teorema!sulle successioni monotòne}
Ogni successione monotòna ammette limite. In particolare, ogni successione monotòna e
limitata è convergente, cioè ammette limite finito. Più precisamente:
\begin{equation}\label{eq:lim-monotone}
\begin{aligned}
  a_{n}\ \text{crescente} &\implica \lim_{n\to+\infty}a_{n}=\sup_{n\in\N}a_{n};\\
  a_{n}\ \text{decrescente} &\implica \lim_{n\to+\infty}a_{n}=\inf_{n\in\N}a_{n}.
\end{aligned}
\end{equation}
\end{teorema}

\begin{schema}
\Ipotesi{$a_{n}$ monotòna (ad esempio crescente).}
\Tesi{$a_{n}$ è regolare, e $\lim a_{n}=\sup_{n}a_{n}$ (che è finito se e solo se $a_{n}$
è limitata superiormente, e vale $+\infty$ altrimenti).}
\end{schema}

\begin{strategia}
Il candidato limite di una successione crescente è il «tetto» dei suoi valori, cioè
$\ell=\sup_{n}a_{n}$. Per dimostrare che è il limite servono due cose, e ciascuna viene da
uno dei due ingredienti:
\begin{itemize}
  \item i termini non superano mai $\ell$, perché $\ell$ è un maggiorante; quindi
  $a_{n}<\ell+\epsilon$ per ogni $n$, gratis;
  \item i termini si avvicinano a $\ell$ \emph{da sotto}: per la caratterizzazione
  dell'estremo superiore (\cref{prop:caratt-sup}) c'è \emph{almeno un} termine $a_{\nu}$
  maggiore di $\ell-\epsilon$; ed essendo la successione crescente, \emph{tutti} i
  termini successivi lo sono a maggior ragione.
\end{itemize}
È la monotonia che trasforma «almeno un termine» in «definitivamente tutti»
(\cref{fig:monotona-sup}): senza di essa la successione potrebbe avvicinarsi a $\ell$ e poi allontanarsene. Se la successione
non è limitata superiormente, lo stesso ragionamento con un livello $M$ al posto di
$\ell-\epsilon$ dà la divergenza a $+\infty$. Le successioni decrescenti si trattano in
modo simmetrico, con l'estremo inferiore.
\end{strategia}

\begin{dimostrazione}
Consideriamo il caso di una successione $a_{n}$ crescente e limitata superiormente.
Posto $\ell=\sup_{n}a_{n}$, che esiste ed è finito per il \cref{teo:esistenza-sup},
fissato $\epsilon>0$, per le proprietà dell'estremo superiore
(\cref{prop:caratt-sup}) esiste $\nu\in\N$ tale che
\begin{equation}\label{eq:monotone-passo1}
  \ell-\epsilon<a_{\nu}.
\end{equation}
Per $n>\nu$ risulta $a_{\nu}\le a_{n}$, e d'altra parte $a_{n}\le\ell$ perché $\ell$ è un
maggiorante; dunque
\begin{equation}\label{eq:monotone-passo2}
  \ell-\epsilon<a_{\nu}\le a_{n}\le\ell<\ell+\epsilon \qquad\perogni n>\nu,
\end{equation}
da cui $\lim_{n\to+\infty}a_{n}=\ell$.

Consideriamo ora il caso di una successione $a_{n}$ crescente e non limitata
(superiormente). Fissato $M>0$, poiché $M$ non è un maggiorante dei valori, esiste
$\nu\in\N$ tale che $a_{\nu}>M$. Dato che $a_{n}$ è crescente, per ogni $n>\nu$ risulta
\begin{equation}\label{eq:monotone-passo3}
  a_{n}\ge a_{\nu}>M,
\end{equation}
da cui $\lim_{n\to+\infty}a_{n}=+\infty=\sup_{n}a_{n}$.

In modo analogo si trattano i casi relativi a successioni decrescenti; oppure ci si
riconduce ai precedenti osservando che, se $a_{n}$ è decrescente, la successione $-a_{n}$
è crescente, e che $\inf_{n}a_{n}=-\sup_{n}(-a_{n})$.
\end{dimostrazione}

\begin{figure}[htb]
\centering
\begin{tikzpicture}
\begin{axis}[width=0.92\textwidth, height=5cm, xmin=0, xmax=21, ymin=0, ymax=2.35,
  xtick={5,10,15,20}, extra x ticks={7}, extra x tick labels={$\nu$},
  ytick={2}, yticklabels={$\ell=\sup a_{n}$}, xlabel={$n$}]
  \fill[verde!15] (axis cs:0,1.7) rectangle (axis cs:21,2.3);
  \addplot[tratteggio, domain=0:21] {2};
  \addplot[only marks, mark=*, mark size=1.4pt, blu, samples at={1,...,20}] {2-1.6*0.78^x};
  \addplot[tratteggio] coordinates {(7,0) (7,1.7)};
  \node[font=\scriptsize, verde!50!black, anchor=east] at (axis cs:20.8,1.78) {$\ell-\epsilon$};
\end{axis}
\end{tikzpicture}
\caption{Una successione crescente e limitata: appena un termine $a_{\nu}$ supera
$\ell-\epsilon$, la monotonia costringe tutti i termini successivi a restare nella
striscia $(\ell-\epsilon,\ell]$.}
\label{fig:monotona-sup}
\end{figure}

\begin{osservazione}
La successione \eqref{eq:succ-alt-n} $a_{n}=(-1)^{n}/n$ non è monotòna; infatti i termini di
posto pari sono positivi, mentre quelli di posto dispari sono negativi: i termini della
successione oscillano intorno allo zero. La successione $(-1)^{n}/n$ è quindi un
\emph{esempio di successione convergente pur non essendo monotòna} (osserviamo che ciò non
contraddice il teorema: infatti nel teorema non si afferma che ogni successione
convergente è monotòna!). La successione non regolare \eqref{eq:succ-alt} $a_{n}=(-1)^{n}$
non è monotòna; ciò è in accordo con il teorema sulle successioni monotòne, perché se la
successione $(-1)^{n}$ fosse monotòna dovrebbe avere limite.
\end{osservazione}

Ricordando che una successione si dice regolare se essa ammette limite (finito o
infinito), il teorema precedente afferma che \emph{ogni successione monotòna è regolare}.
Due osservazioni tecniche ne estendono l'uso.

\begin{osservazione}[Successioni definitivamente monotòne e successioni traslate]\label{oss:traslata}
\leavevmode
\begin{enumerate}
  \item Il teorema vale anche per le successioni \emph{definitivamente} monotòne, cioè
  monotòne a partire da un certo indice $n_{0}$: basta applicarlo alla successione
  $a_{n_{0}},a_{n_{0}+1},\dots$, perché i primi termini non influiscono sul limite. In tal
  caso, però, il limite non è più necessariamente l'estremo superiore di tutti i termini.
  \item Se $a_{n}\to a$ (finito o infinito), allora anche la successione «traslata»
  $a_{n+1}$ tende ad $a$: se $\abs{a_{n}-a}<\epsilon$ per $n>\nu$, a maggior ragione
  $\abs{a_{n+1}-a}<\epsilon$ per $n>\nu$. Lo stesso vale per $a_{n+k}$, con $k\in\N$ fisso.
\end{enumerate}
\end{osservazione}

\subsection{Un'applicazione: il metodo di Erone per le radici quadrate}

Il teorema sulle successioni monotòne è lo strumento naturale per studiare le
\emph{successioni definite per ricorrenza}\index{successione!definita per ricorrenza},
in cui ogni termine è calcolato a partire dal precedente. È così che funzionano quasi
tutti gli algoritmi numerici iterativi.

\begin{esempio}[Il metodo di Erone]\label{es:erone}\index{metodo di Erone}
Consideriamo la successione definita da
\begin{equation}\label{eq:erone}
  x_{1}=2, \qquad x_{n+1}=\frac12\Bigl(x_{n}+\frac2{x_{n}}\Bigr) \quad (n\in\N).
\end{equation}
L'idea è semplice: se $x_{n}$ è un'approssimazione di $\sqrt2$ per eccesso, allora
$2/x_{n}$ lo è per difetto (il loro prodotto vale $2$), e la loro media è
un'approssimazione migliore. Dimostriamo che $x_{n}\to\sqrt2$.

\emph{Tutti i termini sono positivi}: per induzione, se $x_{n}>0$ anche $x_{n+1}>0$.

\emph{Per $n\ge2$ si ha $x_{n}\ge\sqrt2$}: infatti, per ogni $n$,
\[
  x_{n+1}-\sqrt2=\frac{x_{n}^{2}-2\sqrt2\,x_{n}+2}{2x_{n}}=\frac{(x_{n}-\sqrt2)^{2}}{2x_{n}}\ge0
\]
(ed è $x_{1}=2\ge\sqrt2$ anche per $n=1$).

\emph{La successione è decrescente}: per ogni $n$,
\[
  x_{n+1}-x_{n}=\frac{2-x_{n}^{2}}{2x_{n}}\le0,
\]
perché $x_{n}\ge\sqrt2$, cioè $x_{n}^{2}\ge2$.

Per il \cref{teo:monotone}, essendo decrescente e limitata inferiormente da $\sqrt2$, la
successione converge a un limite $L$, e per il \cref{cor:confronto-limiti} $L\ge\sqrt2>0$.
Passiamo al limite nella \eqref{eq:erone}: il primo membro $x_{n+1}$ tende a $L$
(\cref{oss:traslata}), e il secondo, per il \cref{teo:operazioni-limiti}, tende a
$\frac12\bigl(L+\frac2L\bigr)$. Per l'unicità del limite
$L=\frac12\bigl(L+\frac2L\bigr)$, cioè $L^{2}=2$, e quindi $L=\sqrt2$.
\end{esempio}

\begin{attenzione}[Prima l'esistenza, poi il valore]
Il passaggio al limite nella relazione di ricorrenza è lecito \emph{solo dopo} aver
dimostrato che il limite esiste. Senza questa precauzione si possono ottenere assurdità:
per la successione $x_{1}=1$, $x_{n+1}=2x_{n}$ (cioè $x_{n}=2^{n-1}$), «passando al
limite» si otterrebbe $L=2L$, cioè $L=0$; ma la successione diverge a $+\infty$.
\end{attenzione}

\begin{python}[La convergenza del metodo di Erone]
from decimal import Decimal, getcontext
getcontext().prec = 60              # lavoriamo con 60 cifre decimali

radice_vera = Decimal(2).sqrt()
x = Decimal(2)
for n in range(1, 8):
    errore = x - radice_vera
    print(f"x_{n} = {str(x)[:34]}   errore = {float(errore):.1e}")
    x = (x + 2 / x) / 2
\end{python}
\begin{risultato}
x_1 = 2   errore = 5.9e-01
x_2 = 1.5   errore = 8.6e-02
x_3 = 1.41666666666666666666666666666666   errore = 2.5e-03
x_4 = 1.41421568627450980392156862745098   errore = 2.1e-06
x_5 = 1.41421356237468991062629557889013   errore = 1.6e-12
x_6 = 1.41421356237309504880168962350253   errore = 9.0e-25
x_7 = 1.41421356237309504880168872420969   errore = 2.9e-49
\end{risultato}

Il programma mostra un fenomeno notevole: a ogni passo il numero di cifre esatte
\emph{raddoppia} (l'errore passa da $10^{-6}$ a $10^{-12}$, poi a $10^{-24}$). È la
cosiddetta \emph{convergenza quadratica}, che si legge direttamente nella formula
$x_{n+1}-\sqrt2=\frac{(x_{n}-\sqrt2)^{2}}{2x_{n}}$: l'errore al passo successivo è
proporzionale al \emph{quadrato} dell'errore attuale (si veda anche la \cref{fig:erone}). Il
metodo di Erone è un caso
particolare del metodo di Newton per la ricerca degli zeri, che studieremo con le derivate
e che è alla base dei solutori non lineari usati nel calcolo strutturale e nella
fluidodinamica computazionale.

\begin{figure}[htb]
\centering
\includegraphics[width=0.62\textwidth]{erone}
\caption{Il diagramma «a ragnatela» del metodo di Erone. Per rendere visibili più passi
si parte da $x_{1}=3$: dal punto $(x_{1},x_{1})$ della bisettrice $y=x$ si scende
verticalmente alla curva $y=\frac12(x+\frac2x)$, trovando $x_{2}$, poi ci si sposta
orizzontalmente fino alla bisettrice, e così via. Le iterate scendono verso il punto
fisso $\sqrt2$, dove la curva incontra la bisettrice (\texttt{python/fig\_erone.py}).}
\label{fig:erone}
\end{figure}

\begin{notastorica}[Le radici quadrate nell'antichità]
Il metodo porta il nome di Erone di Alessandria (I secolo d.C.), che lo descrive nei
\emph{Metrica} per calcolare $\sqrt{720}$; ma era noto molto prima. La tavoletta
babilonese YBC 7289, conservata a Yale e risalente a circa il 1800 a.C., mostra un
quadrato con le diagonali e, scritto in base $60$, il numero $1;24,51,10$, cioè
$1+\frac{24}{60}+\frac{51}{60^{2}}+\frac{10}{60^{3}}\approx1{,}41421296$: un'approssimazione
di $\sqrt2=1{,}41421356\ldots$ con un errore di circa $6\cdot10^{-7}$ (cinque decimali
esatti). Curiosamente, essa coincide con il troncamento a tre cifre sessagesimali di
$x_{4}=577/408=1;24,51,10,35\ldots$, il quarto termine del metodo di Erone a partire da
$x_{1}=2$. Secondo alcuni storici i Babilonesi conoscevano già un procedimento iterativo
di questo tipo, che sarebbe allora il più antico algoritmo della storia.
\end{notastorica}

%% ----------------------------------------------------------------------------
%%  §3.8  Il numero e
%% ----------------------------------------------------------------------------
\section{Il numero $e$}
\label{sec:numero-e}

Il teorema sulle successioni monotòne è utile per definire il \emph{numero di Nepero}
$e$ come limite di una particolare successione monotòna e limitata. Infatti,
introduciamo tale numero mediante il limite
\begin{equation}\label{eq:def-e}
  e=\lim_{n\to+\infty}a_{n}, \qquad\text{con}\qquad a_{n}=\Bigl(1+\frac1n\Bigr)^{n}.
\end{equation}
La \eqref{eq:def-e} è la definizione del \emph{numero di Nepero}\index{numero di Nepero} $e$.
Tale definizione è giustificata dal fatto che, come proveremo tra poco, \emph{la
successione $a_{n}=(1+1/n)^{n}$ è strettamente crescente e limitata}; quindi esiste, ed è
un numero reale, il limite per $n\to+\infty$ di $a_{n}$.

\begin{idea}[Due tendenze opposte]
La successione $(1+1/n)^{n}$ combina fra loro due tendenze opposte: la tendenza della
base a un limite uguale a $1$, e la tendenza dell'esponente a un limite infinito. Se
prevalesse la prima, il limite sarebbe $1$ (come per $1^{n}$); se prevalesse la seconda,
sarebbe $+\infty$ (come per $1{,}01^{n}$, \cref{teo:lim-geometrica}). Il risultato è
invece un numero intermedio tra $1$ e $+\infty$, che non è né l'uno né l'altro: è
proprio questo equilibrio delicato a rendere $e$ così importante. La forma $1^{\infty}$ è
una nuova forma indeterminata, come vedremo alla fine del paragrafo.
\end{idea}

\subsection{Monotonia e limitatezza}

Ci serve una versione «stretta» della disuguaglianza di Bernoulli.

\begin{lemma}[Bernoulli in forma stretta]\label{lem:bernoulli-stretta}
Se $x\ge-1$, $x\neq0$ e $n\ge2$ è un numero naturale, allora $(1+x)^{n}>1+nx$.
\end{lemma}

\begin{schema}
\Ipotesi{$x\ge-1$, \ $x\neq0$; \ $n\in\N$, \ $n\ge2$.}
\Tesi{$(1+x)^{n}>1+nx$.}
\end{schema}

\begin{strategia}
Si ripete la dimostrazione per induzione della disuguaglianza di Bernoulli, partendo da
$n=2$, dove la disuguaglianza è stretta perché $x^{2}>0$. Nel passo induttivo si
moltiplica per $1+x$, che deve essere \emph{positivo} perché la disuguaglianza stretta
si conservi: il caso $x=-1$, in cui $1+x=0$, va quindi controllato a parte.
\end{strategia}

\begin{dimostrazione}
Per $n=2$: $(1+x)^{2}=1+2x+x^{2}>1+2x$, perché $x\neq0$. Supponiamo la tesi vera per un
certo $n\ge2$. Se $x=-1$, si ha $(1+x)^{n+1}=0>1-(n+1)=1+(n+1)x$. Se $x>-1$, moltiplicando
l'ipotesi induttiva per $1+x>0$,
\[
  (1+x)^{n+1}>(1+nx)(1+x)=1+(n+1)x+nx^{2}\ge1+(n+1)x .
\]
Per il principio di induzione (a partire da $n=2$) la tesi vale per ogni $n\ge2$.
\end{dimostrazione}

\begin{teorema}[Esistenza del numero $e$]\label{teo:numero-e}\index{numero di Nepero!esistenza}
Le successioni
\begin{equation}\label{eq:an-bn}
  a_{n}=\Bigl(1+\frac1n\Bigr)^{n}, \qquad b_{n}=\Bigl(1+\frac1n\Bigr)^{n+1}
\end{equation}
sono, rispettivamente, strettamente crescente e strettamente decrescente, e risulta
$a_{n}<b_{n}$ per ogni $n$. Di conseguenza entrambe convergono, allo stesso limite $e$, e
\begin{equation}\label{eq:e-stima}
  a_{n}<e<b_{n} \qquad\perogni n\in\N .
\end{equation}
\end{teorema}

\begin{schema}
\Ipotesi{$a_{n}$ e $b_{n}$ definite dalla \eqref{eq:an-bn}, \ $n\in\N$.}
\Tesi{(i) $a_{n}<a_{n+1}$; \ (ii) $b_{n}>b_{n+1}$; \ (iii) $a_{n}<b_{n}$; \
(iv) $a_{n}$ e $b_{n}$ convergono allo stesso limite $e$, e $a_{n}<e<b_{n}$.}
\end{schema}

\begin{strategia}
Il cuore della dimostrazione sono le due monotonie, ed entrambe si ottengono con lo
stesso trucco. Si scrive il rapporto tra due termini consecutivi, si mettono insieme le
potenze con lo stesso esponente, e si ottiene una disuguaglianza della forma
$\bigl(1+x\bigr)^{n}\gtrless1+(\text{qualcosa})$, con $x=-\frac1{n^{2}}$ per $a_{n}$ e
$x=\frac1{n^{2}-1}$ per $b_{n}$: esattamente la forma della disuguaglianza di Bernoulli.
Poi la conclusione è quasi automatica: $a_{n}$ cresce, $b_{n}$ decresce, e $a_{n}<b_{n}$;
quindi ogni $a_{n}$ è minore di ogni $b_{m}$ (una «tenaglia» come quella di Archimede nel
\cref{sec:premessa-limiti}). In particolare $a_{n}$ è limitata superiormente da $b_{1}=4$,
e il teorema sulle successioni monotòne dà la convergenza. Che i due limiti coincidano
segue dal fatto che $b_{n}=a_{n}(1+\frac1n)$ e $1+\frac1n\to1$.
\end{strategia}

\begin{dimostrazione}
(i) \emph{$a_{n}$ è strettamente crescente.} La tesi è
\begin{equation}\label{eq:an-cresce-1}
  a_{n}=\Bigl(1+\frac1n\Bigr)^{n}>a_{n-1}=\Bigl(1+\frac1{n-1}\Bigr)^{n-1} \qquad\perogni n\ge2 .
\end{equation}
Eseguendo la somma delle frazioni, si può scrivere in modo equivalente
\begin{equation}\label{eq:an-cresce-2}
  \Bigl(\frac{n+1}{n}\Bigr)^{n}>\Bigl(\frac{n}{n-1}\Bigr)^{n}\cdot\frac{n-1}{n},
\end{equation}
cioè, dividendo per $\bigl(\frac{n}{n-1}\bigr)^{n}>0$ e raccogliendo le potenze,
\begin{equation}\label{eq:an-cresce-3}
  \Bigl(\frac{n^{2}-1}{n^{2}}\Bigr)^{n}>\frac{n-1}{n}.
\end{equation}
È conveniente isolare come addendo $1$, nel modo seguente:
\begin{equation}\label{eq:an-cresce-4}
  \Bigl(1-\frac1{n^{2}}\Bigr)^{n}>1-\frac1n .
\end{equation}
La \eqref{eq:an-cresce-4} è equivalente alla tesi da provare; e si ottiene dal
\cref{lem:bernoulli-stretta} con $x=-1/n^{2}$, che è maggiore di $-1$ e diverso da zero:
$\bigl(1-\frac1{n^{2}}\bigr)^{n}>1-n\cdot\frac1{n^{2}}=1-\frac1n$.

(ii) \emph{$b_{n}$ è strettamente decrescente.} Procediamo come nella dimostrazione
precedente. La tesi è
\begin{equation}\label{eq:bn-decresce-1}
  b_{n-1}=\Bigl(1+\frac1{n-1}\Bigr)^{n}>b_{n}=\Bigl(1+\frac1n\Bigr)^{n+1} \qquad\perogni n\ge2,
\end{equation}
cioè, equivalentemente,
\begin{equation}\label{eq:bn-decresce-2}
\begin{gathered}
  \Bigl(\frac{n}{n-1}\Bigr)^{n}>\Bigl(\frac{n+1}{n}\Bigr)^{n}\cdot\frac{n+1}{n},\\
  \text{cioè}\qquad\Bigl(\frac{n^{2}}{n^{2}-1}\Bigr)^{n}>\frac{n+1}{n}.
\end{gathered}
\end{equation}
Come nella dimostrazione precedente isoliamo $1$ come addendo:
\begin{equation}\label{eq:bn-decresce-3}
  \Bigl(1+\frac1{n^{2}-1}\Bigr)^{n}>1+\frac1n .
\end{equation}
La disuguaglianza di Bernoulli, con $x=1/(n^{2}-1)$, dà la conclusione:
\begin{equation}\label{eq:bn-decresce-4}
  \Bigl(1+\frac1{n^{2}-1}\Bigr)^{n}\ge1+\frac{n}{n^{2}-1}>1+\frac1n \qquad\perogni n\ge2
\end{equation}
(nell'ultimo passaggio si è utilizzata la disuguaglianza $\frac{n}{n^{2}-1}>\frac1n$, che
è vera perché equivale a $n^{2}>n^{2}-1$).

(iii) Per ogni $n\in\N$ si ha
\begin{equation}\label{eq:bn-maggiore-an}
  b_{n}=\Bigl(1+\frac1n\Bigr)^{n}\cdot\Bigl(1+\frac1n\Bigr)=a_{n}\cdot\Bigl(1+\frac1n\Bigr)>a_{n}.
\end{equation}

(iv) Dalle (i), (ii), (iii) segue che, per ogni $n,m\in\N$, posto $k=\max\{n,m\}$,
\begin{equation}\label{eq:an-minore-bm}
  a_{n}\le a_{k}<b_{k}\le b_{m};
\end{equation}
in particolare $2=a_{1}\le a_{n}<b_{1}=4$ per ogni $n$. La successione $a_{n}$ è dunque
crescente e limitata, e per il \cref{teo:monotone} converge a un numero reale, che
chiamiamo $e$; analogamente $b_{n}$, decrescente e limitata inferiormente da $a_{1}=2$,
converge. Dalla \eqref{eq:bn-maggiore-an} e dal \cref{teo:operazioni-limiti},
$\lim b_{n}=\lim a_{n}\cdot\lim(1+\frac1n)=e\cdot1=e$. Infine, poiché $a_{n}$ è strettamente
crescente, $a_{n}<a_{n+1}\le\sup_{k}a_{k}=e$; e poiché $b_{n}$ è strettamente
decrescente, $e=\inf_{k}b_{k}\le b_{n+1}<b_{n}$ (si veda la \eqref{eq:lim-monotone}).
\end{dimostrazione}

\begin{strategia}[Un'altra via]
Seguendo \cite{giusti}, si può sviluppare la potenza con la formula del binomio, invece di
confrontare due termini consecutivi. Si ottiene una somma di addendi positivi che
crescono con $n$, e il cui numero cresce con $n$: la crescenza diventa evidente. Ogni
addendo è poi minore di un reciproco di fattoriale, e la limitatezza segue dal confronto
con una progressione geometrica. Si perde la tenaglia \eqref{eq:e-stima}, ma si scopre da
dove viene il numero $e$.
\end{strategia}

\begin{dimostrazione}[Seconda dimostrazione della crescenza e della limitatezza di $a_{n}$]
Per la formula del binomio (\cref{teo:binomio}) con $a=1$ e $b=\frac1n$, e per la
\eqref{eq:combinazioni},
\[
  a_{n}=\sum_{k=0}^{n}\binom nk\frac1{n^{k}}
  =\sum_{k=0}^{n}\frac1{k!}\cdot\frac{n(n-1)\cdots(n-k+1)}{n^{k}} .
\]
La seconda frazione è il prodotto dei $k$ fattori $\frac nn,\frac{n-1}n,\dots,\frac{n-k+1}n$.
Isolando gli addendi con $k=0$ e $k=1$, che valgono $1$, si ottiene
\begin{equation}\label{eq:e-binomio}
  a_{n}=1+1+\sum_{k=2}^{n}\frac1{k!}\Bigl(1-\frac1n\Bigr)\Bigl(1-\frac2n\Bigr)\cdots
  \Bigl(1-\frac{k-1}n\Bigr)
\end{equation}
(per $n=1$ la somma è vuota).

\emph{Crescenza.} Passando da $n$ a $n+1$, ogni fattore $1-\frac jn$, con
$1\le j\le k-1\le n-1$, è positivo e aumenta, perché diventa $1-\frac j{n+1}$; quindi
aumenta ogni addendo della somma \eqref{eq:e-binomio}. Inoltre $a_{n+1}$ contiene un
addendo positivo in più, quello con $k=n+1$. Dunque $a_{n}<a_{n+1}$.

\emph{Limitatezza.} Ogni prodotto nella \eqref{eq:e-binomio} è compreso tra $0$ e $1$.
Inoltre per $k\ge1$ si ha $k!\ge2^{k-1}$, perché ciascuno dei $k-1$ fattori $2,3,\dots,k$
è almeno $2$. Per la \cref{prop:geometrica} con $x=\frac12$ si ha allora
\begin{equation}\label{eq:e-binomio-stima}
  a_{n}\le1+\sum_{k=1}^{n}\frac1{k!}\le1+\sum_{k=1}^{n}\frac1{2^{k-1}}
  =1+\frac{1-(1/2)^{n}}{1-1/2}<3 .
\end{equation}
La successione $a_{n}$ è dunque crescente e limitata superiormente da $3$. Per il
\cref{teo:monotone} converge, e il suo limite $e$ verifica $e\le3$.
\end{dimostrazione}

\begin{osservazione}[Che cosa insegna ciascuna via]\label{oss:e-due-vie}
La prima dimostrazione dà di più sul piano quantitativo: la successione decrescente
$b_{n}$ e la tenaglia $a_{n}<e<b_{n}$, cioè un controllo dell'errore. La seconda dice
invece da dove viene $e$. Posto $s_{n}=\sum_{k=0}^{n}\frac1{k!}$, la
\eqref{eq:e-binomio-stima} dice che $a_{n}\le s_{n}<3$; e con gli stessi strumenti si
prova che $s_{n}\to e$. La successione $s_{n}$ è crescente e limitata, quindi converge a
un limite $s$, e da $a_{n}\le s_{n}$ segue $e\le s$ (\cref{cor:confronto-limiti}).
Viceversa, fissato $m$, per $n\ge m$ si possono trascurare nella \eqref{eq:e-binomio}
gli addendi con $k>m$, che sono positivi:
\[
  a_{n}\ge1+1+\sum_{k=2}^{m}\frac1{k!}\Bigl(1-\frac1n\Bigr)\cdots\Bigl(1-\frac{k-1}n\Bigr).
\]
Per $n\to+\infty$, con $m$ fisso, il secondo membro è una somma di un numero \emph{fisso}
di termini, ciascuno dei quali tende a $\frac1{k!}$ (\cref{teo:operazioni-limiti}):
tende quindi a $s_{m}$, e si ottiene $e\ge s_{m}$. Passando al limite per $m\to+\infty$
si ha $e\ge s$. Dunque $e=\lim_{n}s_{n}$ è la somma della serie dei reciproci dei
fattoriali (\cref{es:serie-fattoriali}), ottenuta qui senza la formula di Taylor. Il punto
delicato è l'ordine dei passaggi al limite: prima $n\to+\infty$ con $m$ fisso, poi
$m\to+\infty$. Nella \eqref{eq:e-binomio} non si può passare al limite addendo per
addendo, perché il numero degli addendi cresce con $n$.
\end{osservazione}

\begin{corollario}\label{cor:e-tra-2-e-3}
Risulta $2<e<3$.
\end{corollario}

\begin{schema}
\Ipotesi{la stima \eqref{eq:e-stima}: $a_{n}<e<b_{n}$ per ogni $n\in\N$.}
\Tesi{$2<e<3$.}
\end{schema}

\begin{strategia}
Basta scegliere nella tenaglia \eqref{eq:e-stima} due indici opportuni: $a_{1}=2$ dà la
stima dal basso; per quella dall'alto occorre un $b_{n}$ minore di $3$, e un calcolo
diretto mostra che $n=5$ è sufficiente.
\end{strategia}

\begin{dimostrazione}
Dalla \eqref{eq:e-stima} con $n=1$ si ha $e>a_{1}=2$, e con $n=5$
\begin{equation}\label{eq:b5}
  e<b_{5}=\Bigl(\frac65\Bigr)^{6}=\frac{46656}{15625}=2{,}985984<3. \qedhere
\end{equation}
\end{dimostrazione}

La \eqref{eq:e-stima} è un'ottima «tenaglia» numerica: aumentando $n$ si ottengono stime
sempre più precise. Riportiamo nella \cref{tab:numero-e} alcuni valori approssimati di
$a_{n}$ e $b_{n}$, rappresentati anche nella \cref{fig:numero-e}. Un metodo molto più efficiente per calcolare espressioni decimali
approssimate di $e$ verrà dalla formula di Taylor (\cref{es:calcolo-e}) e dalla serie dei
reciproci dei fattoriali, $e=\sum_{k=0}^{\infty}\frac1{k!}$ (\cref{oss:e-due-vie},
\cref{es:serie-fattoriali} e \eqref{eq:e-serie}); con esso si ottengono facilmente
valori come
\begin{equation}\label{eq:e-decimali}
  e=2{,}71828182845904523536\ldots
\end{equation}

\begin{table}[htb]
\centering
\renewcommand{\arraystretch}{1.3}
\footnotesize\setlength{\tabcolsep}{3pt}
\begin{tabular}{l*{8}{c}}
\toprule
$n$ & $1$ & $10$ & $50$ & $100$ & $500$ & $1000$ & $10^{5}$ & $10^{6}$\\
\midrule
$(1+1/n)^{n}$ & $2$ & $2{,}5937$ & $2{,}6916$ & $2{,}7048$ & $2{,}7156$ & $2{,}7169$ & $2{,}71827$ & $2{,}718280$\\
$(1+1/n)^{n+1}$ & $4$ & $2{,}8531$ & $2{,}7454$ & $2{,}7319$ & $2{,}7210$ & $2{,}7196$ & $2{,}71830$ & $2{,}718283$\\
\bottomrule
\end{tabular}
\caption{Approssimazioni per difetto ($a_{n}$) e per eccesso ($b_{n}$) del numero $e$.
La convergenza è lenta: con $n=10^{6}$ si ottengono appena cinque cifre decimali
esatte.}
\label{tab:numero-e}
\end{table}

\begin{figure}[htb]
\centering
\includegraphics[width=0.85\textwidth]{numero_e}
\caption{Le successioni $a_{n}=(1+1/n)^{n}$ (crescente, in blu) e
$b_{n}=(1+1/n)^{n+1}$ (decrescente, in rame) stringono il numero $e$ (tratteggiato)
come in una tenaglia (\texttt{python/fig\_numero\_e.py}).}
\label{fig:numero-e}
\end{figure}

\subsection{Generalizzazioni}

Nelle applicazioni sono utili anche i limiti seguenti, generalizzazioni della
definizione \eqref{eq:def-e} a successioni di esponenti reali qualsiasi. Ci serve
preliminarmente un lemma, che permette di «comporre» una successione con una successione
di indici.

\begin{lemma}\label{lem:composizione-indici}
Sia $c_{m}$ una successione con $\lim_{m\to+\infty}c_{m}=L\in\R$, e sia $m_{n}$ una
successione di numeri naturali tale che $m_{n}\to+\infty$. Allora $c_{m_{n}}\to L$.
\end{lemma}

\begin{schema}
\Ipotesi{$c_{m}\to L\in\R$ per $m\to+\infty$; \ $m_{n}\in\N$ e $m_{n}\to+\infty$.}
\Tesi{$c_{m_{n}}\to L$ per $n\to+\infty$.}
\end{schema}

\begin{strategia}
Due definizioni di limite vanno incatenate. Fissato $\epsilon$, la convergenza di $c_{m}$
fornisce una soglia $\mu$ per l'indice $m$; la divergenza di $m_{n}$ garantisce che
definitivamente $m_{n}$ supera quella soglia. Si noti che, a differenza delle estratte,
qui gli indici $m_{n}$ possono ripetersi o non essere crescenti.
\end{strategia}

\begin{dimostrazione}
Fissato $\epsilon>0$, esiste $\mu$ tale che $\abs{c_{m}-L}<\epsilon$ per ogni $m>\mu$.
Poiché $m_{n}\to+\infty$, esiste $\nu$ tale che $m_{n}>\mu$ per ogni $n>\nu$ (si applichi
la \eqref{eq:def-piu-inf} con $M=\max\{\mu,1\}$). Per tali $n$ si ha quindi
$\abs{c_{m_{n}}-L}<\epsilon$.
\end{dimostrazione}

\begin{teorema}\label{teo:e-generalizzato}
Per ogni successione di numeri reali $a_{n}$ (con $a_{n}\neq0$ e $1+1/a_{n}>0$) si ha
\begin{align}
  a_{n}\to+\infty \quad&\implica\quad \Bigl(1+\frac1{a_{n}}\Bigr)^{a_{n}}\to e; \label{eq:e-piu-inf}\\
  a_{n}\to-\infty \quad&\implica\quad \Bigl(1+\frac1{a_{n}}\Bigr)^{a_{n}}\to e. \label{eq:e-meno-inf}
\end{align}
\end{teorema}

\begin{schema}
\Ipotesi{$a_{n}\neq0$, \ $1+\frac1{a_{n}}>0$; \ $a_{n}\to+\infty$ oppure
$a_{n}\to-\infty$.}
\Tesi{$\Bigl(1+\dfrac1{a_{n}}\Bigr)^{a_{n}}\to e$.}
\end{schema}

\begin{strategia}
Per la \eqref{eq:e-piu-inf} confrontiamo con il caso intero mediante la parte intera
$m_{n}=\lfloor a_{n}\rfloor$: poiché $m_{n}\le a_{n}<m_{n}+1$, la base $1+\frac1{a_{n}}$ è
compresa tra $1+\frac1{m_{n}+1}$ e $1+\frac1{m_{n}}$, e l'esponente tra $m_{n}$ e
$m_{n}+1$. Prendendo la base più piccola con l'esponente più piccolo, e la base più
grande con l'esponente più grande, otteniamo due carabinieri che sono, a meno di
fattori che tendono a $1$, proprio le successioni $a_{m}$ e $b_{m}$ del teorema
precedente, valutate in $m=m_{n}\to+\infty$. Per la \eqref{eq:e-meno-inf} un cambiamento
di variabile, $b_{n}=-a_{n}-1$, la riconduce alla precedente.
\end{strategia}

\begin{dimostrazione}[Dimostrazione della \eqref{eq:e-piu-inf}]
Poiché $a_{n}\to+\infty$, definitivamente $a_{n}\ge1$; consideriamo solo tali indici.
Indichiamo con $m_{n}=\lfloor a_{n}\rfloor$ la parte intera di $a_{n}$
(\cref{prop:parte-intera}), il più grande intero minore o uguale ad $a_{n}$: risulta
$m_{n}\in\N$, $m_{n}\le a_{n}<m_{n}+1$, e quindi $m_{n}>a_{n}-1\to+\infty$. Si ha
\begin{equation}\label{eq:e-carabinieri}
  \Bigl(1+\frac1{m_{n}+1}\Bigr)^{m_{n}}\le\Bigl(1+\frac1{a_{n}}\Bigr)^{a_{n}}\le\Bigl(1+\frac1{m_{n}}\Bigr)^{m_{n}+1},
\end{equation}
perché la base e l'esponente del primo membro sono minori o uguali rispettivamente della
base e dell'esponente del secondo membro, e lo stesso accade tra secondo e terzo membro:
si usano la monotonia della potenza rispetto all'esponente per basi maggiori di $1$
\eqref{eq:pot5} e la monotonia rispetto alla base per esponenti positivi \eqref{eq:pot3}.
Il risultato segue dal teorema dei carabinieri, perché il primo e il terzo membro tendono
ad $e$. Infatti, per il terzo membro, $\bigl(1+\frac1{m}\bigr)^{m+1}=b_{m}\to e$; per il
primo,
\begin{equation}\label{eq:e-primo-membro}
  \Bigl(1+\frac1{m+1}\Bigr)^{m}=\frac{\bigl(1+\frac1{m+1}\bigr)^{m+1}}{1+\frac1{m+1}}
  =\frac{a_{m+1}}{1+\frac1{m+1}}\to\frac e1=e ;
\end{equation}
in entrambi i casi per $m\to+\infty$, e quindi anche per $m=m_{n}$, per il
\cref{lem:composizione-indici}.
\end{dimostrazione}

\begin{dimostrazione}[Dimostrazione della \eqref{eq:e-meno-inf}]
Poniamo
\begin{equation}\label{eq:e-cambio}
  b_{n}=-a_{n}-1 .
\end{equation}
Dato che $a_{n}\to-\infty$, la successione $b_{n}$ diverge a $+\infty$; in particolare
definitivamente $b_{n}>0$. Inoltre $a_{n}=-(b_{n}+1)$, e quindi
\begin{equation}\label{eq:e-meno-inf-passo}
\begin{split}
  \Bigl(1+\frac1{a_{n}}\Bigr)^{a_{n}}&=\Bigl(1-\frac1{b_{n}+1}\Bigr)^{-(b_{n}+1)}
  =\Bigl(\frac{b_{n}}{b_{n}+1}\Bigr)^{-(b_{n}+1)}\\
  &=\Bigl(\frac{b_{n}+1}{b_{n}}\Bigr)^{b_{n}+1}
  =\Bigl(1+\frac1{b_{n}}\Bigr)^{b_{n}}\cdot\Bigl(1+\frac1{b_{n}}\Bigr).
\end{split}
\end{equation}
Ricordando che $b_{n}\to+\infty$, l'ultimo membro converge a $e\cdot1=e$, a causa della
\eqref{eq:e-piu-inf} e della \eqref{eq:inf5}; ciò prova la \eqref{eq:e-meno-inf}.
\end{dimostrazione}

\begin{osservazione}
Nelle due dimostrazioni, e in alcune delle successive, abbiamo usato le proprietà delle
potenze a esponente reale. La monotonia rispetto all'esponente \eqref{eq:pot5} è stata
dimostrata nel \cref{teo:esponenziale}(ii). Le regole algebriche \eqref{eq:pot1} e la
monotonia rispetto alla base \eqref{eq:pot3}, invece, per esponenti irrazionali sono
state soltanto discusse nel \cref{sec:esponenziale-reale}: la loro dimostrazione
completa verrà data quando avremo a disposizione la continuità della funzione
esponenziale.
\end{osservazione}

\begin{esempio}\label{es:limiti-e}
Dal \cref{teo:e-generalizzato} e dalle operazioni con i limiti seguono molti limiti
utili. Ad esempio:
\begin{itemize}
  \item $\bigl(1+\frac1n\bigr)^{2n}=\bigl[\bigl(1+\frac1n\bigr)^{n}\bigr]^{2}\to e^{2}$;
  \item $\bigl(1-\frac1n\bigr)^{n}\to\frac1e$: infatti, per la \eqref{eq:e-meno-inf} con
  $a_{n}=-n$ (per $n\ge2$, così che $1+1/a_{n}>0$), $\bigl(1-\frac1n\bigr)^{-n}\to e$, e il
  reciproco tende a $1/e$;
  \item $\bigl(\frac{n+1}{n-1}\bigr)^{n}=\frac{(1+1/n)^{n}}{(1-1/n)^{n}}\to\frac{e}{1/e}=e^{2}$
  (anche qui $n\ge2$).
\end{itemize}
\end{esempio}

\subsection{Nuove forme indeterminate}

La forma $1^{\infty}$ è una nuova forma indeterminata, come pure sono indeterminate le
forme $\infty^{0}$ e $0^{0}$. Quindi all'elenco \eqref{eq:forme-indeterminate} vanno
aggiunte le \emph{forme indeterminate}\index{forma indeterminata!esponenziale}
\begin{equation}\label{eq:forme-esponenziali}
  1^{(+\infty)}, \qquad 1^{(-\infty)}, \qquad (+\infty)^{0}, \qquad 0^{0} .
\end{equation}
Che la forma $1^{\infty}$ sia davvero indeterminata lo mostrano tre successioni in cui la
base tende a $1$ e l'esponente a $+\infty$, con limiti diversi:
\begin{gather*}
  \Bigl(1+\frac1n\Bigr)^{n}\to e, \qquad
  \Bigl(1+\frac1n\Bigr)^{n^{2}}\ge1+n^{2}\cdot\frac1n=1+n\to+\infty,\\
  1\le\Bigl(1+\frac1{n^{2}}\Bigr)^{n}=\Bigl[\Bigl(1+\frac1{n^{2}}\Bigr)^{n^{2}}\Bigr]^{1/n}<3^{1/n}\to1 .
\end{gather*}
Nella seconda si è usata la disuguaglianza di Bernoulli, nella terza la stima
$a_{n^{2}}<3$ del \cref{cor:e-tra-2-e-3} e la \cref{prop:radice-n-a}.

\begin{notastorica}[Il numero $e$ da Eulero a oggi]
Nell'\emph{Introductio in analysin infinitorum} (1748) Eulero mostrò che il limite
\eqref{eq:def-e} coincide con la somma della serie
$1+1+\frac1{2!}+\frac1{3!}+\frac1{4!}+\cdots$, che converge molto più rapidamente (la
ritroveremo nel \cref{cap:serie}, formula \eqref{eq:e-serie}), e ne calcolò le prime $23$ cifre decimali; nella
stessa opera introdusse la funzione esponenziale $e^{x}$ e il logaritmo naturale come sua
inversa, nella forma che usiamo ancora oggi. Il fatto che la successione $(1+1/n)^{n}$
converga così lentamente (si veda la \cref{tab:numero-e}) è una curiosità ben nota agli
studenti di calcolo numerico: con i numeri in virgola mobile a doppia precisione, per
$n$ molto grande il calcolo di $1+1/n$ perde tutte le cifre significative di $1/n$, e il
risultato «converge» a $1$ invece che ad $e$ (si veda il programma qui sotto).
\end{notastorica}

\begin{python}[Il numero $e$ con la tenaglia, e un'insidia numerica]
import math

for n in [10, 1000, 10**5, 10**7]:
    a = (1 + 1/n) ** n
    b = (1 + 1/n) ** (n + 1)
    print(f"n = {n:>8d}:  {a:.10f} < e < {b:.10f}"
          f"   (ampiezza {b - a:.1e})")
print(f"math.e        = {math.e:.10f}")

# Insidia della virgola mobile: per n enorme 1 + 1/n si arrotonda.
# Confrontiamo il calcolo "ingenuo" con quello basato su log1p.
for k in [9, 12, 15, 16]:
    n = 10**k
    ingenuo = (1 + 1/n) ** n
    accurato = math.exp(n * math.log1p(1/n))
    print(f"n = 1e{k}:  ingenuo = {ingenuo:.10f}"
          f"   con log1p = {accurato:.10f}")
\end{python}
\begin{risultato}
n =       10:  2.5937424601 < e < 2.8531167061   (ampiezza 2.6e-01)
n =     1000:  2.7169239322 < e < 2.7196408562   (ampiezza 2.7e-03)
n =   100000:  2.7182682372 < e < 2.7182954199   (ampiezza 2.7e-05)
n = 10000000:  2.7182816941 < e < 2.7182819660   (ampiezza 2.7e-07)
math.e        = 2.7182818285
n = 1e9:  ingenuo = 2.7182820520   con log1p = 2.7182818271
n = 1e12:  ingenuo = 2.7185234960   con log1p = 2.7182818285
n = 1e15:  ingenuo = 3.0350352065   con log1p = 2.7182818285
n = 1e16:  ingenuo = 1.0000000000   con log1p = 2.7182818285
\end{risultato}

Si noti che già per $n=10^{9}$ il valore calcolato «ingenuamente» supera $e$, in
contraddizione con la \eqref{eq:e-stima}: l'errore di arrotondamento nel calcolo di
$1+1/n$, amplificato dall'elevamento alla potenza $n$, ha superato l'errore di
troncamento. Per $n=10^{15}$ il risultato è addirittura maggiore di $3$, e per $n=10^{16}$
vale $1$, perché $1+10^{-16}$ viene arrotondato a $1$. La matematica dice che la
successione converge ad $e$; il calcolatore, se usato senza cautela, dice altro. Il
rimedio, in questo caso, è riscrivere la potenza come $e^{\,n\log(1+1/n)}$ e calcolare
$\log(1+x)$ con una funzione apposita (\texttt{math.log1p}), che non perde cifre per $x$
piccolo: la terza colonna mostra che così i valori restano corretti.

%% ----------------------------------------------------------------------------
%%  §3.9  Appendice: infiniti di ordine crescente
%% ----------------------------------------------------------------------------
\section{Appendice: infiniti di ordine crescente}
\label{sec:infiniti}

Molte successioni importanti divergono a $+\infty$, ma con «velocità» molto diverse. In
questo paragrafo impariamo a confrontarle; lo strumento principale è un criterio che
sfrutta il teorema sulle successioni monotòne, con lo stesso metodo utilizzato per il
metodo di Erone.

\begin{teorema}[Criterio del rapporto per le successioni]\label{teo:criterio-rapporto}\index{criterio del rapporto!per le successioni}
Sia $a_{n}$ una successione a termini positivi. Definiamo $b_{n}=a_{n+1}/a_{n}$. Se la
successione $b_{n}$ converge a un limite $b<1$, allora la successione $a_{n}$ tende a
zero.
\end{teorema}

\begin{schema}
\Ipotesi{$a_{n}>0$ per ogni $n$; \ $\dfrac{a_{n+1}}{a_{n}}\to b$, con $b<1$.}
\Tesi{$a_{n}\to0$.}
\end{schema}

\begin{strategia}
Se il rapporto tra un termine e il precedente tende a un numero minore di $1$, da un
certo punto in poi ogni termine è più piccolo del precedente: la successione è
definitivamente decrescente e, essendo positiva, limitata inferiormente da $0$; quindi
converge (\cref{teo:monotone}) a un limite $a\ge0$. Per identificare $a$ si passa al
limite nella relazione $a_{n+1}=b_{n}a_{n}$, come nel metodo di Erone: se fosse $a>0$, si
otterrebbe $b=a/a=1$, contro l'ipotesi.
\end{strategia}

\begin{dimostrazione}
Per il teorema della permanenza del segno (\cref{teo:permanenza-segno}), applicato alla
successione $1-b_{n}$, che converge a $1-b>0$, esiste un indice $\nu$ per cui $b_{n}<1$
per ogni $n>\nu$. Quindi $a_{n+1}/a_{n}<1$, cioè $a_{n+1}<a_{n}$, per ogni $n>\nu$. Il
teorema sulle successioni monotòne (nella forma della \cref{oss:traslata}) assicura
l'esistenza del limite $a$, che è un numero reale non negativo, dato che la successione
è definitivamente decrescente e a termini positivi (\cref{cor:limite-non-negativo}).
Supponendo per assurdo $a\neq0$ e passando al limite per $n\to+\infty$ nella relazione
$b_{n}=a_{n+1}/a_{n}$ (si ricordi che $a_{n+1}\to a$, \cref{oss:traslata}), si ottiene
$b=a/a=1$, in contrasto con l'ipotesi $b<1$. Pertanto risulta $a=0$.
\end{dimostrazione}

\begin{osservazione}
Se invece $b>1$, lo stesso ragionamento mostra che $a_{n}$ è definitivamente crescente,
diciamo per $n\ge n_{0}$, e quindi ha limite. Tale limite non può essere $0$, perché è
maggiore o uguale di $a_{n_{0}}>0$, e non può essere un numero positivo $a$, perché si
otterrebbe di nuovo $b=a/a=1$: dunque $a_{n}\to+\infty$. Allo stesso risultato si arriva
applicando il criterio alla successione $1/a_{n}$, il cui rapporto
$\frac{1/a_{n+1}}{1/a_{n}}=\frac1{b_{n}}$ tende a $\frac1b<1$: quindi $1/a_{n}\to0$ e,
poiché $a_{n}>0$, $a_{n}\to+\infty$ (si veda il riquadro «$1/0$ non è $+\infty$» nel
\cref{sec:forme-indeterminate}). È la versione «per divergenza» del criterio. Se infine $b=1$, il criterio non dà alcuna
informazione: $a_{n}=n$, $a_{n}=1$ e $a_{n}=1/n$ hanno tutte $b_{n}\to1$, ma
comportamenti diversi.
\end{osservazione}

\subsection{La gerarchia degli infiniti}

Applichiamo il criterio del rapporto al confronto delle successioni
\begin{equation}\label{eq:gerarchia}
  \log n; \qquad n^{b}; \qquad a^{n}; \qquad n!; \qquad n^{n},
\end{equation}
dove abbiamo scelto $b>0$ e $a>1$. Il simbolo $n!$ ($n$ \emph{fattoriale}) indica il
prodotto dei primi $n$ numeri naturali (\cref{sec:combinatoria}). Per $n\to+\infty$ le
cinque successioni \eqref{eq:gerarchia} tendono tutte a $+\infty$ (per $\log n$ e
$n^{b}$ lo verificheremo nel \cref{lem:potenze-log}, per $a^{n}$ è il
\cref{teo:lim-geometrica}, per $n!\ge n$ e $n^{n}\ge n$ è ovvio). Possiamo però dire che
sono \emph{infiniti di ordine crescente}\index{infiniti di ordine crescente}, nel senso
che i limiti dei rapporti valgono
\begin{equation}\label{eq:gerarchia-limiti}
  \lim_{n\to+\infty}\frac{\log n}{n^{b}}=\lim_{n\to+\infty}\frac{n^{b}}{a^{n}}
  =\lim_{n\to+\infty}\frac{a^{n}}{n!}=\lim_{n\to+\infty}\frac{n!}{n^{n}}=0 .
\end{equation}

\begin{lemma}\label{lem:potenze-log}
\leavevmode
\begin{enumerate}
  \item Per ogni $c>0$ si ha $n^{c}\to+\infty$, e quindi $n^{-c}\to0$.
  \item Per ogni $x\in\R$ si ha $e^{x}>x$; di conseguenza $\log y<y$ per ogni $y>0$.
  \item $\log n\to+\infty$.
  \item Per ogni $b>0$ si ha $(1+\frac1n)^{b}\to1$.
\end{enumerate}
\end{lemma}

\begin{schema}
\Ipotesi{$c>0$, \ $b>0$, \ $x\in\R$, \ $y>0$.}
\Tesi{(i) $n^{c}\to+\infty$, $n^{-c}\to0$; \ (ii) $e^{x}>x$, $\log y<y$;\\
(iii) $\log n\to+\infty$; \ (iv) $(1+\frac1n)^{b}\to1$.}
\end{schema}

\begin{strategia}
Sono quattro fatti elementari, che servono come mattoni nel teorema successivo. Per (i)
e (iii) basta «invertire» la disuguaglianza $n^{c}>M$ o $\log n>M$ con la funzione
inversa, che è crescente. Per (ii) si confronta $e^{x}$ con la potenza di esponente
intero $e^{\lfloor x\rfloor}$, stimata dal basso con Bernoulli. Per (iv) si intrappola
l'esponente reale $b$ tra $0$ e un intero $k\ge b$, e si usano i carabinieri.
\end{strategia}

\begin{dimostrazione}
(i) Fissato $M>0$, per la monotonia delle potenze \eqref{eq:pot3} si ha $n^{c}>M$ se e
solo se $n>M^{1/c}$. La seconda affermazione segue dalla \eqref{eq:inf5}.

(ii) Se $x\le0$ la tesi è ovvia, perché $e^{x}>0$. Sia $x>0$ e $m=\lfloor x\rfloor\ge0$.
Per la disuguaglianza di Bernoulli, ricordando che $e-1>1$ (\cref{cor:e-tra-2-e-3}),
$e^{m}=\bigl(1+(e-1)\bigr)^{m}\ge1+m(e-1)\ge1+m$ (per $m=0$ entrambi i membri valgono
$1$). Quindi, per la monotonia dell'esponenziale, $e^{x}\ge e^{m}\ge1+m>x$. Per la
seconda affermazione, posto $x=\log y$, si ha $y=e^{\log y}>\log y$.

(iii) Fissato $M>0$, poiché il logaritmo è crescente (\cref{prop:log-monotono}), si ha
$\log n>M$ se e solo se $n>e^{M}$.

(iv) Sia $k\in\N$ con $k\ge b$. Poiché $1+\frac1n>1$, per le \eqref{eq:pot5} si ha
$1\le(1+\frac1n)^{b}\le(1+\frac1n)^{k}$; il terzo membro è il prodotto di $k$ successioni
che tendono a $1$, e tende a $1$. Si conclude con il teorema dei carabinieri.
\end{dimostrazione}

\begin{teorema}[Gerarchia degli infiniti]\label{teo:gerarchia}
Per ogni $b>0$ e ogni $a>1$ valgono i limiti \eqref{eq:gerarchia-limiti}.
\end{teorema}

\begin{schema}
\Ipotesi{$b>0$, \ $a>1$.}
\Tesi{$\dfrac{\log n}{n^{b}}\to0$, \ $\dfrac{n^{b}}{a^{n}}\to0$, \
$\dfrac{a^{n}}{n!}\to0$, \ $\dfrac{n!}{n^{n}}\to0$.}
\end{schema}

\begin{strategia}
Gli ultimi tre limiti sono tutti della forma $a_{n}\to0$ con $a_{n}>0$, e il criterio del
rapporto è lo strumento ideale: il rapporto $a_{n+1}/a_{n}$ di quozienti di potenze e
fattoriali si semplifica in modo spettacolare (ad esempio $\frac{(n+1)!}{n!}=n+1$). Per il
primo limite, invece, il criterio del rapporto non serve (il rapporto tende a $1$), e il
testo di riferimento \cite{marcellini-sbordone} rimanda la dimostrazione a un capitolo
successivo. Possiamo però darla subito, usando la disuguaglianza elementare $\log y<y$
del \cref{lem:potenze-log}: applicata a $y=n^{b/2}$, essa dice che $\log n$ è più piccolo
di un multiplo di $n^{b/2}$, e quindi $\frac{\log n}{n^{b}}$ è più piccolo di un multiplo di
$n^{-b/2}$, che è infinitesimo. Il trucco di usare «metà dell'esponente» è tipico: si
sacrifica un po' della potenza per ottenere una stima semplice.
\end{strategia}

\begin{dimostrazione}
\emph{Primo limite.} Per il \cref{lem:potenze-log}(ii) applicato a $y=n^{b/2}$, e per la
\eqref{eq:log3},
\[
  \frac b2\log n=\log n^{b/2}<n^{b/2},
\]
quindi
\[
  0\le\frac{\log n}{n^{b}}<\frac2b\cdot\frac{n^{b/2}}{n^{b}}=\frac2b\,n^{-b/2}
  \qquad\perogni n\in\N .
\]
Per il \cref{lem:potenze-log}(i) il terzo membro tende a zero, e la tesi segue dal
teorema dei carabinieri.

\emph{Secondo limite.} Poniamo
\begin{equation}\label{eq:gerarchia-2}
  a_{n}=\frac{n^{b}}{a^{n}}; \qquad
  r_{n}=\frac{a_{n+1}}{a_{n}}=\Bigl(\frac{n+1}{n}\Bigr)^{b}\frac1a\to\frac1a<1,
\end{equation}
dove si è usato il \cref{lem:potenze-log}(iv). Dal criterio del rapporto otteniamo che
$n^{b}/a^{n}\to0$.

\emph{Terzo limite.} Poniamo
\begin{equation}\label{eq:gerarchia-3}
  a_{n}=\frac{a^{n}}{n!}; \qquad r_{n}=\frac{a_{n+1}}{a_{n}}=\frac{a^{n+1}}{(n+1)!}\cdot\frac{n!}{a^{n}}=\frac{a}{n+1}\to0 .
\end{equation}
Ancora dal criterio del rapporto segue che $a^{n}/n!\to0$.

\emph{Quarto limite.} Poniamo
\begin{equation}\label{eq:gerarchia-4}
\begin{split}
  a_{n}=\frac{n!}{n^{n}}; \qquad
  r_{n}&=\frac{a_{n+1}}{a_{n}}=\frac{(n+1)!}{(n+1)^{n+1}}\cdot\frac{n^{n}}{n!}\\
  &=\frac{n^{n}}{(n+1)^{n}}=\frac{1}{\bigl(\frac{n+1}{n}\bigr)^{n}}\to\frac1e<1,
\end{split}
\end{equation}
e di nuovo, per il criterio del rapporto, la successione $a_{n}=n!/n^{n}$ converge a zero
per $n\to+\infty$.
\end{dimostrazione}

\begin{osservazione}
Il primo limite vale per i logaritmi in qualsiasi base $c>1$, perché
$\log_{c}n=\log n/\log c$ \eqref{eq:log4}. Il risultato si ricorda con la gerarchia
\[
  \log n\ \ll\ n^{b}\ \ll\ a^{n}\ \ll\ n!\ \ll\ n^{n},
\]
dove $x_{n}\ll y_{n}$ significa $x_{n}/y_{n}\to0$. Si noti che la gerarchia non dipende
dalle costanti: anche $\log n$ è trascurabile rispetto a $n^{0{,}001}$, e $n^{1000}$
rispetto a $1{,}001^{n}$. Ma il «sorpasso» può avvenire molto tardi, come mostra
l'\cref{es:n10-su-1-1n}.
\end{osservazione}

\begin{idea}[Il costo degli algoritmi]
La gerarchia degli infiniti è il linguaggio con cui gli informatici classificano il costo
degli algoritmi. Un problema di dimensione $n$ (ad esempio il numero di celle di una
griglia di calcolo, o di elementi di un modello strutturale) si risolve con un numero di
operazioni che può crescere come $\log n$ (ricerca binaria), come $n$ o $n\log n$
(ordinamento), come $n^{2}$ o $n^{3}$ (eliminazione di Gauss su una matrice piena), come
$2^{n}$ (esplorare tutti i sottoinsiemi) o come $n!$ (provare tutti i percorsi del
problema del commesso viaggiatore). Raddoppiare la potenza del calcolatore permette di
affrontare problemi molto più grandi se il costo è polinomiale, ma aggiunge appena una
unità alla dimensione trattabile se il costo è $2^{n}$. La \cref{fig:infiniti} rende
visibile la differenza.
\end{idea}

\begin{python}[La gerarchia degli infiniti in cifre]
import math

print(f"{'n':>4} {'log n':>8} {'n^2':>8} {'2^n':>12}"
      f" {'n!':>12} {'n^n':>12}")
for n in (5, 10, 20, 50, 100):
    print(f"{n:>4} {math.log(n):8.3f} {n**2:8d} {2.0**n:12.3e} "
          f"{float(math.factorial(n)):12.3e} {float(n)**n:12.3e}")
\end{python}
\begin{risultato}
   n    log n      n^2          2^n           n!          n^n
   5    1.609       25    3.200e+01    1.200e+02    3.125e+03
  10    2.303      100    1.024e+03    3.629e+06    1.000e+10
  20    2.996      400    1.049e+06    2.433e+18    1.049e+26
  50    3.912     2500    1.126e+15    3.041e+64    8.882e+84
 100    4.605    10000    1.268e+30   9.333e+157   1.000e+200
\end{risultato}

\begin{figure}[htb]
\centering
\includegraphics[width=0.85\textwidth]{infiniti}
\caption{Le successioni $\log n$, $n^{2}$, $2^{n}$, $n!$, $n^{n}$ in scala logaritmica
(\texttt{python/fig\_infiniti.py}): le curve si separano sempre di più, ciascuna
diventando trascurabile rispetto alla successiva.}
\label{fig:infiniti}
\end{figure}

%% ----------------------------------------------------------------------------
%%  §3.10  Appendice: successioni estratte, Bolzano–Weierstrass, Cauchy
%% ----------------------------------------------------------------------------
\section[Appendice: successioni estratte, teorema di Bolzano--Weierstrass, successioni di Cauchy]%
  {Appendice: successioni estratte, teorema di\\ Bolzano--Weierstrass, successioni di Cauchy}
\label{sec:bw-cauchy}

\subsection{Successioni estratte}

\begin{definizione}[Successione estratta]\index{successione!estratta}
Sia $a_{n}$ una successione di numeri reali e sia $n_{k}$ una successione strettamente
crescente di numeri naturali. La successione $a_{n_{k}}$ definita da
\begin{equation}\label{eq:estratta}
  k\in\N\ \longmapsto\ a_{n_{k}}
\end{equation}
prende il nome di \emph{successione estratta} (o \emph{sottosuccessione}) da $a_{n}$, di
indici $n_{k}$.
\end{definizione}

In altre parole, un'estratta si ottiene «saltando» alcuni termini della successione (anche
infiniti), ma conservando l'ordine. Ad esempio, se $n_{k}=2k$, la successione estratta da
$a_{n}$ di indici $2k$, cioè di indici pari, è
\begin{equation}\label{eq:estratta-pari}
  a_{2},\ a_{4},\ \dots,\ a_{2k},\ \dots
\end{equation}
Se invece $n_{k}=2k-1$, si ottiene l'estratta da $a_{n}$ di indici dispari:
\begin{equation}\label{eq:estratta-dispari}
  a_{1},\ a_{3},\ \dots,\ a_{2k-1},\ \dots
\end{equation}

\begin{lemma}\label{lem:nk-maggiore-k}
Per ogni successione $n_{k}$ strettamente crescente di numeri naturali si ha
\begin{equation}\label{eq:nk-k}
  n_{k}\ge k \qquad\perogni k\in\N .
\end{equation}
\end{lemma}

\begin{schema}
\Ipotesi{$n_{k}\in\N$, \ $n_{k}<n_{k+1}$ per ogni $k$.}
\Tesi{$n_{k}\ge k$ per ogni $k\in\N$.}
\end{schema}

\begin{strategia}
Induzione su $k$: il primo indice è almeno $1$, e a ogni passo l'indice cresce almeno di
un'unità, perché tra due naturali consecutivi non ci sono altri naturali.
\end{strategia}

\begin{dimostrazione}
Per $k=1$ si ha ovviamente $n_{1}\ge1$. Inoltre, supponendo valida la \eqref{eq:nk-k},
proviamo che risulta $n_{k+1}\ge k+1$, da cui, per il principio di induzione, la
\eqref{eq:nk-k} risulterà vera per ogni $k$. Per ipotesi è $n_{k+1}>n_{k}\ge k$, ovvero
$n_{k+1}>k$, e perciò $n_{k+1}\ge k+1$ (\cref{lem:struttura-N}(iii)).
\end{dimostrazione}

Dalla \eqref{eq:nk-k} si ricava facilmente la seguente

\begin{proposizione}\label{prop:estratta-limite}
Se $a_{n}$ ha limite $a$ (finito o infinito), allora ogni estratta $a_{n_{k}}$ ha lo
stesso limite $a$.
\end{proposizione}

\begin{schema}
\Ipotesi{$a_{n}\to a$, con $a\in\R$ oppure $a=\pm\infty$; \ $n_{k}$ strettamente
crescente.}
\Tesi{$a_{n_{k}}\to a$ per $k\to+\infty$.}
\end{schema}

\begin{strategia}
L'estratta prende i suoi termini sempre più avanti nella successione: per il lemma
precedente il $k$-esimo termine dell'estratta ha indice $n_{k}\ge k$. Quindi la soglia
$\nu$ che funziona per la successione funziona anche per l'estratta.
\end{strategia}

\begin{dimostrazione}
Sia $a$ finito. Fissato $\epsilon>0$ esiste $k_{0}$ tale che $\abs{a_{n}-a}<\epsilon$ per
ogni $n>k_{0}$. Se $k>k_{0}$, essendo $n_{k}\ge k$ per il lemma precedente, si ha anche
$n_{k}>k_{0}$, e perciò $\abs{a_{n_{k}}-a}<\epsilon$. Se $a=\pm\infty$ si ragiona allo
stesso modo, con la disuguaglianza $a_{n}>M$ (o $a_{n}<-M$) al posto di
$\abs{a_{n}-a}<\epsilon$.
\end{dimostrazione}

\begin{corollario}\label{cor:estratte-limiti-diversi}
Se una successione ha due estratte con limiti diversi, allora non ha limite.
\end{corollario}

\begin{schema}
\Ipotesi{$a_{n_{k}}\to\ell_{1}$, \ $a_{m_{k}}\to\ell_{2}$ (finiti o infiniti), con
$\ell_{1}\neq\ell_{2}$.}
\Tesi{$a_{n}$ non ha limite, né finito né infinito.}
\end{schema}

\begin{strategia}
È la \cref{prop:estratta-limite} letta «al contrario»: se la successione avesse un
limite, tutte le estratte avrebbero quello stesso limite, e per l'unicità del limite non
potrebbero averne due diversi.
\end{strategia}

\begin{dimostrazione}
Se fosse $a_{n}\to a$ (finito o infinito), per la \cref{prop:estratta-limite} si avrebbe
$a_{n_{k}}\to a$ e $a_{m_{k}}\to a$. Per l'unicità del limite (\cref{teo:unicita-limite}
e, per i limiti infiniti, l'osservazione che segue la definizione di successione regolare
nel \cref{sec:def-limite}) sarebbe $\ell_{1}=a=\ell_{2}$, contro l'ipotesi.
\end{dimostrazione}

Ad esempio, $(-1)^{n}$ non ha limite perché l'estratta di indici pari tende a $1$ e quella
di indici dispari tende a $-1$: è una dimostrazione della
\cref{prop:alterna-non-converge} più rapida di quella data a suo tempo. Il
\cref{teo:lim-geometrica} per $a<-1$ si ottiene allo stesso modo.

\subsection{Il teorema di Bolzano--Weierstrass}

Nel \cref{sec:operazioni-limiti} abbiamo dimostrato che ogni successione convergente è
limitata, cioè esiste $M>0$ tale che $\abs{a_{n}}\le M$ per ogni $n\in\N$. Il viceversa non
sussiste, perché, ad esempio, la successione $a_{n}=(-1)^{n}$ è limitata ma non
convergente. Tuttavia, in essa si possono individuare estratte convergenti (quella dei
termini di posto pari, ad esempio). Ciò non è un caso: sussiste un notevole teorema,
che nello schema dei grandi teoremi di esistenza della \cref{fig:schema-teoremi} segue
immediatamente quello sulle successioni monotòne.

Per enunciarlo e dimostrarlo conviene precisare il linguaggio. Diciamo che un intervallo
$I$ \emph{contiene termini della successione per infiniti indici} se l'insieme
$\{n\in\N : a_{n}\in I\}$ non è limitato superiormente, cioè se per ogni $m\in\N$ esiste
$n>m$ con $a_{n}\in I$. Se $I=I'\cup I''$ e né $I'$ né $I''$ contengono termini per
infiniti indici, allora gli insiemi $\{n : a_{n}\in I'\}$ e $\{n : a_{n}\in I''\}$ sono
limitati superiormente, e quindi lo è anche la loro unione $\{n : a_{n}\in I\}$:
\emph{se $I$ contiene termini per infiniti indici, almeno una delle due parti $I'$, $I''$
ne contiene per infiniti indici}.

\begin{teorema}[di Bolzano--Weierstrass]\label{teo:bolzano-weierstrass}\index{teorema!di Bolzano--Weierstrass}
Sia $a_{n}$ una successione limitata. Allora esiste almeno una sua estratta convergente.
\end{teorema}

\begin{schema}
\Ipotesi{$A\le a_{n}\le B$ per ogni $n\in\N$.}
\Tesi{$\esiste n_{1}<n_{2}<\cdots<n_{k}<\cdots$ e $\ell\in[A,B]$ tali che $a_{n_{k}}\to\ell$.}
\end{schema}

\begin{strategia}
È il metodo di bisezione, già incontrato per approssimare $\sqrt2$, usato questa volta per
\emph{catturare} un valore attorno al quale si addensano infiniti termini della
successione. La successione ha infiniti termini, e sono
tutti nell'intervallo $[A,B]$. Dividiamo l'intervallo a metà (\cref{fig:bolzano-weierstrass}): almeno una delle due metà
contiene termini della successione per infiniti indici (se ciascuna ne contenesse solo un
numero finito, i termini sarebbero in tutto un numero finito). Scegliamo quella metà e
ripetiamo: otteniamo una successione di intervalli incapsulati $[A_{k},B_{k}]$, ciascuno
lungo la metà del precedente e ciascuno contenente termini della successione per infiniti
indici. Gli estremi sinistri $A_{k}$ formano una successione crescente e limitata, che
converge (\cref{teo:monotone}) a un numero $\ell$: qui entra in gioco la completezza.
Infine scegliamo un termine $a_{n_{k}}$ in ciascun intervallo, con indici crescenti
(possibile proprio perché ogni intervallo ne contiene per infiniti indici): l'estratta
così ottenuta è stretta tra $A_{k}$ e $B_{k}$, che tendono entrambi a $\ell$, e i
carabinieri concludono.
\end{strategia}

\begin{dimostrazione}
Per ipotesi la successione $a_{n}$ è limitata; pertanto esistono due costanti $A,B\in\R$,
con $A<B$, tali che
\begin{equation}\label{eq:bw-limitata}
  A\le a_{n}\le B \qquad\perogni n\in\N .
\end{equation}
L'intervallo $[A,B]$ contiene tutti i termini, e quindi termini per infiniti indici
(perché $\N$ non è limitato superiormente, \cref{teo:archimede}).

\emph{Costruzione degli intervalli.} Suddividiamo l'intervallo $[A,B]$ mediante il punto
di mezzo $C=(A+B)/2$ e consideriamo i due intervalli $[A,C]$, $[C,B]$. Per quanto
osservato sopra, uno almeno dei due contiene termini della successione per infiniti
indici; indichiamolo con $[A_{1},B_{1}]$ (se entrambi lo fanno, ne scegliamo
uno qualsiasi). Risulta
\begin{equation}\label{eq:bw-passo1}
  A\le A_{1}, \qquad B_{1}\le B, \qquad B_{1}-A_{1}=\frac{B-A}{2}.
\end{equation}
Suddividiamo l'intervallo $[A_{1},B_{1}]$ mediante il punto di mezzo
$C_{1}=(A_{1}+B_{1})/2$. Per lo stesso motivo almeno uno tra i due intervalli
$[A_{1},C_{1}]$, $[C_{1},B_{1}]$ contiene termini della successione per infiniti indici;
indichiamolo con $[A_{2},B_{2}]$:
\begin{equation}\label{eq:bw-passo2}
  A_{1}\le A_{2}, \qquad B_{2}\le B_{1}, \qquad B_{2}-A_{2}=\frac{B_{1}-A_{1}}{2}=\frac{B-A}{2^{2}}.
\end{equation}
Iterando il procedimento (cioè definendo $[A_{k+1},B_{k+1}]$ a partire da $[A_{k},B_{k}]$,
per induzione) si generano due successioni $A_{k}$, $B_{k}$ $(k\in\N)$ tali che
\begin{gather}
  A\le A_{k}\le A_{k+1}<B_{k+1}\le B_{k}\le B \qquad\perogni k\in\N, \label{eq:bw-incapsulati}\\
  B_{k}-A_{k}=\frac{B-A}{2^{k}} \qquad\perogni k\in\N, \label{eq:bw-lunghezza}
\end{gather}
e inoltre ogni intervallo $[A_{k},B_{k}]$ contiene termini della successione $a_{n}$ per
infiniti indici.

\emph{Scelta degli indici.} L'intervallo $[A_{1},B_{1}]$ contiene termini della
successione $a_{n}$; quindi esiste il primo intero $n_{1}$ tale che
$a_{n_{1}}\in[A_{1},B_{1}]$ (\cref{teo:buon-ordinamento}). Per lo stesso motivo esiste un
primo intero $n_{2}$, fra tutti i numeri naturali più grandi di $n_{1}$, per cui
$a_{n_{2}}\in[A_{2},B_{2}]$: ne esistono, perché $[A_{2},B_{2}]$ contiene termini per
infiniti indici. Iterando, per $k=3,4,5,\dots$ determiniamo una successione
strettamente crescente di interi
\begin{equation}\label{eq:bw-indici}
  n_{1}<n_{2}<n_{3}<\dots<n_{k}<n_{k+1}<\dots
\end{equation}
per cui $a_{n_{k}}\in[A_{k},B_{k}]$ per ogni $k\in\N$. Dato che
$B_{k}-A_{k}=(B-A)/2^{k}$, abbiamo quindi
\begin{equation}\label{eq:bw-carabinieri}
  A_{k}\le a_{n_{k}}\le B_{k}=A_{k}+\frac{B-A}{2^{k}} \qquad\perogni k\in\N .
\end{equation}

\emph{Passaggio al limite.} Per la \eqref{eq:bw-incapsulati} la successione $A_{k}$ è
monotòna (crescente) e limitata; per il teorema sulle successioni monotòne $A_{k}$ ammette
limite finito per $k\to+\infty$. Indichiamo con $\ell\in\R$ il valore di tale limite. Dato
che $(B-A)/2^{k}=(B-A)\bigl(\frac12\bigr)^{k}\to0$ per $k\to+\infty$
(\cref{teo:lim-geometrica}), sia il primo sia l'ultimo membro della
\eqref{eq:bw-carabinieri} convergono a $\ell$. Inoltre $A\le A_{k}\le B$ per ogni $k$, per
la \eqref{eq:bw-incapsulati}, e quindi $A\le\ell\le B$ per il \cref{cor:confronto-limiti}
(applicato confrontando $A_{k}$ con le successioni costanti $A$ e $B$): dunque
$\ell\in[A,B]$. Dal teorema dei carabinieri si ottiene allora la conclusione
\begin{equation}\label{eq:bw-tesi}
  \lim_{k\to+\infty}a_{n_{k}}=\ell. \qedhere
\end{equation}
\end{dimostrazione}

\begin{figure}[htb]
\centering
\begin{tikzpicture}[x=10cm, every node/.style={font=\footnotesize}]
  % termini di una successione limitata (valori di sen n riscalati in [0,1])
  \foreach \row/\A/\B in {0/0/1, 1/0.5/1, 2/0.5/0.75, 3/0.625/0.75, 4/0.625/0.6875} {
    \begin{scope}[yshift={-\row*0.85cm}]
      \draw[grigio!60] (0,0) -- (1,0);
      \fill[blu!18] (\A,-0.12) rectangle (\B,0.12);
      \draw[blu, thick] (\A,-0.12) -- (\A,0.12); \draw[blu, thick] (\B,-0.12) -- (\B,0.12);
      \ifnum\row=0 \node[left] at (0,0) {$[A,B]$}; \else \node[left] at (0,0) {$[A_{\row},B_{\row}]$}; \fi
    \end{scope}}
  \foreach \n in {1,...,60} {
    \pgfmathsetmacro{\v}{0.5+0.5*sin(deg(\n))}
    \fill[rame] (\v,0) circle (0.9pt);}
  \draw[verde!60!black, thick, dashed] (0.6553,0.3) -- (0.6553,-3.8) node[below] {$\ell$};
\end{tikzpicture}
\caption{Il procedimento di bisezione nella dimostrazione del teorema di
Bolzano--Weierstrass. I punti color rame sono i termini (riscalati) di una successione
limitata; a ogni passo si conserva una metà dell'intervallo che contiene termini per
infiniti indici. Gli intervalli si stringono attorno al limite $\ell$ dell'estratta.}
\label{fig:bolzano-weierstrass}
\end{figure}

\begin{osservazione}
Il teorema afferma l'esistenza di un'estratta convergente, ma non ne fornisce, in
generale, una costruzione effettiva: la scelta della metà «buona» a ogni passo richiede
di sapere quale delle due contiene infiniti termini, cosa che nessun calcolatore può
verificare. Inoltre l'estratta non è unica, e successioni diverse possono avere estratte
con limiti diversi: per $\sen n$, ad esempio, si può dimostrare che \emph{ogni} numero di
$[-1,1]$ è limite di qualche estratta. Infine, l'ipotesi di limitatezza è essenziale: la
successione $a_{n}=n$ non ha alcuna estratta convergente, perché ogni estratta diverge a
$+\infty$ (\cref{prop:estratta-limite}).
\end{osservazione}

\subsection{Successioni di Cauchy}

Il teorema sulle successioni monotòne permette di dimostrare che una successione converge
senza conoscerne il limite; ma richiede la monotonia. Il criterio di Cauchy offre un
analogo strumento per successioni qualsiasi.

\begin{definizione}[Successione di Cauchy]\index{successione!di Cauchy}
Sia $a_{n}$ una successione di numeri reali. Si dice che $a_{n}$ è una \emph{successione di
Cauchy} se, per ogni $\epsilon>0$, esiste un indice $\nu$ tale che per $h,k>\nu$ risulti
\begin{equation}\label{eq:cauchy}
  \abs{a_{k}-a_{h}}<\epsilon .
\end{equation}
\end{definizione}

La condizione di Cauchy non menziona alcun limite: chiede soltanto che i termini, da un
certo punto in poi, siano \emph{vicini tra loro}. Dimostriamo in primo luogo la seguente

\begin{proposizione}\label{prop:convergente-cauchy}
Ogni successione convergente è di Cauchy.
\end{proposizione}

\begin{schema}
\Ipotesi{$a_{n}\to a\in\R$.}
\Tesi{$\perogni\epsilon>0\ \esiste\nu:\ \abs{a_{k}-a_{h}}<\epsilon$ per ogni $h,k>\nu$.}
\end{schema}

\begin{strategia}
Se due termini sono entrambi vicini al limite $a$, sono vicini tra loro: la disuguaglianza
triangolare, passando per $a$, somma le due distanze. Per ottenere $\epsilon$ alla fine
si chiede a ciascuna distanza di essere minore di $\epsilon/2$.
\end{strategia}

\begin{dimostrazione}
Se $a_{n}$ converge verso $a$, allora, per ogni $\epsilon>0$, esiste $\nu$ tale che
\begin{equation}\label{eq:conv-cauchy-1}
  \abs{a_{n}-a}<\frac\epsilon2 \qquad\perogni n>\nu .
\end{equation}
Dalla disuguaglianza triangolare segue allora, per $h,k>\nu$:
\begin{equation}\label{eq:conv-cauchy-2}
  \abs{a_{k}-a_{h}}\le\abs{a_{k}-a}+\abs{a-a_{h}}<\frac\epsilon2+\frac\epsilon2=\epsilon. \qedhere
\end{equation}
\end{dimostrazione}

Per dimostrare che, viceversa, ogni successione di Cauchy è convergente, premettiamo alcuni
lemmi.

\begin{lemma}\label{lem:cauchy-limitata}
Una successione di Cauchy è limitata.
\end{lemma}

\begin{schema}
\Ipotesi{$a_{n}$ di Cauchy.}
\Tesi{$\esiste A,B\in\R$: \ $A\le a_{n}\le B$ per ogni $n\in\N$.}
\end{schema}

\begin{strategia}
Con $\epsilon=1$, da un certo indice in poi tutti i termini distano meno di $1$ da un
termine fissato $a_{h_{0}}$: sono quindi definitivamente in un intervallo limitato.
Restano fuori un numero finito di termini, che si aggiungono al calcolo del minimo e del
massimo.
\end{strategia}

\begin{dimostrazione}
Sia $\epsilon=1$; per ipotesi esiste $\nu\in\N$ tale che
\begin{equation}\label{eq:cauchy-lim-1}
  \abs{a_{k}-a_{h}}<1 \qquad\perogni h,k>\nu .
\end{equation}
Fissiamo un indice $h_{0}>\nu$. Allora, dalla \eqref{eq:cauchy-lim-1}, per le proprietà
del valore assoluto, segue
\begin{equation}\label{eq:cauchy-lim-2}
  a_{h_{0}}-1<a_{k}<a_{h_{0}}+1 \qquad\perogni k>\nu .
\end{equation}
Posto
\[
  A=\min\{a_{1},\dots,a_{\nu},a_{h_{0}}-1\}, \qquad B=\max\{a_{1},\dots,a_{\nu},a_{h_{0}}+1\},
\]
evidentemente risulta $A\le a_{k}\le B$ per ogni $k\in\N$, e perciò la successione è
limitata. (È lo stesso schema del \cref{teo:convergente-limitata}: controllo definitivo
più un numero finito di eccezioni.)
\end{dimostrazione}

\begin{lemma}\label{lem:cauchy-estratta}
Se una successione di Cauchy $a_{n}$ contiene un'estratta $a_{n_{k}}$ convergente verso
$\ell$, allora anche $a_{n}$ converge verso $\ell$.
\end{lemma}

\begin{schema}
\Ipotesi{$a_{n}$ di Cauchy; \ $a_{n_{k}}\to\ell\in\R$.}
\Tesi{$a_{n}\to\ell$.}
\end{schema}

\begin{strategia}
Per stimare $\abs{a_{n}-\ell}$ si passa per un termine dell'estratta abbastanza avanzato,
$a_{n_{k_{0}}}$: questo è vicino a $\ell$ perché l'estratta converge, ed è vicino ad
$a_{n}$ perché la successione è di Cauchy. La disuguaglianza triangolare, con il solito
$\epsilon/2+\epsilon/2$, conclude.
\end{strategia}

\begin{dimostrazione}
Fissato $\epsilon>0$, sia $\nu\in\N$ tale che
\begin{equation}\label{eq:cauchy-estr-1}
  \abs{a_{k}-a_{h}}<\frac\epsilon2 \qquad\perogni h,k>\nu .
\end{equation}
Sia inoltre $k_{0}>\nu$ tale che
\begin{equation}\label{eq:cauchy-estr-2}
  \abs{a_{n_{k}}-\ell}<\frac\epsilon2 \qquad\perogni k\ge k_{0}.
\end{equation}
Poiché si ha $n_{k_{0}}\ge k_{0}>\nu$ (\cref{lem:nk-maggiore-k}), per ogni $n>\nu$ risulta
\begin{equation}\label{eq:cauchy-estr-3}
  \abs{a_{n}-\ell}\le\abs{a_{n}-a_{n_{k_{0}}}}+\abs{a_{n_{k_{0}}}-\ell}<\frac\epsilon2+\frac\epsilon2=\epsilon. \qedhere
\end{equation}
\end{dimostrazione}

Combinando i risultati precedenti si dimostra il seguente

\begin{teorema}[Criterio di convergenza di Cauchy]\label{teo:criterio-cauchy}\index{criterio di Cauchy!per le successioni}
Una successione $a_{n}$ di numeri reali è convergente se e solo se è di Cauchy.
\end{teorema}

\begin{schema}
\Ipotesi{$a_{n}$ successione di numeri reali.}
\Tesi{$a_{n}$ converge in $\R$ \ $\sse$ \ $a_{n}$ è di Cauchy, cioè per ogni $\epsilon>0$
esiste $\nu$ tale che $\abs{a_{k}-a_{h}}<\epsilon$ per ogni $h,k>\nu$.}
\end{schema}

\begin{strategia}
Una delle due implicazioni è la \cref{prop:convergente-cauchy}. Per l'altra, il problema è
che non conosciamo il candidato limite; lo facciamo produrre al teorema di
Bolzano--Weierstrass. Una successione di Cauchy è limitata (\cref{lem:cauchy-limitata}),
quindi ha un'estratta convergente a qualche $\ell$; e poiché i termini di una successione
di Cauchy sono definitivamente vicini tra loro, se alcuni di essi si avvicinano a $\ell$,
vi si avvicinano tutti (\cref{lem:cauchy-estratta}).
\end{strategia}

\begin{dimostrazione}
Con la \cref{prop:convergente-cauchy} abbiamo già provato che le successioni convergenti
sono di Cauchy. Viceversa, se $a_{n}$ è di Cauchy, per il \cref{lem:cauchy-limitata} è
anche limitata. In base al teorema di Bolzano--Weierstrass (\cref{teo:bolzano-weierstrass})
$a_{n}$ ammette una successione estratta $a_{n_{k}}$ convergente. Per il
\cref{lem:cauchy-estratta}, $a_{n}$ è convergente.
\end{dimostrazione}

\begin{osservazione}[Cauchy e la completezza]
Nel campo dei razionali il criterio di Cauchy è falso: la successione dei troncamenti
decimali di $\sqrt2$,
\[
  1;\quad 1{,}4;\quad 1{,}41;\quad 1{,}414;\quad 1{,}4142;\quad\dots,
\]
è formata da numeri razionali ed è di Cauchy (due termini di indice maggiore di $k$
distano meno di $10^{-k}$), ma non converge ad alcun numero razionale, perché il suo
limite in $\R$ è $\sqrt2\notin\Q$. La dimostrazione del criterio usa infatti, attraverso
il teorema di Bolzano--Weierstrass e quello sulle successioni monotòne, l'assioma di
completezza. Si può anzi dimostrare che, in un campo ordinato archimedeo, il criterio di
Cauchy è \emph{equivalente} all'assioma di completezza \cite{giusti}. Lo strumento è
ancora la bisezione. Dati $A$ e $B$ come nell'assioma, si sceglie $a_{0}\in A$ e
$b_{0}\in B$ e si dimezza ripetutamente l'intervallo $[a_{0},b_{0}]$, conservando a ogni
passo una metà il cui estremo destro è un maggiorante di $A$ e il cui estremo sinistro non
supera qualche elemento di $A$. Gli estremi formano due successioni di Cauchy, e il loro
limite comune separa $A$ da $B$. La proprietà di Archimede serve a garantire che le
lunghezze $(b_{0}-a_{0})/2^{k}$ tendano a zero.
\end{osservazione}

\begin{python}[Una successione di Cauchy di razionali che non converge in $\Q$]
# Troncamenti decimali di sqrt(2): sono razionali, distano sempre
# meno, ma il loro limite sqrt(2) non è razionale.
from fractions import Fraction as F
from decimal import Decimal, getcontext
getcontext().prec = 40
cifre = str(Decimal(2).sqrt())            # "1.414213562373095..."

tronc = [F(cifre[:2 + k]) for k in range(1, 12)]   # 1.4, 1.41, ...
print("k               q_k   max |q_j - q_k|, j > k   q_k^2 = 2 ?")
for k in range(1, 8):
    q = tronc[k - 1]
    dist = max(abs(tronc[j] - q) for j in range(k, len(tronc)))
    print(f"{k}  {str(q):>16s}   {float(dist):22.1e}   {q * q == 2}")
\end{python}
\begin{risultato}
k               q_k   max |q_j - q_k|, j > k   q_k^2 = 2 ?
1               7/5                  1.4e-02   False
2           141/100                  4.2e-03   False
3           707/500                  2.1e-04   False
4         7071/5000                  1.4e-05   False
5     141421/100000                  3.6e-06   False
6   1414213/1000000                  5.6e-07   False
7   2828427/2000000                  6.2e-08   False
\end{risultato}

\begin{idea}[I numeri reali come completamento dei razionali]
Il programma mostra una successione di Cauchy di razionali che «vorrebbe» convergere ma
non trova in $\Q$ il suo limite. Charles Méray (1869) e, indipendentemente, Georg Cantor
(1872) rovesciarono il problema: se il limite manca, lo si \emph{definisce} a partire
dalla successione stessa. Due successioni di Cauchy di razionali $a_{n}$ e $b_{n}$ si
dicono equivalenti se $a_{n}-b_{n}\to0$: ad esempio i troncamenti decimali di $\sqrt2$ e
i termini $2,\frac32,\frac{17}{12},\frac{577}{408},\dots$ del metodo di Erone
(\cref{es:erone}). Un numero reale è, per definizione, una classe di successioni
equivalenti, e il razionale $q$ si identifica con la classe della successione costante
$q$. Somma e prodotto si definiscono termine a termine; un numero è $\ge0$ se ha un
rappresentante definitivamente $\ge0$. Con un lavoro lungo ma elementare (la parte più
delicata riguarda l'ordinamento) si verifica che l'insieme così costruito soddisfa tutti
gli assiomi del \cref{sec:assiomi}, compreso quello di completezza \cite{giusti}. Anche
$e$ trova posto: è la classe della successione di razionali $(1+\frac1n)^{n}$, alla
quale è equivalente quella delle somme $\sum_{k=0}^{n}\frac1{k!}$
(\cref{oss:e-due-vie}). È la costruzione seguita in \cite{tao1}, alternativa alle
sezioni di Dedekind ricordate nella nota storica del \S\ref{sec:assiomi}.
\end{idea}

\begin{notastorica}[Bolzano, Cauchy, Weierstrass]
L'idea della dimostrazione del teorema che chiamiamo «di Bolzano--Weierstrass», la
bisezione, compare già nel saggio di Bernard Bolzano del 1817 ricordato nel
\cref{sec:sup}, dove serve a dimostrare l'esistenza dell'estremo superiore e, da questa,
il teorema dei valori intermedi. Nello stesso lavoro Bolzano enunciò
anche, quattro anni prima di Cauchy, il criterio che porta il nome di quest'ultimo, e
cercò di dimostrarlo. Il suo argomento si riduceva però a osservare che l'ipotesi di un
limite «non contiene nulla di impossibile», perché, se il limite c'è, il suo valore si
può determinare con tutta la precisione che si vuole: così si descrive il limite, ma non
se ne prova l'esistenza. In mancanza di una teoria dei numeri reali la dimostrazione non
poteva che essere illusoria \cite{giusti}. Il saggio, d'altra parte, rimase quasi
sconosciuto, e fu Karl Weierstrass a riscoprire il risultato e a
farne, nelle lezioni berlinesi degli anni Sessanta dell'Ottocento, uno dei pilastri
dell'Analisi. Augustin-Louis Cauchy enunciò il suo criterio nel \emph{Cours d'analyse}
(1821), ma ne considerò ovvia la parte «difficile» (che una successione di Cauchy
converge): ovvia lo è, se si ha in mente la retta geometrica; dimostrarla richiede una
teoria dei numeri reali, che né a Cauchy né a Bolzano era disponibile. Fu proprio per
colmare questa lacuna che, cinquant'anni dopo, Méray, Dedekind e Cantor costruirono i
numeri reali (si veda il riquadro precedente).
\end{notastorica}

\subsection{Approfondimento: massimo e minimo limite}
\label{subsec:massimo-minimo-limite}
\index{massimo limite}\index{minimo limite}

Una successione limitata può non avere limite, come $(-1)^{n}$. Il teorema di
Bolzano--Weierstrass garantisce che abbia estratte convergenti, ma non dice quali, e i
loro limiti possono essere diversi. Vi sono però due numeri, associati a ogni
successione limitata, che esistono \emph{sempre} e che si riducono al limite quando
questo c'è: il \emph{massimo} e il \emph{minimo limite} \cite{giusti}. L'idea è guardare
la successione «da un certo punto in poi», cioè le sue \emph{code}, e misurarne il
«tetto» e il «pavimento».

Sia $a_{n}$ una successione limitata, diciamo $A\le a_{n}\le B$ per ogni $n$. Per ogni
$n\in\N$ la coda $\{a_{k} : k\ge n\}$ è un insieme non vuoto e limitato: per il
\cref{teo:esistenza-sup} ha estremo superiore ed estremo inferiore reali, che indichiamo
con
\begin{equation}\label{eq:code-limsup}
  s_{n}=\sup_{k\ge n}a_{k}, \qquad i_{n}=\inf_{k\ge n}a_{k} .
\end{equation}
Poiché $a_{n}$ appartiene alla propria coda, si ha $A\le i_{n}\le a_{n}\le s_{n}\le B$.
Passando da $n$ a $n+1$, inoltre, la coda perde il termine $a_{n}$. Il numero $s_{n}$,
che maggiora la coda di indice $n$, maggiora a maggior ragione quella di indice $n+1$;
e poiché $s_{n+1}$ è il \emph{minimo} dei maggioranti di quest'ultima, $s_{n+1}\le s_{n}$.
Allo stesso modo $i_{n+1}\ge i_{n}$. La successione $s_{n}$ è dunque decrescente, la
$i_{n}$ crescente, ed entrambe sono limitate: per il \cref{teo:monotone} convergono.

\begin{definizione}[Massimo e minimo limite]\label{def:limsup}
Sia $a_{n}$ una successione limitata, e siano $s_{n}$ e $i_{n}$ definite dalla
\eqref{eq:code-limsup}. I numeri reali
\begin{equation}\label{eq:limsup-def}
\begin{aligned}
  \limsup_{n\to+\infty}a_{n}&=\lim_{n\to+\infty}s_{n}=\inf_{n\in\N}s_{n},\\
  \liminf_{n\to+\infty}a_{n}&=\lim_{n\to+\infty}i_{n}=\sup_{n\in\N}i_{n}
\end{aligned}
\end{equation}
si chiamano, rispettivamente, \emph{massimo limite} e \emph{minimo limite} (o
\emph{limite superiore} e \emph{limite inferiore}) della successione $a_{n}$. Si usano
anche i simboli $\overline{\lim}\,a_{n}$ e $\underline{\lim}\,a_{n}$.
\end{definizione}

Il massimo limite è dunque il tetto della successione «visto da lontano»: $s_{n}$ è il
più basso dei livelli che i termini da $a_{n}$ in poi non superano, e
$\limsup a_{n}$ è il valore verso cui questi livelli si abbassano al crescere di $n$.

\begin{esempio}[Una successione alternante]\label{es:limsup-alternante}
Sia $a_{n}=(-1)^{n}\bigl(1+\frac1n\bigr)$. I termini di indice pari, $1+\frac1n$,
decrescono verso $1$; quelli di indice dispari, $-1-\frac1n$, crescono verso $-1$
(\cref{fig:limsup-code}). Il più grande termine della coda di indice $n$ è il primo
termine di indice pari che vi compare: $s_{n}=1+\frac1n$ se $n$ è pari,
$s_{n}=1+\frac1{n+1}$ se $n$ è dispari. Analogamente $i_{n}=-1-\frac1n$ se $n$ è dispari
e $i_{n}=-1-\frac1{n+1}$ se $n$ è pari. Quindi
\[
  \limsup_{n\to+\infty}a_{n}=1, \qquad \liminf_{n\to+\infty}a_{n}=-1,
\]
mentre $\sup_{n}a_{n}=a_{2}=\frac32$ e $\inf_{n}a_{n}=a_{1}=-2$. Gli estremi
superiore e inferiore dipendono dai primi termini; il massimo e il minimo limite no,
perché modificando un numero finito di termini le code da un certo indice in poi non
cambiano.
\end{esempio}

\begin{figure}[htb]
\centering
\begin{tikzpicture}
\begin{axis}[width=0.92\textwidth, height=5.6cm, xmin=0, xmax=25, ymin=-2.25, ymax=1.85,
  xtick={5,10,15,20}, ytick={-2,-1,1,1.5}, yticklabels={$-2$,$\lambda=-1$,$L=1$,$\frac32$},
  xlabel={$n$}, legend columns=3,
  legend style={font=\scriptsize, draw=none, fill=none, at={(0.5,0.69)}, anchor=center}]
  \addplot[tratteggio, domain=0:25, forget plot] {1};
  \addplot[tratteggio, domain=0:25, forget plot] {-1};
  \addplot[only marks, mark=*, mark size=1.4pt, blu, samples at={1,...,24}] {(-1)^x*(1+1/x)};
  \addplot[rame, thick, const plot, samples at={1,...,25}] {1+1/(2*ceil(x/2))};
  \addplot[verde, thick, const plot, samples at={1,...,25}] {-1-1/(2*floor(x/2)+1)};
  \legend{$a_{n}$, $s_{n}=\sup_{k\ge n}a_{k}$, $i_{n}=\inf_{k\ge n}a_{k}$}
\end{axis}
\end{tikzpicture}
\caption{La successione $a_{n}=(-1)^{n}\bigl(1+\frac1n\bigr)$
dell'\cref{es:limsup-alternante}, con i «tetti» $s_{n}$ (in rame) e i «pavimenti»
$i_{n}$ (in verde) delle sue code, disegnati a gradini. I tetti scendono verso
$L=\limsup a_{n}=1$, i pavimenti salgono verso $\lambda=\liminf a_{n}=-1$; l'estremo
superiore $\frac32$ e l'estremo inferiore $-2$ sono invece fissati dai primi termini.}
\label{fig:limsup-code}
\end{figure}

\begin{idea}[Il picco a regime di una risposta oscillante]
In una prova di vibrazione, o nella simulazione della risposta di un'ala a una raffica
periodica, si registra a istanti successivi una grandezza, ad esempio lo spostamento
della punta alare. La successione dei campioni contiene un \emph{transitorio}, spesso
con un picco iniziale elevato, seguito da un'oscillazione a regime. L'estremo superiore
dei campioni risente del picco iniziale; il massimo limite no: misura il picco \emph{a
regime}, quello che si ripete a ogni ciclo e che conta, ad esempio, per la fatica del
materiale. Allo stesso modo il minimo limite misura la valle a regime, e la semidifferenza
$\frac12(L-\lambda)$ l'ampiezza dell'oscillazione stabilizzata. Nell'\cref{es:limsup-alternante}
il valore $\sup a_{n}=\frac32$ è dovuto al secondo termine, mentre a regime la
successione oscilla tra valori sempre più vicini a $-1$ e a $1$.
\end{idea}

\begin{proposizione}[Proprietà del massimo e del minimo limite]\label{prop:limsup-proprieta}
Sia $a_{n}$ una successione limitata, e siano $L=\limsup a_{n}$ e $\lambda=\liminf a_{n}$.
Allora:
\begin{enumerate}
  \item $\inf_{n}a_{n}\le\lambda\le L\le\sup_{n}a_{n}$;
  \item per ogni $\epsilon>0$ si ha $a_{n}<L+\epsilon$ definitivamente, e
  $a_{n}>L-\epsilon$ per infiniti indici;
  \item per ogni $\epsilon>0$ si ha $a_{n}>\lambda-\epsilon$ definitivamente, e
  $a_{n}<\lambda+\epsilon$ per infiniti indici.
\end{enumerate}
\end{proposizione}

\begin{schema}
\Ipotesi{$a_{n}$ limitata; \ $L=\limsup a_{n}$, \ $\lambda=\liminf a_{n}$; \ $\epsilon>0$.}
\Tesi{(i) $\inf a_{n}\le\lambda\le L\le\sup a_{n}$; \ (ii) $a_{n}<L+\epsilon$
definitivamente, e $a_{n}>L-\epsilon$ per infiniti indici; \ (iii) le affermazioni
simmetriche per $\lambda$.}
\end{schema}

\begin{strategia}
La (i) si ottiene passando al limite nella catena $i_{1}\le i_{n}\le s_{n}\le s_{1}$. La
(ii) traduce le caratterizzazioni dell'estremo inferiore e dell'estremo superiore.
Poiché $L$ è l'estremo inferiore dei tetti $s_{n}$, qualche tetto $s_{\nu}$ è più basso
di $L+\epsilon$, e da $\nu$ in poi i termini non lo superano. D'altra parte ogni tetto
$s_{m}$ è almeno $L$: quindi in ogni coda qualche termine supera $L-\epsilon$, e ciò
accade oltre qualunque indice $m$. La (iii) è simmetrica.
\end{strategia}

\begin{dimostrazione}
(i) Per ogni $n$ si ha $i_{1}\le i_{n}\le s_{n}\le s_{1}$, dove $i_{1}=\inf_{k}a_{k}$ e
$s_{1}=\sup_{k}a_{k}$ sono gli estremi di tutta la successione. Passando al limite per
$n\to+\infty$ (\cref{cor:confronto-limiti}, con il confronto anche con le successioni
costanti $i_{1}$ e $s_{1}$) si ottiene la (i).

(ii) Sia $\epsilon>0$. Poiché $L=\inf_{n}s_{n}$, per la \eqref{eq:caratt-inf} esiste
$\nu\in\N$ con $s_{\nu}<L+\epsilon$. Per ogni $n>\nu$ il termine $a_{n}$ appartiene alla
coda di indice $\nu$, e quindi $a_{n}\le s_{\nu}<L+\epsilon$: la prima affermazione è
provata. Per la seconda fissiamo $m\in\N$. Poiché $L$ è un minorante dei tetti,
$s_{m+1}\ge L>L-\epsilon$; per la caratterizzazione dell'estremo superiore
(\cref{prop:caratt-sup}) applicata alla coda di indice $m+1$, esiste $k>m$ con
$a_{k}>L-\epsilon$. Dato che $m$ è arbitrario, l'insieme degli indici $k$ con
$a_{k}>L-\epsilon$ non è limitato superiormente: la disuguaglianza vale per infiniti
indici.

(iii) Si ragiona allo stesso modo, scambiando i ruoli dell'estremo superiore e
dell'estremo inferiore.
\end{dimostrazione}

\begin{osservazione}[La definizione di \cite{giusti}]\label{oss:limsup-giusti}
Diciamo \emph{maggiorante definitivo} della successione un numero $M$ tale che
$a_{n}\le M$ definitivamente. La (ii) dice che ogni $M>L$ è un maggiorante definitivo
(con $\epsilon=M-L$), mentre nessun $M<L$ lo è (con $\epsilon=L-M$ si ha $a_{n}>M$ per
infiniti indici). Quindi $L$ è un minorante dell'insieme dei maggioranti definitivi, e
$L+\epsilon$ vi appartiene per ogni $\epsilon>0$: per la \eqref{eq:caratt-inf}, $L$ ne è
l'estremo inferiore. È la definizione di massimo limite adottata in \cite{giusti}, che è
dunque equivalente alla nostra; allo stesso modo il minimo limite è l'estremo superiore
dei minoranti definitivi. Lo stesso ragionamento mostra che le due condizioni della (ii)
\emph{caratterizzano} $L$: un numero che le soddisfa entrambe è l'estremo inferiore dei
maggioranti definitivi, e quindi coincide con $L$.
\end{osservazione}

\begin{teorema}[Massimo e minimo limite ed estratte]\label{teo:limsup-estratte}
Sia $a_{n}$ una successione limitata, e siano $L=\limsup a_{n}$ e $\lambda=\liminf a_{n}$.
\begin{enumerate}
  \item Esistono un'estratta di $a_{n}$ che converge a $L$ e un'estratta che converge
  a $\lambda$.
  \item Se un'estratta $a_{n_{k}}$ converge a un numero $\ell$, allora
  $\lambda\le\ell\le L$.
\end{enumerate}
In altre parole, $L$ e $\lambda$ sono il più grande e il più piccolo dei limiti delle
estratte convergenti di $a_{n}$.
\end{teorema}

\begin{schema}
\Ipotesi{$a_{n}$ limitata; \ $L=\limsup a_{n}$, \ $\lambda=\liminf a_{n}$.}
\Tesi{(i) esistono estratte $a_{n_{k}}\to L$ e $a_{m_{k}}\to\lambda$; \ (ii) se
$a_{n_{k}}\to\ell$, allora $\lambda\le\ell\le L$.}
\end{schema}

\begin{strategia}
Per la (i) costruiamo l'estratta un termine alla volta, come nella dimostrazione del
teorema di Bolzano--Weierstrass, ma senza bisezione. Per la
\cref{prop:limsup-proprieta}, per ogni $k$ gli indici $n$ con
$L-\frac1k<a_{n}<L+\frac1k$ sono infiniti, e se ne può sempre scegliere uno maggiore del
precedente. Per la (ii) basta osservare che $a_{n_{k}}$, avendo indice $n_{k}\ge k$, sta
nella coda di indice $k$: è quindi compreso tra $i_{k}$ e $s_{k}$, e si passa al limite.
\end{strategia}

\begin{dimostrazione}
(i) Sia $k\in\N$. Per la \cref{prop:limsup-proprieta}(ii) con $\epsilon=\frac1k$ esiste
$\nu_{k}$ tale che $a_{n}<L+\frac1k$ per ogni $n>\nu_{k}$, mentre $a_{n}>L-\frac1k$ per
infiniti indici, cioè per indici maggiori di qualunque numero fissato. Ne segue che, per
ogni $m$, esiste $n>\max\{m,\nu_{k}\}$ tale che
\begin{equation}\label{eq:limsup-scelta}
  L-\frac1k<a_{n}<L+\frac1k .
\end{equation}
Scegliamo allora un indice $n_{1}$ che soddisfi la \eqref{eq:limsup-scelta} con $k=1$,
poi un indice $n_{2}>n_{1}$ che la soddisfi con $k=2$, e così via. Per induzione si
ottiene una successione strettamente crescente di indici $n_{k}$ con
$L-\frac1k<a_{n_{k}}<L+\frac1k$ per ogni $k$. Poiché $L\pm\frac1k\to L$, per il teorema
dei carabinieri $a_{n_{k}}\to L$. In modo analogo, usando la (iii), si costruisce
un'estratta convergente a $\lambda$.

(ii) Per il \cref{lem:nk-maggiore-k} si ha $n_{k}\ge k$, quindi $a_{n_{k}}$ appartiene
alla coda di indice $k$, e
\[
  i_{k}\le a_{n_{k}}\le s_{k} \qquad\perogni k\in\N .
\]
Per $k\to+\infty$ si ha $i_{k}\to\lambda$, $s_{k}\to L$ e $a_{n_{k}}\to\ell$; per il
\cref{cor:confronto-limiti} risulta $\lambda\le\ell\le L$.
\end{dimostrazione}

\begin{osservazione}[Un'altra dimostrazione del teorema di Bolzano--Weierstrass]\label{oss:bw-limsup}
Il punto (i) del \cref{teo:limsup-estratte} contiene una nuova dimostrazione del
\cref{teo:bolzano-weierstrass}: una successione limitata ha un'estratta convergente, ad
esempio quella che converge al massimo limite. È la via seguita da \cite{giusti}, e non è
circolare: abbiamo usato l'esistenza degli estremi superiore e inferiore, il teorema
sulle successioni monotòne e la caratterizzazione dell'estremo superiore, non il teorema
di Bolzano--Weierstrass. Il confronto tra le due vie è istruttivo. La bisezione cattura
un punto attorno al quale si addensano infiniti termini, senza dire quale; i tetti
$s_{n}$ indicano invece un punto preciso, il più grande possibile. In entrambe la
completezza entra nello stesso modo, attraverso il limite di una successione monotòna:
gli estremi $A_{k}$ degli intervalli dimezzati nella prima, i tetti $s_{n}$ nella
seconda.
\end{osservazione}

\begin{teorema}[Limite, massimo e minimo limite]\label{teo:limsup-limite}
Sia $a_{n}$ una successione limitata e sia $\ell\in\R$. Allora
\begin{equation}\label{eq:limsup-limite}
  \lim_{n\to+\infty}a_{n}=\ell \quad\sse\quad
  \limsup_{n\to+\infty}a_{n}=\liminf_{n\to+\infty}a_{n}=\ell .
\end{equation}
\end{teorema}

\begin{schema}
\Ipotesi{$a_{n}$ limitata; \ $\ell\in\R$.}
\Tesi{$a_{n}\to\ell$ \ $\sse$ \ $\limsup a_{n}=\liminf a_{n}=\ell$.}
\end{schema}

\begin{strategia}
Se la successione converge, ogni sua estratta converge allo stesso limite; ma per il
teorema precedente ci sono estratte che tendono a $L$ e a $\lambda$, che sono quindi
uguali a $\ell$. Viceversa, ogni termine sta tra il pavimento e il tetto della propria
coda, $i_{n}\le a_{n}\le s_{n}$: se pavimento e tetto tendono allo stesso limite, i
carabinieri concludono.
\end{strategia}

\begin{dimostrazione}
($\implica$) Se $a_{n}\to\ell$, per la \cref{prop:estratta-limite} ogni estratta di
$a_{n}$ tende a $\ell$. Per il \cref{teo:limsup-estratte}(i) esiste un'estratta che
tende a $L=\limsup a_{n}$, e per l'unicità del limite (\cref{teo:unicita-limite}) è
$L=\ell$. Allo stesso modo $\lambda=\liminf a_{n}=\ell$.

($\Leftarrow$) Per ogni $n$ si ha $i_{n}\le a_{n}\le s_{n}$, e per ipotesi
$i_{n}\to\lambda=\ell$ e $s_{n}\to L=\ell$. Per il teorema dei carabinieri
(\cref{teo:carabinieri}) $a_{n}\to\ell$.
\end{dimostrazione}

\begin{osservazione}\label{oss:limsup-illimitate}
Poiché ogni successione convergente è limitata (\cref{teo:convergente-limitata}), il
teorema si legge anche così: \emph{una successione converge se e solo se è limitata e
il suo massimo limite coincide con il minimo limite}. La differenza $L-\lambda\ge0$
misura quanto la successione «oscilla» all'infinito. Per la successione
dell'\cref{es:limsup-alternante} vale $2$, e la successione non converge; per
$(-1)^{n}/n$ si ha $0<s_{n}\le\frac1n$ e $-\frac1n\le i_{n}<0$, quindi $L=\lambda=0$, che
è il limite. Con le convenzioni $\sup=+\infty$ e $\inf=-\infty$ per gli insiemi non
limitati, le definizioni si estendono a successioni qualsiasi. Ad esempio, per
$a_{n}=(-1)^{n}n$ ogni coda è illimitata sia superiormente sia inferiormente, e si pone
$\limsup a_{n}=+\infty$, $\liminf a_{n}=-\infty$. In questo senso il massimo e il minimo
limite esistono sempre, anche quando il limite non esiste.
\end{osservazione}

%% ----------------------------------------------------------------------------
%%  §3.11  Esercizi di riepilogo
%% ----------------------------------------------------------------------------
\section{Esercizi di riepilogo}
\label{sec:esercizi-cap3}

Come nei capitoli precedenti, ogni soluzione è accompagnata da un commento sulle
\emph{scelte logiche} che la guidano.

\subsection*{La definizione di limite}

\begin{esercizio}
Verificare con la definizione che $\displaystyle\lim_{n\to+\infty}\frac{3n-1}{n+2}=3$,
determinando esplicitamente un indice $\nu(\epsilon)$.
\begin{soluzione}
Si ha $\bigl|\frac{3n-1}{n+2}-3\bigr|=\bigl|\frac{3n-1-3n-6}{n+2}\bigr|=\frac7{n+2}$, e
$\frac7{n+2}<\epsilon$ equivale a $n>\frac7\epsilon-2$: basta $\nu=\frac7\epsilon-2$.
\emph{Scelta logica:} per verificare un limite con la definizione si parte sempre da
$\abs{a_{n}-a}$, lo si semplifica, e si risolve la disuguaglianza $\abs{a_{n}-a}<\epsilon$
rispetto a $n$. Non serve trovare il $\nu$ \emph{più piccolo}: basta un $\nu$ che funzioni,
e spesso conviene maggiorare $\abs{a_{n}-a}$ con qualcosa di più semplice prima di
risolvere.
\end{soluzione}
\end{esercizio}

\begin{esercizio}
Dimostrare che se $a_{n}\to0$ allora anche la media aritmetica
$m_{n}=\frac{a_{1}+\dots+a_{n}}{n}$ tende a zero.
\begin{soluzione}
Fissato $\epsilon>0$, sia $\nu\in\N$ tale che $\abs{a_{n}}<\epsilon/2$ per $n>\nu$, e sia
$S=\abs{a_{1}}+\dots+\abs{a_{\nu}}$ (un numero fisso). Per $n>\nu$,
\[
  \abs{m_{n}}\le\frac{S}{n}+\frac{\abs{a_{\nu+1}}+\dots+\abs{a_{n}}}{n}
  <\frac Sn+\frac{n-\nu}{n}\cdot\frac\epsilon2<\frac Sn+\frac\epsilon2,
\]
e $\frac Sn<\frac\epsilon2$ per $n>2S/\epsilon$. Quindi $\abs{m_{n}}<\epsilon$ per
$n>\max\{\nu,2S/\epsilon\}$. Il viceversa è falso: per $a_{n}=(-1)^{n}$ la media tende a
zero ma $a_{n}$ non converge. \emph{Scelta logica:} si spezza la somma in una «testa»
finita, che pesa sempre meno perché divisa per $n$, e una «coda» in cui i termini sono
piccoli; a ciascuna parte si assegna metà della tolleranza, come nel trucco
dell'$\epsilon/2$ (\cref{es:sup-somma}).
\end{soluzione}
\end{esercizio}

\subsection*{Calcolo di limiti}

\begin{esercizio}
Calcolare $\displaystyle\lim_{n\to+\infty}\frac{2^{n}+3^{n}}{3^{n+1}-2^{n}}$.
\begin{soluzione}
Dividendo numeratore e denominatore per $3^{n}$, il termine che cresce più rapidamente:
$\frac{(2/3)^{n}+1}{3-(2/3)^{n}}\to\frac{0+1}{3-0}=\frac13$, per il
\cref{teo:lim-geometrica} con $a=2/3$. \emph{Scelta logica:} in una forma $\infty/\infty$ si
divide per l'infinito di ordine più alto, qui $3^{n}$ perché $2^{n}/3^{n}=(2/3)^{n}\to0$;
nei casi meno evidenti la gerarchia del \cref{sec:infiniti} dice quale sia.
\end{soluzione}
\end{esercizio}

\begin{esercizio}
Calcolare $\displaystyle\lim_{n\to+\infty}\bigl(\sqrt{n^{2}+n}-n\bigr)$.
\begin{soluzione}
È una forma $\infty-\infty$. Razionalizzando,
$\sqrt{n^{2}+n}-n=\frac{n}{\sqrt{n^{2}+n}+n}=\frac{1}{\sqrt{1+1/n}+1}$. Poiché
$1\le\sqrt{1+1/n}\le1+\frac1n$ (la seconda disuguaglianza perché
$(1+\frac1n)^{2}\ge1+\frac1n$), per i carabinieri $\sqrt{1+1/n}\to1$, e il limite vale
$\frac12$. \emph{Scelta logica:} come nell'\cref{es:razionalizzare}, si moltiplica per il
«coniugato» per trasformare la differenza di due infiniti in un quoziente; il risultato,
$\frac12$, mostra che una forma $\infty-\infty$ può avere un limite finito non nullo.
\end{soluzione}
\end{esercizio}

\begin{esercizio}\label{es:radice-continua}
Dimostrare che se $c_{n}\ge0$ e $c_{n}\to c$, allora $\sqrt{c_{n}}\to\sqrt c$. Dedurne
che $\bigl(1+\frac1{2n}\bigr)^{n}\to\sqrt e$.
\begin{soluzione}
Per il \cref{cor:limite-non-negativo} $c\ge0$. Se $c>0$:
$\abs{\sqrt{c_{n}}-\sqrt c}=\frac{\abs{c_{n}-c}}{\sqrt{c_{n}}+\sqrt c}\le\frac{\abs{c_{n}-c}}{\sqrt c}$,
che tende a zero. Se $c=0$: $\sqrt{c_{n}}<\epsilon$ se e solo se $c_{n}<\epsilon^{2}$, che è
vero definitivamente. Infine
$\bigl(1+\frac1{2n}\bigr)^{n}=\sqrt{\bigl(1+\frac1{2n}\bigr)^{2n}}$, e
$\bigl(1+\frac1{2n}\bigr)^{2n}$ è l'estratta di indici pari della successione che definisce
$e$, quindi tende ad $e$ (\cref{prop:estratta-limite}). \emph{Scelta logica:} la stima
della differenza di due radici si ottiene ancora razionalizzando; il caso $c=0$ va trattato
a parte perché la stima divide per $\sqrt c$.
\end{soluzione}
\end{esercizio}

\begin{esercizio}\label{es:quesiti-b}
Riprendiamo l'osservazione che segue il \cref{cor:radice-n-nb}. Posto
$b_{n}=n^{\beta/n}-1$ con $\beta>0$, la disuguaglianza di Bernoulli dà
$n^{\beta}=(1+b_{n})^{n}\ge1+nb_{n}$. Per quali valori di $\beta$ se ne deduce che
$b_{n}\to0$? Perché la scelta $\beta=1$ non funziona?
\begin{soluzione}
Si ottiene $0\le b_{n}\le\frac{n^{\beta}-1}{n}<n^{\beta-1}$, e l'ultimo membro tende a
zero se e solo se $\beta<1$ (\cref{lem:potenze-log}(i)). Per $\beta=1$ la stima diventa
$b_{n}\le\frac{n-1}{n}$, che tende a $1$ e non dice nulla. Quindi ogni $\beta\in(0,1)$ va
bene, e $\beta=\frac12$ è semplicemente la scelta più comoda. Con $\beta=\frac1m$
($m\ge2$ naturale) il caso di un esponente intero $k$ segue poi da
$n^{k/n}=\bigl(n^{1/(mn)}\bigr)^{mk}$, e il caso di $b$ reale con la parte intera, come
nella dimostrazione del \cref{cor:radice-n-nb}.
\emph{Scelta logica:} quando una stima «quasi funziona», conviene chiedersi quale parametro
la rende troppo debole, e indebolire l'obiettivo (qui: dimostrare il risultato per un
esponente più piccolo) per poi ricostruire il caso generale.
\end{soluzione}
\end{esercizio}

\subsection*{Successioni monotòne, numero $e$, criterio del rapporto}

\begin{esercizio}\label{es:n10-su-1-1n}
Dimostrare che $\dfrac{n^{10}}{1{,}1^{n}}\to0$, e calcolare con Python il massimo valore
assunto dalla successione.
\begin{soluzione}
Per il \cref{teo:gerarchia} con $b=10$ e $a=1{,}1$. Il rapporto tra termini consecutivi
vale $\bigl(1+\frac1n\bigr)^{10}/1{,}1$, che è maggiore di $1$ finché
$\bigl(1+\frac1n\bigr)^{10}>1{,}1$, cioè per $n$ fino a circa $104$:
\begin{python}
a = lambda n: n**10 / 1.1**n
n_max = max(range(1, 1000), key=a)
print("massimo per n =", n_max, f"  a(n) = {a(n_max):.3e}")
print(f"a(500) = {a(500):.3e},  a(1000) = {a(1000):.3e}")
\end{python}
\begin{risultato}
massimo per n = 105   a(n) = 7.339e+15
a(500) = 1.965e+06,  a(1000) = 4.049e-12
\end{risultato}
La successione cresce fino a $n=105$, dove supera $7\cdot10^{15}$, e solo dopo comincia a
decrescere verso zero: a $n=500$ vale ancora circa $2\cdot10^{6}$, a $n=1000$ circa
$4\cdot10^{-12}$. \emph{Scelta logica:} un
limite descrive il comportamento \emph{definitivo}, e non dice nulla su quanto presto esso
si manifesti; chi usa i limiti per stimare grandezze concrete deve tenerlo presente.
\end{soluzione}
\end{esercizio}

\begin{esercizio}
Dimostrare che $\dfrac{2^{n}\,n!}{n^{n}}\to0$, mentre $\dfrac{3^{n}\,n!}{n^{n}}\to+\infty$.
\begin{soluzione}
Il rapporto tra termini consecutivi della prima vale
$2\cdot\frac{n^{n}}{(n+1)^{n}}=\frac{2}{(1+1/n)^{n}}\to\frac2e<1$, e il criterio del
rapporto dà la tesi. Per la seconda il rapporto tende a $\frac3e>1$, e per l'osservazione
che segue il \cref{teo:criterio-rapporto} la successione diverge. \emph{Scelta logica:} il
valore $e$ compare in modo naturale in tutti i confronti con $n^{n}$; il confine tra
convergenza e divergenza è proprio $a=e$ (per $a=e$ il criterio non decide).
\end{soluzione}
\end{esercizio}

\begin{esercizio}[Il problema di Basilea]
Sia $s_{n}=\displaystyle\sum_{k=1}^{n}\frac1{k^{2}}$. Dimostrare che $s_{n}$ converge.
\begin{soluzione}
La successione è strettamente crescente. Per $k\ge2$ si ha
$\frac1{k^{2}}<\frac1{k(k-1)}$, e quindi, per la somma telescopica
(\cref{prop:telescopica}),
\[
  s_{n}\le1+\sum_{k=2}^{n}\frac1{(k-1)k}=1+\frac{n-1}{n}<2
\]
(con l'uguaglianza solo per $n=1$, quando la somma a destra è vuota).
Per il teorema sulle successioni monotòne $s_{n}$ converge a un limite $s\le2$. Il valore
del limite, $s=\pi^{2}/6\approx1{,}6449$, fu trovato da Eulero nel 1734, risolvendo un
problema posto da Pietro Mengoli nel 1650 e rimasto aperto per oltre ottant'anni.
\emph{Scelta logica:} non sappiamo calcolare $s_{n}$ in forma chiusa, ma per la
convergenza basta una maggiorazione con una somma che sappiamo calcolare.
\end{soluzione}
\end{esercizio}

\begin{esercizio}[La serie armonica]
Sia $H_{n}=1+\frac12+\dots+\frac1n$. Dimostrare che $H_{2n}-H_{n}\ge\frac12$ per ogni $n$,
e dedurne che $H_{n}$ non è di Cauchy e che $H_{n}\to+\infty$.
\begin{soluzione}
$H_{2n}-H_{n}=\frac1{n+1}+\dots+\frac1{2n}$ è una somma di $n$ termini, ciascuno
$\ge\frac1{2n}$; quindi $H_{2n}-H_{n}\ge n\cdot\frac1{2n}=\frac12$. Con $\epsilon=\frac12$,
per ogni $\nu$ gli indici $h=n$, $k=2n$ con $n>\nu$ violano la condizione di Cauchy: la
successione non è di Cauchy, e quindi non converge (\cref{teo:criterio-cauchy}). Essendo
crescente, per il \cref{teo:monotone} diverge a $+\infty$. La divergenza è però
lentissima: $H_{n}$ supera $10$ solo per $n>12\,000$ circa. Ne riparleremo nel
\cref{sec:serie-armonica}. \emph{Scelta logica:} i termini $\frac1n$ tendono a zero,
eppure la somma diverge; il confronto con l'esercizio precedente mostra che conta la
\emph{velocità} con cui i termini tendono a zero.
\end{soluzione}
\end{esercizio}

\begin{esercizio}[Con Python]
Generalizzare il metodo di Erone: per $a>0$ e $x_{1}>0$ qualsiasi, la successione
$x_{n+1}=\frac12\bigl(x_{n}+\frac a{x_{n}}\bigr)$ converge a $\sqrt a$. Dimostrarlo
adattando l'\cref{es:erone} (attenzione: se $x_{1}<\sqrt a$, la successione non è
monotòna fin dal primo termine), e verificare numericamente quante iterazioni servono per
calcolare $\sqrt{10}$ con $15$ cifre partendo da $x_{1}=1$.
\begin{soluzione}
Come nell'esempio, $x_{n+1}-\sqrt a=\frac{(x_{n}-\sqrt a)^{2}}{2x_{n}}\ge0$, quindi
$x_{n}\ge\sqrt a$ per $n\ge2$ (anche se $x_{1}<\sqrt a$), e da $n=2$ in poi la successione
è decrescente: è \emph{definitivamente} monotòna, e la \cref{oss:traslata} permette di
concludere come nel testo.
\begin{python}
import math
a, x = 10.0, 1.0
r = math.sqrt(a)
for n in range(1, 9):
    print(f"x_{n} = {x:.15f}   errore = {abs(x - r):.1e}")
    x = (x + a / x) / 2
\end{python}
\begin{risultato}
x_1 = 1.000000000000000   errore = 2.2e+00
x_2 = 5.500000000000000   errore = 2.3e+00
x_3 = 3.659090909090909   errore = 5.0e-01
x_4 = 3.196005081874647   errore = 3.4e-02
x_5 = 3.162455622803890   errore = 1.8e-04
x_6 = 3.162277665175675   errore = 5.0e-09
x_7 = 3.162277660168379   errore = 4.4e-16
x_8 = 3.162277660168379   errore = 4.4e-16
\end{risultato}
Il primo passo peggiora l'approssimazione (da $x_{1}=1$, per difetto, si salta a
$x_{2}=5{,}5$, per eccesso); poi scatta la convergenza quadratica, e l'errore passa da
circa $0{,}5$ a $3\cdot10^{-2}$, $2\cdot10^{-4}$, $5\cdot10^{-9}$: alla settima iterazione si
raggiunge la precisione della macchina. \emph{Scelta logica:} l'identità
$x_{n+1}-\sqrt a=\frac{(x_{n}-\sqrt a)^{2}}{2x_{n}}$ fa tutto il lavoro: dice che dal
secondo termine in poi si sta sopra $\sqrt a$, e quindi che la successione è
definitivamente decrescente e limitata, cioè proprio le ipotesi del \cref{teo:monotone}.
\end{soluzione}
\end{esercizio}

%% ----------------------------------------------------------------------------
%%  §3.12  Punti fondamentali del capitolo
%% ----------------------------------------------------------------------------
\section{Punti fondamentali del capitolo}
\label{sec:sintesi-cap03}

L'idea centrale del capitolo è che il comportamento «all'infinito» di una successione si
può descrivere con un numero finito di parole: per ogni tolleranza $\epsilon$ esiste un
indice oltre il quale i termini la rispettano. Da questa definizione discendono le regole
di calcolo, i teoremi di confronto e, grazie all'assioma di completezza, i teoremi di
esistenza che permettono di affermare che un limite c'è senza conoscerlo in anticipo.

\begin{puntichiave}
\item \textbf{Definizione di limite.} $a_{n}\to a\in\R$ se per ogni $\epsilon>0$ esiste
$\nu$ tale che $\abs{a_{n}-a}<\epsilon$ per ogni $n>\nu$ (\cref{def:limite}); $a_{n}\to+\infty$
se per ogni $M>0$ esiste $\nu$ tale che $a_{n}>M$ per ogni $n>\nu$
(\cref{def:limite-infinito}). L'indice $\nu$ dipende da $\epsilon$: scambiare i
quantificatori cambia il significato. Il limite, se esiste, è unico
(\cref{teo:unicita-limite}).
\item \textbf{Convergenza e limitatezza.} Ogni successione convergente è limitata
(\cref{teo:convergente-limitata}), ma non viceversa: $(-1)^{n}$ è limitata e non ha
limite (\cref{prop:alterna-non-converge}).
\item \textbf{Operazioni con i limiti.} Se $a_{n}\to a$ e $b_{n}\to b$, con $a,b\in\R$,
allora $a_{n}\pm b_{n}\to a\pm b$, $a_{n}b_{n}\to ab$ e, se $b_{n}\neq0$ e $b\neq0$,
$a_{n}/b_{n}\to a/b$ (\cref{teo:operazioni-limiti}). Con limiti infiniti valgono regole
analoghe (\cref{teo:operazioni-infinite}).
\item \textbf{Forme indeterminate.} Nelle forme $\infty-\infty$, $0\cdot\infty$,
$\frac\infty\infty$, $\frac00$ e $1^{\infty}$, $\infty^{0}$, $0^{0}$ i limiti dei
singoli termini non bastano: occorre prima trasformare l'espressione (raccogliere il
termine dominante, razionalizzare), come negli \cref{es:quoziente-polinomi} e
\cref{es:razionalizzare}. Attenzione: se $b_{n}\to0$, soltanto $\abs{1/b_{n}}$ diverge;
$1/b_{n}\to+\infty$ solo se $b_{n}>0$ definitivamente.
\item \textbf{Permanenza del segno e confronto.} Se $a_{n}\to a>0$, allora $a_{n}>0$
definitivamente (\cref{teo:permanenza-segno}). Le disuguaglianze larghe passano al
limite (\cref{cor:confronto-limiti}); quelle strette si indeboliscono: $\frac1n>0$, ma il
limite è $0$.
\item \textbf{Teorema dei carabinieri.} Se $a_{n}\le c_{n}\le b_{n}$ definitivamente e
$a_{n}$, $b_{n}$ tendono allo stesso limite $a$, allora anche $c_{n}$ ha limite, uguale
ad $a$ (\cref{teo:carabinieri}). Ne segue che il prodotto di una successione limitata per
una infinitesima è infinitesimo (\cref{teo:limitata-infinitesima}), ad esempio
$\frac{\sen n}{n}\to0$.
\item \textbf{Limiti notevoli.} $a^{n}$ tende a $+\infty$, $1$ o $0$ secondo che $a>1$,
$a=1$ o $\abs a<1$, e non ha limite per $a\le-1$ (\cref{teo:lim-geometrica});
$\sqrt[n]a\to1$ per $a>0$ e $\sqrt[n]n\to1$ (\cref{prop:radice-n-a},
\cref{prop:radice-n-n}). Lo strumento comune sono le stime dal basso di $(1+x)^{n}$ date
da Bernoulli e dal binomio.
\item \textbf{Il limite $\frac{\sen a_{n}}{a_{n}}$.} Se $a_{n}\to0$ con $a_{n}\neq0$,
allora $\frac{\sen a_{n}}{a_{n}}\to1$ (\cref{teo:sen-x-su-x}); ne segue che l'area del
cerchio di raggio $1$ è $\pi$ (\cref{teo:area-cerchio}). Non va confuso con
$\frac{\sen n}{n}\to0$, dove l'argomento diverge.
\item \textbf{Successioni monotòne.} Ogni successione monotòna ha limite: il $\sup$ se è
crescente, l'$\inf$ se è decrescente; se è anche limitata, converge (\cref{teo:monotone}).
Per le successioni definite per ricorrenza si dimostra prima che il limite esiste, poi lo
si calcola passando al limite nella ricorrenza (\cref{es:erone}).
\item \textbf{Il numero $e$.} $(1+\frac1n)^{n}$ è strettamente crescente,
$(1+\frac1n)^{n+1}$ strettamente decrescente, ed entrambe tendono ad $e$, con
$(1+\frac1n)^{n}<e<(1+\frac1n)^{n+1}$ per ogni $n$ (\cref{teo:numero-e}) e $2<e<3$
(\cref{cor:e-tra-2-e-3}). Più in generale $(1+\frac1{a_{n}})^{a_{n}}\to e$ se
$a_{n}\to\pm\infty$ (\cref{teo:e-generalizzato}).
\item \textbf{Criterio del rapporto e gerarchia degli infiniti.} Se $a_{n}>0$ e
$\frac{a_{n+1}}{a_{n}}\to b<1$, allora $a_{n}\to0$ (\cref{teo:criterio-rapporto}); per
$b>1$ la successione diverge, per $b=1$ il criterio tace. Per $b>0$ e $a>1$ vale
$\log n\ll n^{b}\ll a^{n}\ll n!\ll n^{n}$ (\cref{teo:gerarchia}), anche se il sorpasso
può avvenire molto tardi (\cref{es:n10-su-1-1n}).
\item \textbf{Estratte e teorema di Bolzano--Weierstrass.} Se $a_{n}$ ha limite, ogni
estratta ha lo stesso limite (\cref{prop:estratta-limite}); due estratte con limiti
diversi escludono l'esistenza del limite. Ogni successione limitata ha un'estratta
convergente (\cref{teo:bolzano-weierstrass}).
\item \textbf{Criterio di Cauchy.} Una successione reale converge se e solo se è di
Cauchy (\cref{teo:criterio-cauchy}): la convergenza si verifica senza conoscere il
limite. In $\Q$ il criterio è falso, perché poggia sulla completezza di $\R$.
\item \textbf{Limiti e calcolo numerico.} La definizione di limite è la tolleranza sul
residuo di un solutore iterativo; il metodo di Erone, caso particolare del metodo di
Newton, converge quadraticamente; il limite $\frac{\sen\alpha}{\alpha}\to1$ giustifica le
approssimazioni $W\sen\gamma\approx W\gamma$ e $C_{L}=2\pi\alpha$ per piccoli angoli in
radianti. Ma il calcolatore va usato con cautela: in doppia precisione
$(1+10^{-16})^{10^{16}}$ risulta $1$.
\end{puntichiave}

\paragraph{Il filo del discorso.} Il capitolo ha studiato il comportamento di una
quantità che dipende da un indice naturale. Restano aperte due questioni: che cosa
accade a $f(x)$ quando la variabile \emph{reale} $x$ si avvicina a un punto, e la
«continuità» di esponenziali e potenze, che qui abbiamo dovuto aggirare. Il
\cref{cap:funzioni-continue} risponde a entrambe definendo il limite di una funzione
proprio mediante le successioni, così da ereditare tutti i risultati di questo capitolo;
il metodo di bisezione e il teorema di Bolzano--Weierstrass diventeranno il teorema degli
zeri e il teorema di Weierstrass.

